% Leçon 22 — Grammaire : phrase complexes — Chinois Tle A — Cours signé OnBuch+
% Programme officiel MINESEC. Compiler avec : tectonic ZH22-grammaire-phrase-complexes.tex
\newcommand{\DOCMATIERE}{Chinois}
\newcommand{\DOCNIVEAU}{Tle A}
\newcommand{\DOCMODULE}{Module 3 : Citoyenneté et ouverture au monde}
\newcommand{\DOCLECON}{ZH22}
\newcommand{\DOCTITRE}{Leçon 22 — Grammaire : phrase complexes}
% OnBuch+ — préambule « cours signés » ENRICHI (compilé avec tectonic / XeLaTeX)
% Système visuel « Pop » d'OnBuch+ : fond crème (jamais blanc), bordures encre
% épaisses, ombres « dures » décalées, titres Archivo Black en capitales.
% Version MATHÉMATIQUES : graphes (pgfplots/tikz), boîtes pédagogiques
% (définition, propriété, méthode, attention, le savais-tu, exercice résolu,
% exercice, TP, bilan), page de garde et signature OnBuch+ ancrée en bas.
%
% Macros attendues dans main.tex AVANT \input{preamble} :
%   \DOCMATIERE  \DOCNIVEAU  \DOCMODULE  \DOCLECON (numéro)  \DOCTITRE
\documentclass[11pt]{article}

\usepackage{fontspec}
\usepackage[french]{babel} % français = langue principale ; le chinois via \zh{..} / textebox (xeCJK)
\usepackage{xeCJK}
\usepackage{pifont} % \ding{51} \ding{55} utilisés par les modèles
\usepackage[a4paper,margin=2cm,top=3.4cm,bottom=2.8cm,headheight=1.4cm,footskip=1.3cm]{geometry}
\usepackage{amsmath,amssymb,amsfonts}
\usepackage{mathtools} % \xmapsto, \coloneqq... souvent utilisés par les modèles
\usepackage{cancel} % simplifications barrées (bilans de forces)
% \abs{z} (module, valeur absolue) et \norm{u} (norme), utilisés par les modèles
\providecommand{\abs}{}\let\abs\relax\DeclarePairedDelimiter{\abs}{\lvert}{\rvert}
\providecommand{\norm}{}\let\norm\relax\DeclarePairedDelimiter{\norm}{\lVert}{\rVert}
\usepackage[table]{xcolor} % [table] : fournit aussi \rowcolors (alternance de lignes)
\usepackage{pagecolor}
\usepackage{tikz}
\usetikzlibrary{arrows.meta,positioning,calc,shapes.geometric,shapes.misc,shapes.arrows,decorations.pathmorphing,decorations.pathreplacing,decorations.markings,patterns,shadows,mindmap,trees,fit,backgrounds,matrix,intersections,angles,quotes}
\usepackage{pgfplots}
\usepackage[version=4]{mhchem} % équations nucléaires \ce{^{235}_{92}U}
\usepackage[european,siunitx]{circuitikz} % schémas électriques
\pgfplotsset{compat=1.18}
\usepgfplotslibrary{fillbetween,groupplots} % aires entre courbes (calcul intégral)
\usepackage{enumitem}
\usepackage{fancyhdr}
% [most] charge toutes les bibliothèques tcolorbox (listings, minted, raster,
% documentation...) et gonfle inutilement la mémoire TeX sur un document long
% (~300 boîtes) : on ne charge que ce qui est réellement utilisé.
\usepackage[skins,breakable]{tcolorbox}
\usepackage{graphicx}
\usepackage{colortbl}
\usepackage{booktabs}
\usepackage{tabularx}
\usepackage{diagbox} % case d'en-tête diagonale des tableaux à double entrée
% Colonnes extensibles centrée / gauche / droite (C, L, R), souvent utilisées par les modèles
\newcolumntype{C}{>{\centering\arraybackslash}X}
\newcolumntype{L}{>{\raggedright\arraybackslash}X}
\newcolumntype{R}{>{\raggedleft\arraybackslash}X}
\usepackage{array}
\usepackage{multicol}
\usepackage{multirow}
\usepackage{siunitx}
% parse-numbers=false : accepte aussi \SI{2,5 \times 10^{-3}}{...}
\sisetup{locale=FR,output-decimal-marker={,},group-minimum-digits=4,parse-numbers=false}
\DeclareSIUnit\ohm{\ensuremath{\symup{\Omega}}} % U+2126 absent de la police
\DeclareSIUnit\division{div} % réglages d'oscilloscope (ms/div, V/div)
\DeclareSIUnit\tr{tr} % tours (vitesse de rotation, ex. \SI{1500}{\tr\per\minute})
\DeclareSIUnit\molar{M} % concentration molaire (ex. \SI{3}{\molar}, \SI{1}{\micro\molar})
\DeclareSIUnit\an{an} % année (le modèle confond parfois avec \year, un compteur TeX) — ex. \SI{5}{\kilo\meter\per\an}
\DeclareSIUnit\FCFA{FCFA} % franc CFA (ex. \SI{500}{\FCFA\per\kilo\gram})
\DeclareSIUnit\degreeBrix{\textdegree Brix} % teneur en sucre d'un sirop/jus (réfractométrie)
\DeclareSIUnit\month{mois} % le modèle utilise parfois \month, un compteur TeX, comme unité de durée
\usepackage{qrcode}
% \degree : commande fréquemment utilisée par les modèles pour un angle ou une
% température en dehors de \SI{}{} (non fournie par les paquets chargés).
\providecommand{\degree}{^{\circ}}
% \qed (fin de démonstration), utilisé par les modèles sans amsthm
\providecommand{\qed}{\hfill$\square$}
% \barre : double barre des valeurs interdites dans un tableau de variations (tkz-tab)
\providecommand{\barre}{\ensuremath{\|}}
% environnement proof (amsthm non chargé) utilisé par les modèles
\ifdefined\proof\else\newenvironment{proof}[1][Démonstration]{\par\noindent\textit{#1.}\ }{\qed\par}\fi
\usepackage{lastpage}
\usepackage{hyperref}
\hypersetup{hidelinks}

% ============================================================
% Palette « Pop » OnBuch+ — mêmes hex que la classe P (plus_ui.dart)
% ============================================================
\definecolor{poporange}{HTML}{F59321}
\definecolor{popdark}{HTML}{E07A0C}
\definecolor{popcream}{HTML}{FFF6E8}
\definecolor{popink}{HTML}{1C1714}
\definecolor{poppurple}{HTML}{7A5AE0}
\definecolor{popgreen}{HTML}{2E9E6A}
\definecolor{poppink}{HTML}{E0457B}
\definecolor{popblue}{HTML}{2F7DE1}
\definecolor{popgold}{HTML}{C98A12}
\definecolor{popmuted}{HTML}{8A7F76}
\definecolor{popline}{HTML}{EADFD2}
\definecolor{popgray}{HTML}{8A7F76}
\colorlet{popgrey}{popgray}
% Alias tolérants (le modèle nomme parfois les couleurs différemment)
\colorlet{popwhite}{white}
\colorlet{popblack}{popink}
\colorlet{popred}{poppink}
\colorlet{popyellow}{popgold}
% Teintes claires (fonds de tableaux/encadrés) — le modèle nomme parfois
% les couleurs avec un suffixe L (ex. popblueL) sans les avoir définies.
\colorlet{poporangeL}{poporange!20}
\colorlet{popdarkL}{popdark!20}
\colorlet{poppurpleL}{poppurple!20}
\colorlet{popgreenL}{popgreen!20}
\colorlet{poppinkL}{poppink!20}
\colorlet{popblueL}{popblue!20}
\colorlet{popgoldL}{popgold!20}
\colorlet{popredL}{popred!20}
\colorlet{popgrayL}{popgray!20}
% Teintes claires pour fonds de tableaux / zones
\colorlet{popblueL}{popblue!10!white}
\colorlet{poporangeL}{poporange!14!white}
\colorlet{popgreenL}{popgreen!12!white}
\colorlet{poppurpleL}{poppurple!12!white}
\colorlet{poppinkL}{poppink!12!white}
\colorlet{popgoldL}{popgold!14!white}

\pagecolor{popcream}

% ============================================================
% Typographie
% ============================================================
\defaultfontfeatures{Ligatures=TeX}
\setmainfont{PlusJakartaSans-Medium.ttf}[Path=../../../fonts/,BoldFont=PlusJakartaSans-Bold.ttf]
% Symboles, API et emojis absents de la police principale : repli sur DejaVu Sans (caractères actifs)
\newfontface\symfont{DejaVuSans.ttf}[Path=../../../fonts/]
\makeatletter
\newcommand\symactive[1]{\catcode#1=13 \begingroup\lccode`\~=#1 \lowercase{\endgroup\def~}{{\symfont\char#1}}}
\newcommand\symblank[1]{\catcode#1=13 \begingroup\lccode`\~=#1 \lowercase{\endgroup\def~}{}}
\newcommand\symhyph[1]{\catcode#1=13 \begingroup\lccode`\~=#1 \lowercase{\endgroup\def~}{\mbox{-}}}
\newcount\symc
\newcommand\symrange[2]{\symc=#1 \@whilenum\symc<#2 \do{\expandafter\symactive\expandafter{\the\symc}\advance\symc by 1 }}
\newcommand\symblankrange[2]{\symc=#1 \@whilenum\symc<#2 \do{\expandafter\symblank\expandafter{\the\symc}\advance\symc by 1 }}
\makeatother
\symrange{"0250}{"02FF}\symactive{"2713}\symactive{"2714}\symactive{"2717}\symactive{"2718}\symactive{"2610}\symactive{"2611}\symactive{"2612}
\symrange{"2600}{"27BF}
\catcode"2705=13 \begingroup\lccode`\~="2705 \lowercase{\endgroup\def~}{{\symfont\char"2713}}\symblank{"26BD}\symblank{"26A1}
\symrange{"2460}{"257F}\symactive{"223C}\symactive{"2248}\symactive{"2260}
\symactive{"01F8}\symactive{"01F9}\symactive{"1E3E}\symactive{"1E3F}
\symhyph{"2011}\symhyph{"2E1A}\symblankrange{"1F300}{"1FAFF}\symblank{"FE0F}
% Caractères chinois : WenQuanYi Zen Hei (sous-ensemble GB2312 + pinyin), gras/italique simulés
\setCJKmainfont{WenQuanYiZenHei-subset.ttf}[Path=../../../fonts/,AutoFakeBold=3,AutoFakeSlant=0.2]
\xeCJKsetup{CJKspace=false}
% Pinyin : voyelles à caron (ǎ ǐ ǒ ǔ ǖ ǘ ǚ ǜ) absentes de la police principale -> caractères actifs
\newfontface\pyfont{WenQuanYiZenHei-subset.ttf}[Path=../../../fonts/]
\newcommand\pyactive[1]{\catcode#1=13 \begingroup\lccode`\~=#1 \lowercase{\endgroup\def~}{{\pyfont\char#1}}}
\pyactive{"01CD}\pyactive{"01CE}\pyactive{"01CF}\pyactive{"01D0}\pyactive{"01D1}\pyactive{"01D2}\pyactive{"01D3}\pyactive{"01D4}\pyactive{"01D5}\pyactive{"01D6}\pyactive{"01D7}\pyactive{"01D8}\pyactive{"01D9}\pyactive{"01DA}\pyactive{"01DB}\pyactive{"01DC}
% Maths en sans-serif (Fira Math) avec chiffres et lettres droites de Plus
% Jakarta, pour que les formules se fondent dans le texte.
\usepackage[math-style=ISO,bold-style=ISO]{unicode-math}
\setmathfont{FiraMath-Regular.otf}[Path=../../../fonts/]
\setmathfont{PlusJakartaSans-Medium.ttf}[Path=../../../fonts/,range={up/num,up/{Latin,latin},\mathup}]
\usepackage{icomma} % virgule décimale sans espace en mode math
% Caractères mathématiques Unicode absents des polices de texte (Plus Jakarta,
% Archivo Black) : rendus par leur équivalent mathématique, en texte comme en
% formule — sinon ils s'affichent en carré vide (ex. « Arithmétique dans ℤ »).
\usepackage{newunicodechar}
\newunicodechar{ℝ}{\ensuremath{\mathbb{R}}}
\newunicodechar{ℤ}{\ensuremath{\mathbb{Z}}}
\newunicodechar{ℂ}{\ensuremath{\mathbb{C}}}
\newunicodechar{ℕ}{\ensuremath{\mathbb{N}}}
\newunicodechar{ℚ}{\ensuremath{\mathbb{Q}}}
\newunicodechar{𝔻}{\ensuremath{\mathbb{D}}}
\newunicodechar{α}{\ensuremath{\alpha}}
\newunicodechar{β}{\ensuremath{\beta}}
\newunicodechar{γ}{\ensuremath{\gamma}}
\newunicodechar{δ}{\ensuremath{\delta}}
\newunicodechar{ε}{\ensuremath{\varepsilon}}
\newunicodechar{θ}{\ensuremath{\theta}}
\newunicodechar{λ}{\ensuremath{\lambda}}
\newunicodechar{μ}{\ensuremath{\mu}}
\newunicodechar{σ}{\ensuremath{\sigma}}
\newunicodechar{τ}{\ensuremath{\tau}}
\newunicodechar{φ}{\ensuremath{\varphi}}
\newunicodechar{ψ}{\ensuremath{\psi}}
\newunicodechar{ω}{\ensuremath{\omega}}
\newunicodechar{Σ}{\ensuremath{\Sigma}}
% majuscules produites par \MakeUppercase dans les titres (α-aminés -> Α)
\newunicodechar{Α}{\ensuremath{\alpha}}
\newunicodechar{Β}{\ensuremath{\beta}}
\newunicodechar{Γ}{\ensuremath{\Gamma}}
\newunicodechar{Δ}{\ensuremath{\Delta}}
\newunicodechar{Ε}{\ensuremath{\varepsilon}}
\newunicodechar{Θ}{\ensuremath{\Theta}}
\newunicodechar{Λ}{\ensuremath{\Lambda}}
\newunicodechar{Μ}{\ensuremath{\mu}}
\newunicodechar{Τ}{\ensuremath{\tau}}
\newunicodechar{Φ}{\ensuremath{\Phi}}
\newunicodechar{Ψ}{\ensuremath{\Psi}}
\newunicodechar{Ω}{\ensuremath{\Omega}}
\newunicodechar{↦}{\ensuremath{\mapsto}}
\newunicodechar{∈}{\ensuremath{\in}}
\newunicodechar{∉}{\ensuremath{\notin}}
\newunicodechar{∧}{\ensuremath{\wedge}}
\newunicodechar{⇔}{\ensuremath{\Leftrightarrow}}
\newunicodechar{⟺}{\ensuremath{\Longleftrightarrow}}
\newunicodechar{⇒}{\ensuremath{\Rightarrow}}
\newunicodechar{⟹}{\ensuremath{\Longrightarrow}}
\newunicodechar{∘}{\ensuremath{\circ}}
\newunicodechar{∪}{\ensuremath{\cup}}
\newunicodechar{∩}{\ensuremath{\cap}}
\newunicodechar{≡}{\ensuremath{\equiv}}
\newunicodechar{⊂}{\ensuremath{\subset}}
\newunicodechar{⊥}{\ensuremath{\perp}}
\newunicodechar{∃}{\ensuremath{\exists}}
\newunicodechar{∀}{\ensuremath{\forall}}
\newunicodechar{∛}{\ensuremath{\sqrt[3]{\,}}}
\newunicodechar{∜}{\ensuremath{\sqrt[4]{\,}}}
\newunicodechar{⃗}{\ensuremath{{}^{\to}}}
\newunicodechar{ⁿ}{\ifmmode^{n}\else\textsuperscript{\ensuremath{n}}\fi}
\newunicodechar{ˣ}{\ifmmode^{x}\else\textsuperscript{\ensuremath{x}}\fi}
\newunicodechar{ᵇ}{\ifmmode^{b}\else\textsuperscript{\ensuremath{b}}\fi}
\newunicodechar{ʳ}{\ifmmode^{r}\else\textsuperscript{\ensuremath{r}}\fi}
\newunicodechar{ᵘ}{\ifmmode^{u}\else\textsuperscript{\ensuremath{u}}\fi}
\newunicodechar{ʸ}{\ifmmode^{y}\else\textsuperscript{\ensuremath{y}}\fi}
\newunicodechar{ᵃ}{\ifmmode^{a}\else\textsuperscript{\ensuremath{a}}\fi}
\newunicodechar{ᵉ}{\ifmmode^{e}\else\textsuperscript{\ensuremath{e}}\fi}
\newunicodechar{ᵏ}{\ifmmode^{k}\else\textsuperscript{\ensuremath{k}}\fi}
\newunicodechar{ᵖ}{\ifmmode^{p}\else\textsuperscript{\ensuremath{p}}\fi}
\newunicodechar{ᴺ}{\ifmmode^{N}\else\textsuperscript{\ensuremath{N}}\fi}
\newunicodechar{ᶜ}{\ifmmode^{c}\else\textsuperscript{\ensuremath{c}}\fi}
\newunicodechar{ʰ}{\ifmmode^{h}\else\textsuperscript{\ensuremath{h}}\fi}
\newunicodechar{ᵠ}{\ifmmode^{q}\else\textsuperscript{\ensuremath{q}}\fi}
\newunicodechar{ₙ}{\ifmmode_{n}\else\textsubscript{\ensuremath{n}}\fi}
\newunicodechar{ₐ}{\ifmmode_{a}\else\textsubscript{\ensuremath{a}}\fi}
\newunicodechar{ᵢ}{\ifmmode_{i}\else\textsubscript{\ensuremath{i}}\fi}
\newunicodechar{ᵣ}{\ifmmode_{r}\else\textsubscript{\ensuremath{r}}\fi}
\newunicodechar{ₖ}{\ifmmode_{k}\else\textsubscript{\ensuremath{k}}\fi}
\newunicodechar{ᵦ}{\ifmmode_{\beta}\else\textsubscript{\ensuremath{\beta}}\fi}
\newunicodechar{ₓ}{\ifmmode_{x}\else\textsubscript{\ensuremath{x}}\fi}
\newfontfamily\popdisplay{ArchivoBlack-Regular.ttf}[Path=../../../fonts/]
\newfontfamily\poplogo{SpaceGrotesk-Bold.ttf}[Path=../../../fonts/]
\newfontfamily\popsemi{PlusJakartaSans-SemiBold.ttf}[Path=../../../fonts/]
\newfontfamily\popxbold{PlusJakartaSans-ExtraBold.ttf}[Path=../../../fonts/]
\newfontfamily\popxboldls{PlusJakartaSans-ExtraBold.ttf}[Path=../../../fonts/,LetterSpace=3.0]

\setlist[enumerate]{leftmargin=1.7em,itemsep=5pt,topsep=4pt}
\setlist[itemize]{leftmargin=1.7em,itemsep=4pt,topsep=4pt,label={\color{poporange}\textbullet}}
\setlength{\parindent}{0pt}
\setlength{\parskip}{5pt}
\renewcommand{\arraystretch}{1.25}

% pgfplots : style Pop par défaut
\pgfplotsset{
  every axis/.append style={axis line style={popink,line width=1pt},tick style={popink},
    grid style={popline,line width=0.5pt},label style={font=\small},tick label style={font=\footnotesize}},
  popcurve/.style={poporange,line width=1.8pt,no markers,smooth},
}
% Même style disponible dans un \draw TikZ simple (hors pgfplots)
\tikzset{popcurve/.style={poporange,line width=1.8pt}}
\tikzset{popgrid/.style={popline,very thin}}
\colorlet{popcurve}{poporange} % le nom du style est parfois utilisé comme couleur

% ============================================================
% Pill / squiggle
% ============================================================
\newcommand{\poppill}[3]{%
  \tikz[baseline=-0.55ex]\node[draw=popink,line width=1.2pt,rounded corners=2.6pt,%
    fill=#2,inner xsep=6.5pt,inner ysep=3pt]{%
    \popxboldls\fontsize{7}{7}\selectfont\color{#3}\MakeUppercase{#1}};%
}
\newcommand{\popsquiggle}[1]{%
  \tikz[baseline=-0.4ex]{
    \draw[#1,line width=2.6pt,line cap=round]
      plot [smooth] coordinates {(0,0.09) (0.34,0.01) (0.68,0.21) (1.02,0.01) (1.36,0.10)};
  }%
}
% Mot-clé surligné (marqueur orange sous le texte)
\newcommand{\cle}[1]{{\popxbold\color{popdark}#1}}

% ============================================================
% En-tête / pied
% ============================================================
\pagestyle{fancy}
\fancyhf{}
\renewcommand{\headrulewidth}{0pt}
\renewcommand{\footrulewidth}{0pt}
\lhead{\raisebox{-0.15ex}{\poplogo\fontsize{13}{13}\selectfont\color{popink}OnBuch\color{poporange}+}}
\rhead{\poppill{\DOCMATIERE\ \textbullet\ \DOCNIVEAU\ \textbullet\ Leçon \DOCLECON}{white}{popink}}
\lfoot{\popxbold\fontsize{7.6}{7.6}\selectfont\color{popmuted}\MakeUppercase{Cours signé OnBuch+}\ \textbullet\ {\popsemi\DOCTITRE}}
\rfoot{\tikz[baseline=-0.6ex]\node[rounded corners=8.5pt,fill=popink,minimum height=17pt,minimum width=17pt,inner xsep=5pt,inner sep=0]{\popxbold\fontsize{8}{8}\selectfont\color{popcream}\thepage\,/\,\pageref*{LastPage}};}

% ============================================================
% Titres
% ============================================================
\newcommand{\chaptitre}[2]{%
  \begin{tcolorbox}[enhanced,colback=poporange,colframe=popink,boxrule=2.6pt,arc=11pt,
      drop shadow=popink,left=15pt,right=15pt,top=12pt,bottom=11pt,before skip=3pt,after skip=18pt]
    \poppill{#2}{popink}{popcream}\par\vspace{9pt}
    {\raggedright\hyphenpenalty=10000\exhyphenpenalty=10000\popdisplay\fontsize{22}{24}\selectfont\color{popcream}\MakeUppercase{#1}\par}
  \end{tcolorbox}
}
% Section numérotée : pastille orange avec le numéro + titre Archivo
\newcounter{popsec}
\newcommand{\coursec}[1]{%
  \stepcounter{popsec}\setcounter{popsub}{0}%
  \par\vspace{16pt}\noindent
  \begin{minipage}{\linewidth}
  \tikz[baseline=-0.7ex]\node[draw=popink,line width=1.6pt,fill=poporange,rounded corners=4pt,minimum size=22pt,inner sep=0]{\popdisplay\fontsize{12}{12}\selectfont\color{popcream}\thepopsec};%
  \hspace{8pt}{\popdisplay\fontsize{15}{17}\selectfont\color{popink}\MakeUppercase{#1}}\par
  \vspace{4pt}\noindent{\color{poporange}\rule{36pt}{3.6pt}}
  \end{minipage}\par\nopagebreak\vspace{8pt}
}
\newcounter{popsub}
\newcommand{\courssub}[1]{%
  \stepcounter{popsub}%
  \par\vspace{9pt}\noindent{\popxbold\fontsize{11.5}{13}\selectfont\color{popink}{\color{poporange}\thepopsec.\thepopsub}\hspace{6pt}#1}\par\nopagebreak\vspace{4pt}
}

% ============================================================
% Boîtes pédagogiques — même recette : fond blanc, bordure épaisse colorée,
% ombre dure encre, badge plein. Titre optionnel [#1].
% ============================================================
\tcbset{popbase/.style={breakable,enhanced,colback=white,boxrule=2.3pt,arc=9pt,
  shadow={3pt}{-3pt}{0pt}{popink},left=14pt,right=14pt,top=11pt,bottom=11pt,before skip=11pt,after skip=15pt,
  fontupper=\popsemi\fontsize{10.6}{14.5}\selectfont\color{popink},
  fontlower=\popsemi\fontsize{10.6}{14.5}\selectfont\color{popink}}}
% Coche dessinée (le glyphe ✓ n'existe pas dans les polices du document)
\newcommand{\popcheck}[1]{\tikz[baseline=-0.5ex]\draw[#1,line width=1.6pt,line cap=round,line join=round](-0.13,0.0)--(-0.04,-0.09)--(0.13,0.1);}
\providecommand{\coche}{\popcheck{popgreen}}
\providecommand{\resultat}[1]{\par\smallskip\noindent\fcolorbox{popgreen}{popgreenL}{\parbox{\dimexpr\linewidth-2\fboxsep-2\fboxrule\relax}{\textbf{Résultat.} #1}}\par\smallskip}
\newcommand{\popboxhead}[3]{\poppill{#1}{#2}{white}\ifx\relax#3\relax\else\hspace{6pt}{\popxbold\fontsize{10.6}{12}\selectfont\color{popink}#3}\fi\par\vspace{7pt}}

\newtcolorbox{aretenir}[1][]{popbase,colframe=popgreen,before upper={\popboxhead{À retenir}{popgreen}{#1}}}
% Texte en chinois (document de lecture, dialogue, extrait) : cadre
% d'étude. Un titre optionnel (source, auteur) s'affiche dans l'en-tête.
\newtcolorbox{textebox}[1][]{popbase,colframe=popblue,before upper={\popboxhead{文本 · Texte}{popblue}{#1}},after upper={}}
% Marques ✓ / ✗ (forme correcte / fautive) souvent utilisées par les modèles
\providecommand{\cross}{\textcolor{poppink}{\ensuremath{\times}}}
\providecommand{\xmark}{\cross}\providecommand{\crossmark}{\cross}\providecommand{\cmark}{\ensuremath{\checkmark}}
% Ligne de réponse / justification à compléter (inventée par les modèles)
\providecommand{\justification}[1]{\par\noindent\textit{Justification~:} \dotfill\par}
% \textipa (package tipa non chargé) : la police ne couvre pas l'API -> simple passage
\providecommand{\textipa}[1]{#1}
% Soulignement (ulem non chargé) : \uline, \ul
\providecommand{\uline}[1]{\underline{#1}}\providecommand{\ul}[1]{\underline{#1}}
% \glossaire{\item ...} inventée par les modèles : liste de glossaire
\providecommand{\glossaire}[1]{\begin{itemize}#1\end{itemize}}
% \alert (beamer) inventée par les modèles : texte mis en évidence
\providecommand{\alert}[1]{\textbf{\textcolor{poppink}{#1}}}
% variantes ASCII d'ordinaux écrites en macro
\providecommand{\iere}{\textsuperscript{re}}\providecommand{\ieme}{\textsuperscript{e}}\providecommand{\ere}{\textsuperscript{re}}\providecommand{\eme}{\textsuperscript{e}}\providecommand{\er}{\textsuperscript{er}}\providecommand{\ie}{\textsuperscript{e}}
% Ordinaux français écrits en macro par les modèles : 1\ière, 3\ième
\newcommand{\ière}{\textsuperscript{re}}\newcommand{\ième}{\textsuperscript{e}}
% \exemple (inventée par les modèles) : étiquette « Exemple : »
\providecommand{\exemple}{\textbf{Exemple~:}\ }
% \square utilisable aussi hors mode math (cases à cocher)
\AtBeginDocument{\let\squareMath\square\renewcommand{\square}{\ensuremath{\squareMath}}}
% Mot ou phrase en chinois à l'intérieur d'un paragraphe français
\newcommand{\zh}[1]{\ifmmode\text{#1}\else{#1}\fi}
\let\eng\zh \let\esp\zh \let\all\zh % alias tolérés
% flèches d'intonation inventées par les modèles
\providecommand{\textsearrow}{}\renewcommand{\textsearrow}{\ensuremath{\searrow}}\providecommand{\textnearrow}{}\renewcommand{\textnearrow}{\ensuremath{\nearrow}}
% \fr{..} : mot français (inventé par les modèles)
\providecommand{\fr}[1]{\textit{#1}}
% zhbox : encadré de texte chinois inventé par les modèles
\newenvironment{zhbox}{\begin{quote}}{\end{quote}}
% \courssubsub : titre de niveau 3 inventé par les modèles
\providecommand{\courssubsub}[1]{\par\medskip\noindent\textbf{#1}\par\smallskip}
% \vs : « versus » inventé par les modèles
\providecommand{\vs}{\textit{vs}\ }
% \blank : case à compléter inventée par les modèles
\providecommand{\blank}[1][]{\underline{\hspace{1.6em}}}
\newtcolorbox{exemplebox}[1][]{popbase,colframe=poporange,before upper={\popboxhead{Exemple}{poporange}{#1}}}
\newtcolorbox{definition}[1][]{popbase,colframe=popblue,before upper={\popboxhead{Définition}{popblue}{#1}}}
\newtcolorbox{propriete}[1][]{popbase,colframe=poppurple,before upper={\popboxhead{Propriété}{poppurple}{#1}}}
\newtcolorbox{methode}[1][]{popbase,colframe=popgold,before upper={\popboxhead{Méthode}{popgold}{#1}}}
\newtcolorbox{attention}[1][]{popbase,colframe=poppink,before upper={\popboxhead{Attention}{poppink}{#1}}}
\newtcolorbox{experience}[1][]{popbase,colframe=popblue,colback=popblueL,before upper={\popboxhead{Expérience}{popblue}{#1}}}
% « Le savais-tu ? » : carte violette pleine (contexte, culture, Cameroun)
\newtcolorbox{savaistu}[1][]{popbase,colback=poppurple,colframe=popink,
  fontupper=\popsemi\fontsize{10.4}{14.2}\selectfont\color{white},
  before upper={\poppill{Le savais-tu\,?}{white}{poppurple}\ifx\relax#1\relax\else\hspace{6pt}{\popxbold\color{white}#1}\fi\par\vspace{7pt}}}
% Exercice résolu : énoncé en haut, \tcblower puis la solution sur fond crème
\newcounter{popexo}\newcounter{popexr}
\newtcolorbox{exoresolu}[1][]{popbase,colframe=poporange,colbacklower=popcream,
  segmentation style={popink,line width=1.2pt,dashed},
  before upper={\stepcounter{popexr}\popboxhead{Exercice résolu \thepopexr}{poporange}{#1}},
  before lower={\poppill{Solution}{popink}{popcream}\par\vspace{6pt}}}
% Exercice (non résolu en place) — la correction est dans la partie Corrigés
\newtcolorbox{exercice}[1][]{popbase,colframe=popink,
  before upper={\stepcounter{popexo}\popboxhead{Exercice \thepopexo}{popink}{#1}}}
% Corrigé
\newtcolorbox{corrige}[1][]{popbase,colframe=popgreen,colback=popgreenL,
  before upper={\popboxhead{Corrigé}{popgreen}{#1}}}
% Objectifs / prérequis / bilan
\newtcolorbox{objectifs}{popbase,colframe=popink,colback=poporangeL,
  before upper={\popboxhead{Objectifs de la leçon}{popink}{}}}
\newtcolorbox{prerequis}{popbase,colframe=popink,colback=popblueL,
  before upper={\popboxhead{Prérequis}{popblue}{}}}
\newtcolorbox{bilan}{popbase,colframe=popink,colback=white,boxrule=2.8pt,
  before upper={\popboxhead{Fiche bilan}{poporange}{L'essentiel à connaître pour l'examen}}}
% Figure encadrée avec légende
\newtcolorbox{popfigure}[1][]{enhanced,breakable=false,colback=white,colframe=popline,boxrule=1.4pt,arc=8pt,
  left=10pt,right=10pt,top=10pt,bottom=8pt,before skip=10pt,after skip=12pt,halign=center,
  lower separated=false,halign lower=center,
  fontlower=\popsemi\fontsize{9}{11.5}\selectfont\color{popmuted},
  #1}
\newcommand{\legende}[1]{\par\vspace{6pt}{\popsemi\fontsize{9}{11.5}\selectfont\color{popmuted}#1}}

% ============================================================
% Signature OnBuch+
% ============================================================
\newcommand{\poplogoinline}[1]{{\poplogo\fontsize{#1}{#1}\selectfont\color{popink}OnBuch\color{poporange}+}}
\newcommand{\signatureonbuch}{%
  \begin{tcolorbox}[enhanced,colback=popink,colframe=popink,boxrule=2.2pt,arc=10pt,
      drop shadow=poporange,left=14pt,right=14pt,top=10pt,bottom=10pt,before skip=0pt,after skip=0pt]
    \tikz[baseline=-0.7ex]\node[circle,fill=popgreen,draw=popcream,line width=1.3pt,minimum size=22pt,inner sep=0]{\popcheck{white}};%
    \hspace{9pt}\begin{minipage}[c]{0.62\linewidth}
      {\popxbold\fontsize{12}{13}\selectfont\color{popcream}Signé\ \textbullet\ {\poplogo OnBuch\color{poporange}+}}\par\vspace{2pt}
      {\raggedright\popsemi\fontsize{8}{10.5}\selectfont\color{popline}\MakeUppercase{Équipe pédagogique OnBuch+}\\\MakeUppercase{Conforme au programme officiel MINESEC}\par}
    \end{minipage}\hfill
    {\poplogo\fontsize{22}{22}\selectfont\color{popcream}OnBuch\color{poporange}+}
  \end{tcolorbox}%
}

% ============================================================
% Page de garde
% #1 titre · #2 matière · #3 niveau · #4 module · #5 n° leçon · #6 durée ·
% #7 description / objectifs (texte ou itemize)
% ============================================================
\newcommand{\pagegarde}[7]{%
  \thispagestyle{empty}
  \newgeometry{a4paper,margin=1.7cm,top=1.6cm,bottom=1.4cm}%
  \begin{tikzpicture}[remember picture,overlay]
    \fill[poporange!18] ([xshift=-3.2cm,yshift=-3.6cm]current page.north east) circle (4.6cm);
    \fill[poppurple!14] ([xshift=2.4cm,yshift=7.6cm]current page.south west) circle (3.8cm);
  \end{tikzpicture}%
  {\poplogo\fontsize{30}{30}\selectfont\color{popink}OnBuch\color{poporange}+}\hfill
  \raisebox{0.6ex}{\poppill{Cours signé \textbullet\ Premium}{popink}{popcream}}\par
  \vspace{1pt}\popsquiggle{poppurple}\par
  \vspace{8pt}
  \begin{tcolorbox}[enhanced,colback=poporange,colframe=popink,boxrule=2.8pt,arc=13pt,
      drop shadow=popink,left=18pt,right=18pt,top=12pt,bottom=13pt,after skip=10pt]
    \poppill{#2 \textbullet\ #3}{popink}{popcream}\hspace{5pt}\poppill{#4}{white}{popink}\par\vspace{8pt}
    {\popdisplay\fontsize{14}{14}\selectfont\color{popink}LEÇON #5}\par\vspace{4pt}
    {\raggedright\hyphenpenalty=10000\exhyphenpenalty=10000\popdisplay\fontsize{25}{27}\selectfont\color{popcream}\MakeUppercase{#1}\par}
    \vspace{8pt}
    {\popxbold\fontsize{9.5}{9.5}\selectfont\color{popink}\MakeUppercase{Durée conseillée : #6}}
  \end{tcolorbox}
  \begin{tcolorbox}[enhanced,colback=white,colframe=popink,boxrule=2.2pt,arc=10pt,
      drop shadow=popink,left=16pt,right=16pt,top=10pt,bottom=10pt,after skip=8pt]
    \poppill{Dans cette leçon}{poppurple}{white}\par\vspace{5pt}
    \popsemi\fontsize{9.3}{12.6}\selectfont\color{popink}#7
  \end{tcolorbox}
  \noindent\begin{minipage}[t]{0.487\linewidth}
    \begin{tcolorbox}[enhanced,colback=popgreen,colframe=popink,boxrule=2.2pt,arc=10pt,
        drop shadow=popink,left=11pt,right=11pt,top=10pt,bottom=10pt,height=3.1cm,valign=center]
      \tcbox[colback=white,colframe=popink,boxrule=1.3pt,arc=5pt,boxsep=0pt,nobeforeafter,
        left=3pt,right=3pt,top=3pt,bottom=3pt,valign=center]{\qrcode[height=1.9cm]{https://play.google.com/store/apps/details?id=com.luvvix.onbuch&hl=fr}}%
      \hspace{8pt}\begin{minipage}[c]{0.5\linewidth}
        {\popdisplay\fontsize{11}{12.5}\selectfont\color{white}\MakeUppercase{Télécharge OnBuch}}\par\vspace{4pt}
        {\raggedright\popsemi\fontsize{8.4}{11}\selectfont\color{white}Gratuit \textbullet\ Google Play\\Tuteur IA, annales, concours, résultats\par}
      \end{minipage}
    \end{tcolorbox}
  \end{minipage}\hfill
  \begin{minipage}[t]{0.487\linewidth}
    \begin{tcolorbox}[enhanced,colback=poppurple,colframe=popink,boxrule=2.2pt,arc=10pt,
        drop shadow=popink,left=11pt,right=11pt,top=10pt,bottom=10pt,height=3.1cm,valign=center]
      \tikz[baseline=-0.7ex]\node[circle,fill=white,draw=popink,line width=1.3pt,minimum size=1.6cm,inner sep=0]{\popdisplay\fontsize{24}{24}\selectfont\color{poppurple}+};%
      \hspace{8pt}\begin{minipage}[c]{0.6\linewidth}
        {\popdisplay\fontsize{11}{12.5}\selectfont\color{white}\MakeUppercase{Passe à OnBuch+}}\par\vspace{4pt}
        {\raggedright\popsemi\fontsize{8.4}{11}\selectfont\color{white}Tous les cours signés, Tuteur IA illimité, fiches PDF — onglet OnBuch+ de l'app\par}
      \end{minipage}
    \end{tcolorbox}
  \end{minipage}
  \vspace{8pt}
  \popcontenu
  \vfill
  \signatureonbuch
  \restoregeometry
  \newpage
}

% Rangée « Ce cours contient » (page de garde)
\newcommand{\poptile}[3]{% #1 couleur #2 gros texte #3 légende
  \begin{tcolorbox}[enhanced,colback=#1,colframe=popink,boxrule=2pt,arc=9pt,shadow={2.5pt}{-2.5pt}{0pt}{popink},
      left=6pt,right=6pt,top=9pt,bottom=9pt,halign=center,valign=center,height=1.75cm,width=\linewidth,nobeforeafter]
    {\popdisplay\fontsize{12.5}{13}\selectfont\color{white}\MakeUppercase{#2}}\par\vspace{3pt}
    {\centering\popsemi\fontsize{7.8}{9.5}\selectfont\color{white}#3\par}
  \end{tcolorbox}}
\newcommand{\popcontenu}{%
  \par\noindent{\popxboldls\fontsize{7.5}{7.5}\selectfont\color{popmuted}CE COURS SIGNÉ CONTIENT}\par\vspace{7pt}
  \noindent\begin{minipage}[t]{0.235\linewidth}\poptile{popblue}{Cours}{complet, illustré, pas à pas}\end{minipage}\hfill
  \begin{minipage}[t]{0.235\linewidth}\poptile{poporange}{Méthodes}{et exercices résolus}\end{minipage}\hfill
  \begin{minipage}[t]{0.235\linewidth}\poptile{poppink}{Exercices}{type examen + corrigés}\end{minipage}\hfill
  \begin{minipage}[t]{0.235\linewidth}\poptile{popgreen}{Fiche bilan}{l'essentiel à retenir}\end{minipage}\par}

% Sommaire
\newtcolorbox{sommaire}{enhanced,colback=white,colframe=popink,boxrule=2.2pt,arc=10pt,
  drop shadow=popink,left=18pt,right=18pt,top=13pt,bottom=13pt,before skip=6pt,after skip=20pt,
  before upper={\poppill{Sommaire}{poppurple}{white}\par\vspace{9pt}\popsemi\fontsize{10.6}{14.5}\selectfont\color{popink}}}

% Signature de fin de cours — poussée tout en bas de la dernière page
\newcommand{\findecours}{%
  \par\vfill\nopagebreak
  \begin{center}\popsquiggle{poporange}\end{center}\vspace{4pt}
  \signatureonbuch
}
\AtBeginDocument{\def\llbracket{\mathopen{[\mkern-3.5mu[}}\def\rrbracket{\mathclose{]\mkern-3.5mu]}}} % intervalles d'entiers (glyphe ⟦ absent de la police)
% Glyphes absents de Fira Math : redéfinis à partir de glyphes disponibles
\AtBeginDocument{%
  \def\setminus{\mathbin{\backslash}}\let\smallsetminus\setminus
  \def\vdots{\mathord{\vbox{\baselineskip3.2pt\lineskiplimit0pt\kern4pt\hbox{.}\hbox{.}\hbox{.}}}}%
  \def\ddots{\mathinner{\mkern1mu\raise6.5pt\hbox{.}\mkern2mu\raise3.6pt\hbox{.}\mkern2mu\raise0.7pt\hbox{.}\mkern1mu}}%
  \def\triangle{\mathord{\scalebox{0.9}{$\symup{\Delta}$}}}%
  \def\nabla{\mathord{\raisebox{\depth}{\scalebox{1}[-1]{$\symup{\Delta}$}}}}%
  \def\checkmark{\popcheck{popdark}}%
  \def\star{\mathbin{\ast}}\def\bigstar{\mathord{\ast}}%
  \def\diamond{\mathbin{\scalebox{0.6}{$\Diamond$}}}%
  \def\bigcup{\mathop{\mathchoice{\raisebox{-0.3em}{\scalebox{1.7}{$\cup$}}}{\scalebox{1.25}{$\cup$}}{\cup}{\cup}}\displaylimits}%
  \def\bigcap{\mathop{\mathchoice{\raisebox{-0.3em}{\scalebox{1.7}{$\cap$}}}{\scalebox{1.25}{$\cap$}}{\cap}{\cap}}\displaylimits}%
  \def\lesseqqgtr{\mathrel{\lessgtr}}\def\bigoplus{\mathop{\oplus}\displaylimits}%
}
\providecommand{\ssi}{si et seulement si} % abréviation fréquente
\AtBeginDocument{\providecommand{\wideparen}{\overparen}} % arc de cercle

%% FIGFIX-BEGIN v3 — correctifs globaux des figures (docs/onbuch-generation/tools/figfix)
\usepackage{adjustbox}\usepackage{etoolbox}
% toute figure TikZ/pgfplots plus large que la ligne est réduite à la largeur disponible
\BeforeBeginEnvironment{tikzpicture}{\begin{adjustbox}{max width=\linewidth}}
\AfterEndEnvironment{tikzpicture}{\end{adjustbox}}
% légendes pgfplots lisibles par-dessus les courbes
\pgfplotsset{every axis/.append style={legend style={fill=white,fill opacity=0.92,text opacity=1,draw=black!15,inner sep=3pt}}}
% étiquettes d'arêtes (syntaxe edge["..."]) avec fond blanc
\tikzset{every edge quotes/.append style={fill=white,inner sep=1.2pt}}
% bulles de cartes mentales : texte à la ligne dans la bulle, bulle un peu plus grande
\tikzset{every concept/.append style={text width=2.0cm,align=center,minimum size=2.3cm,inner sep=2pt}}
%% FIGFIX-END
\begin{document}

\pagegarde{Leçon 22 — Grammaire : phrase complexes}{Chinois}{Tle A}{Module 3 : Citoyenneté et ouverture au monde}{ZH22}{2 h}{Cette leçon de grammaire présente les principales phrases complexes du chinois : coordination, cause, concession, condition et but. À travers l'observation d'exemples traduits, des schémas de structures et une comparaison systématique avec le français, les élèves apprennent à construire des phrases à deux propositions correctes et naturelles.
\vspace{4pt}
\begin{multicols}{2}\small
\begin{itemize}
\item À la fin de cette leçon, je sais identifier une phrase complexe et ses deux propositions en chinois
\item À la fin de cette leçon, je sais exprimer la cause avec 因为...所以... et le but avec 为了
\item À la fin de cette leçon, je sais exprimer la concession avec 虽然...但是...
\item À la fin de cette leçon, je sais exprimer la condition avec 如果...就... et 要是...就...
\item À la fin de cette leçon, je sais placer correctement les corrélats et le sujet dans chaque structure
\item À la fin de cette leçon, je sais éviter les erreurs typiques des francophones (double sujet, ordre des mots, traduction mot à mot)
\end{itemize}\end{multicols}}

\begin{sommaire}
\begin{multicols}{2}
\begin{enumerate}
\item Observation et notion de phrase complexe
\item Cause et conséquence : 因为...所以...
\item Concession : 虽然...但是...
\item Condition et hypothèse : 如果...就... / 要是...就...
\item But : 为了... et coordination simple
\item Synthèse, pièges des francophones et entraînement
\item Activité d'intégration
\item Méthodes et exercices résolus
\item Exercices
\item Corrigés des exercices
\item Fiche bilan
\end{enumerate}
\end{multicols}
\end{sommaire}

\begin{objectifs}
\begin{itemize}
\item À la fin de cette leçon, je sais identifier une phrase complexe et ses deux propositions en chinois
\item À la fin de cette leçon, je sais exprimer la cause avec 因为...所以... et le but avec 为了
\item À la fin de cette leçon, je sais exprimer la concession avec 虽然...但是...
\item À la fin de cette leçon, je sais exprimer la condition avec 如果...就... et 要是...就...
\item À la fin de cette leçon, je sais placer correctement les corrélats et le sujet dans chaque structure
\item À la fin de cette leçon, je sais éviter les erreurs typiques des francophones (double sujet, ordre des mots, traduction mot à mot)
\item À la fin de cette leçon, je sais transformer des phrases simples en phrases complexes et inversement
\item À la fin de cette leçon, je sais employer ces structures dans un court texte cohérent (lettre, récit)
\end{itemize}
\end{objectifs}

\begin{prerequis}
\begin{itemize}
\item Phrase simple chinoise : ordre Sujet–Verbe–Objet et place des compléments (leçons de Première)
\item Adverbes de liaison déjà vus : 也, 都, 就, 才
\item Négation 不 / 没 et particule 了
\item Vocabulaire courant : école, famille, voyages, météo, études
\item Notion de conjonction en français (parce que, mais, si, bien que)
\end{itemize}
\end{prerequis}

\newpage

\coursec{Observation et notion de phrase complexe}

Pendant les grandes vacances, Amina, élève de Terminale A à Yaoundé, s'installe à sa table pour écrire à Li Ming, son correspondant chinois de Shanghai. Elle veut lui raconter sa semaine : parce qu'il a beaucoup plu, elle n'a pas pu aller au village voir sa grand-mère ; bien qu'elle soit un peu déçue, elle en a profité pour réviser son chinois ; si le temps s'améliore, elle fera le voyage la semaine prochaine ; et pour réussir son examen, elle s'entraîne chaque jour. En français, tout cela tient en quelques phrases grâce à « parce que », « mais », « si », « pour »... Mais comment Amina dira-t-elle cela en chinois ? Comment relier deux idées dans une seule phrase ?

\textbf{Problématique :} comment le chinois construit-il la \cle{phrase complexe} pour exprimer la cause, la concession, la condition et le but ? Cette première partie de la leçon pose les fondations : observer, définir, repérer les \cle{corrélats} et comprendre la logique des deux propositions.

\courssub{Première observation : quatre phrases modèles}

Lisons attentivement les quatre phrases suivantes. Chacune est formée de \textbf{deux parties} séparées par une virgule chinoise «，». Surlignons mentalement les mots qui relient ces deux parties.

\begin{exemplebox}[Quatre phrases modèles à observer]
\begin{enumerate}
\item \zh{因为下雨，所以我没去。}\\
\emph{Yīnwèi xiàyǔ, suǒyǐ wǒ méi qù.}\\
« Parce qu'il pleut, je ne suis pas allé. »
\item \zh{虽然很累，但是我很开心。}\\
\emph{Suīrán hěn lèi, dànshì wǒ hěn kāixīn.}\\
« Bien que je sois fatigué, je suis très content. »
\item \zh{如果你努力，你就会进步。}\\
\emph{Rúguǒ nǐ nǔlì, nǐ jiù huì jìnbù.}\\
« Si tu travailles dur, tu progresseras. »
\item \zh{为了学好中文，我每天练习。}\\
\emph{Wèile xuéhǎo Zhōngwén, wǒ měitiān liànxí.}\\
« Pour bien apprendre le chinois, je m'entraîne chaque jour. »
\end{enumerate}
\end{exemplebox}

\begin{experience}[Observe et réponds]
En binôme, répondez à ces questions avant de lire la suite :
\begin{enumerate}
\item Combien de « petites phrases » (propositions) compte chaque phrase modèle ?
\item Quels mots se répètent en tête de chaque proposition ? Soulignez-les.
\item Dans la phrase 1, quelle partie donne la cause ? Quelle partie donne le résultat ?
\item Comparez avec le français : la traduction française utilise-t-elle deux mots de liaison ou un seul ?
\end{enumerate}
\end{experience}

\courssub{Le texte d'Amina : observation en contexte}

Voici le début de la lettre qu'Amina finira par écrire. Observons comment les phrases complexes s'y enchaînent naturellement.

\begin{textebox}[La lettre d'Amina à Li Ming (extrait)]
因为今天下雨，所以我不去学校。\\
\emph{Yīnwèi jīntiān xiàyǔ, suǒyǐ wǒ bú qù xuéxiào.}\\
« Parce qu'il pleut aujourd'hui, je ne vais pas à l'école. »\\[0.4em]
虽然我很失望，但是我在家学习中文。\\
\emph{Suīrán wǒ hěn shīwàng, dànshì wǒ zài jiā xuéxí Zhōngwén.}\\
« Bien que je sois déçue, j'étudie le chinois à la maison. »\\[0.4em]
如果明天天气好，我就去村里看奶奶。\\
\emph{Rúguǒ míngtiān tiānqì hǎo, wǒ jiù qù cūnlǐ kàn nǎinai.}\\
« Si le temps est beau demain, j'irai voir ma grand-mère au village. »\\[0.4em]
为了说好中文，我每天练习一个小时。\\
\emph{Wèile shuōhǎo Zhōngwén, wǒ měitiān liànxí yí ge xiǎoshí.}\\
« Pour bien parler chinois, je m'entraîne une heure par jour. »
\end{textebox}

\begin{center}
\begin{tabularx}{\linewidth}{|X|X|X|X|}\hline
\rowcolor{popblueL} Caractères & Pinyin & Nature & Traduction \\ \hline
失望 & shīwàng & adj. / v. & déçu ; être déçu \\ \hline
天气 & tiānqì & n. & le temps (météo) \\ \hline
村 & cūn & n. & village \\ \hline
奶奶 & nǎinai & n. & grand-mère (paternelle) \\ \hline
小时 & xiǎoshí & n. & heure (durée) \\ \hline
进步 & jìnbù & v. / n. & progresser ; progrès \\ \hline
\end{tabularx}
\end{center}

\courssub{Définition de la phrase complexe}

\begin{definition}[Phrase complexe]
Une \cle{phrase complexe} est une phrase formée de \textbf{deux propositions} (P1 et P2), c'est-à-dire deux groupes de mots construits chacun autour d'un verbe, reliées par des \cle{corrélats} (mots-outils appariés) qui expriment un \textbf{rapport logique} : la \cle{cause}, la \cle{concession}, la \cle{condition} ou le \cle{but}.
\end{definition}

Décortiquons la première phrase modèle :

\begin{center}
\zh{因为下雨} \;，\; \zh{所以我没去。}\\
P1 : \zh{因为} + cause (il pleut) \quad $\longrightarrow$ \quad P2 : \zh{所以} + conséquence (je ne suis pas allé)
\end{center}

\begin{itemize}
\item \textbf{P1} (proposition subordonnée) : \zh{因为下雨} « parce qu'il pleut » — elle donne la \textbf{cause}, elle dépend de P2 pour avoir son plein sens.
\item \textbf{P2} (proposition principale) : \zh{所以我没去} « donc je ne suis pas allé » — elle porte l'\textbf{information principale}, le résultat.
\end{itemize}

\begin{definition}[Corrélat]
Un \cle{corrélat} est un mot-outil (conjonction ou adverbe de liaison) qui travaille \textbf{en paire} avec un autre pour relier deux propositions : \zh{因为} appelle \zh{所以}, \zh{虽然} appelle \zh{但是}, \zh{如果} appelle \zh{就}. On parle de \textbf{corrélats appariés}. \textit{Note : pour le but, \zh{为了} fonctionne seul en tête de P1.}
\end{definition}

\courssub{Schéma de la structure : P1 ， P2}

\begin{popfigure}
\begin{tikzpicture}[node distance=1.1cm and 1.6cm]
\node[draw=popblue, fill=popblueL, rounded corners, thick, minimum width=4.6cm, minimum height=1.5cm, align=center] (p1) {\textbf{P1} : \zh{因为} + cause\\\zh{因为下雨}\\\emph{parce qu'il pleut}};
\node[draw=poporange, fill=poporangeL, rounded corners, thick, minimum width=4.6cm, minimum height=1.5cm, align=center, right=of p1] (p2) {\textbf{P2} : \zh{所以} + conséquence\\\zh{所以我没去}\\\emph{donc je ne suis pas allé}};
\draw[-{Stealth[length=3mm]}, very thick, popdark] (p1) -- (p2) node[midway, above, align=center, font=\small] {，\\rapport logique};
\node[below=0.35cm of p1, font=\small\itshape, align=center, popmuted] {subordonnée\\(cause)};
\node[below=0.35cm of p2, font=\small\itshape, align=center, popmuted] {principale\\(résultat)};
\end{tikzpicture}
\legende{Structure générale de la phrase complexe de cause : P1 (因为 + cause) ， P2 (所以 + conséquence). La flèche matérialise le rapport logique : la cause entraîne la conséquence.}
\end{popfigure}

\begin{propriete}[Schéma général de la phrase complexe]
\textbf{Corrélat 1 + P1 ， Corrélat 2 + P2 。}
\begin{itemize}
\item La \textbf{virgule chinoise «，»} sépare obligatoirement les deux propositions.
\item L'ordre le plus fréquent est : \textbf{subordonnée d'abord, principale ensuite} (cause $\rightarrow$ résultat ; condition $\rightarrow$ conséquence).
\item Le sujet, s'il est le même dans les deux propositions, est souvent placé \textbf{avant le premier corrélat} ou \textbf{après celui-ci}, mais il n'est pas répété inutilement : \zh{我因为下雨，所以没去。} ou \zh{因为下雨，所以我没去。}
\end{itemize}
\end{propriete}

\courssub{Les quatre rapports logiques de la leçon}

\begin{popfigure}
\begin{tikzpicture}[mindmap, grow cyclic, every node/.style={concept, align=center, font=\small, text width=2.2cm},
root concept/.append style={concept color=popblue, fill=popblueL, font=\small\bfseries, text width=3cm},
level 1/.append style={level distance=3.8cm, sibling angle=90},
level 2/.append style={level distance=2.8cm, sibling angle=50, font=\footnotesize}]
\node[root concept, align=center]{Phrase complexe\\P1 ， P2}
  child[concept color=popgreen, fill=popgreenL]{ node[align=center]{Cause\\\zh{因为…所以…}\\\emph{parce que… donc…}} }
  child[concept color=poporange, fill=poporangeL]{ node[align=center]{Concession\\\zh{虽然…但是…}\\\emph{bien que… mais…}} }
  child[concept color=poppurple, fill=poppurpleL]{ node[align=center]{Condition\\\zh{如果…就…}\\\emph{si… alors…}} }
  child[concept color=poppink, fill=poppinkL]{ node[align=center]{But\\\zh{为了…}\\\emph{pour…}} };
\end{tikzpicture}
\legende{Carte mentale des quatre rapports logiques exprimés par la phrase complexe dans cette leçon.}
\end{popfigure}

\begin{center}
\begin{tabularx}{\linewidth}{|X|X|X|X|}\hline
\rowcolor{popblueL} Rapport logique & Corrélats chinois & Exemple & Traduction \\ \hline
Cause & \zh{因为…所以…} \emph{yīnwèi… suǒyǐ…} & \zh{因为下雨，所以我没去。} & Parce qu'il pleut, je ne suis pas allé. \\ \hline
Concession & \zh{虽然…但是…} \emph{suīrán… dànshì…} & \zh{虽然很累，但是我很开心。} & Bien que fatigué, je suis content. \\ \hline
Condition & \zh{如果…就…} \emph{rúguǒ… jiù…} & \zh{如果你努力，你就会进步。} & Si tu travailles, tu progresseras. \\ \hline
But & \zh{为了…} \emph{wèile…} & \zh{为了学好中文，我每天练习。} & Pour apprendre le chinois, je m'entraîne chaque jour. \\ \hline
\end{tabularx}
\end{center}

\courssub{Comparaison avec le français : deux corrélats contre un seul}

C'est ici que se situe la différence la plus importante entre les deux langues, et la source de la plupart des erreurs des élèves francophones.

\begin{propriete}[Les deux corrélats apparaissent ensemble (Cause, Concession, Condition)]
En chinois, pour la cause, la concession et la condition, les deux corrélats apparaissent \textbf{souvent ensemble} dans la même phrase : \zh{因为…所以…}, \zh{虽然…但是…}, \zh{如果…就…}. Le français, lui, \textbf{interdit} cette double marque : on dit « parce qu'il pleut, je ne sors pas » et jamais « parce qu'il pleut, donc je ne sors pas » (tournure fautive en français soigné). Le chinois, au contraire, trouve naturel et clair de marquer les deux bords du rapport logique.
\end{propriete}

\begin{center}
\begin{tabularx}{\linewidth}{|X|X|X|}\hline
\rowcolor{popblueL} Rapport & Chinois (deux corrélats possibles) & Français (un seul mot de liaison) \\ \hline
Cause & \zh{因为} P1 ，\zh{所以} P2 & parce que P1, P2 \\ \hline
Concession & \zh{虽然} P1 ，\zh{但是} P2 & bien que P1, P2 \\ \hline
Condition & \zh{如果} P1 ，P2 \zh{就}… & si P1, P2 \\ \hline
But & \zh{为了} P1 ，P2 & pour P1, P2 \\ \hline
\end{tabularx}
\end{center}

\begin{exemplebox}[Exemples traduits : observez la double marque]
\begin{enumerate}
\item \zh{虽然他很忙，但是他帮助我。}\\
\emph{Suīrán tā hěn máng, dànshì tā bāngzhù wǒ.}\\
« Bien qu'il soit occupé, il m'aide. » — Le français ne garde que « bien que » ; le chinois garde \zh{虽然} ET \zh{但是}.
\item \zh{如果你有时间，我们就去。}\\
\emph{Rúguǒ nǐ yǒu shíjiān, wǒmen jiù qù.}\\
« Si tu as le temps, nous irons. » — Le français ne dit pas « si... alors nous irons » dans la langue courante ; le chinois emploie \zh{如果} ET \zh{就}.
\item \zh{因为他生病了，所以今天没来上课。}\\
\emph{Yīnwèi tā shēngbìng le, suǒyǐ jīntiān méi lái shàngkè.}\\
« Parce qu'il est malade, il n'est pas venu en cours aujourd'hui. »
\end{enumerate}
\end{exemplebox}

\begin{attention}[Erreur n° 1 des francophones : calquer le français]
Sous l'influence du français, les élèves ont tendance à n'utiliser qu'\textbf{un seul} corrélat : \zh{因为下雨，我没去。} Cette phrase n'est pas impossible à l'oral familier, mais à l'écrit et à l'examen, la norme attendue est la \textbf{paire complète} : \zh{因为下雨，所以我没去。} Inversement, en traduisant vers le français, il faut \textbf{supprimer} un des deux corrélats : \zh{虽然…但是…} se traduit « bien que... » seul, jamais « bien que... mais... ».
\end{attention}

\courssub{La ponctuation : la virgule chinoise «，»}

\begin{propriete}[Virgule obligatoire entre les deux propositions]
Entre P1 et P2, on écrit toujours la \textbf{virgule chinoise pleine chasse «，»}. Elle marque nettement la pause entre les deux propositions. La phrase se termine par le point chinois «。». On n'utilise jamais la virgule française « , » dans un texte chinois.
\end{propriete}

\begin{savaistu}[Une virgule plus large]
La virgule chinoise «，» est typographiquement \textbf{plus large} que la virgule française : elle occupe la largeur d'un caractère entier (on dit « pleine chasse »). À l'oral, elle correspond à une vraie pause, un temps d'arrêt qui laisse l'auditeur enregistrer la première idée avant d'entendre la seconde. Il existe aussi la virgule chinoise «、» (dùnhào), réservée aux énumérations de mots : \zh{我喜欢苹果、香蕉和橙子。} « J'aime les pommes, les bananes et les oranges. » Ne confondez pas les deux !
\end{savaistu}

\courssub{Proposition principale et proposition subordonnée : le sens logique}

Comment distinguer la proposition principale de la subordonnée ? Deux tests simples :

\begin{methode}[Identifier P1 et P2]
\begin{enumerate}
\item \textbf{Repérer les corrélats.} \zh{因为}, \zh{虽然}, \zh{如果}, \zh{为了} introduisent la subordonnée (P1) ; \zh{所以}, \zh{但是}, \zh{就} introduisent la principale (P2).
\item \textbf{Poser la question du sens.} Demandez-vous : « quelle est l'information essentielle ? » Dans \zh{因为下雨，所以我没去。}, l'essentiel est « je ne suis pas allé » (P2) ; « il pleut » (P1) ne fait qu'expliquer pourquoi.
\item \textbf{Vérifier l'ordre logique.} Le chinois suit l'ordre du raisonnement : d'abord la cause ou la condition, ensuite le résultat. C'est l'ordre « parce que... donc... », très naturel.
\end{enumerate}
\end{methode}

\begin{exemplebox}[Analyse guidée de trois phrases]
\begin{enumerate}
\item \zh{虽然杜阿拉很热，但是我喜欢这个城市。}\\
\emph{Suīrán Dù'ālā hěn rè, dànshì wǒ xǐhuan zhège chéngshì.}\\
« Bien que Douala soit chaude, j'aime cette ville. »\\
P1 subordonnée : \zh{虽然杜阿拉很热} (concession) ; P2 principale : \zh{但是我喜欢这个城市} (l'essentiel : j'aime la ville).
\item \zh{如果明天下雨，我们就不去雅温得。}\\
\emph{Rúguǒ míngtiān xiàyǔ, wǒmen jiù bú qù Yǎwēndé.}\\
« S'il pleut demain, nous n'irons pas à Yaoundé. »\\
P1 subordonnée : la condition ; P2 principale : la conséquence.
\item \zh{为了通过考试，她每天学习两个小时。}\\
\emph{Wèile tōngguò kǎoshì, tā měitiān xuéxí liǎng ge xiǎoshí.}\\
« Pour réussir l'examen, elle étudie deux heures par jour. »\\
P1 subordonnée : le but ; P2 principale : l'action réelle.
\end{enumerate}
\end{exemplebox}

\courssub{Exercice résolu et entraînement}

\begin{exoresolu}[Identifier et compléter]
\textbf{Énoncé.} 1) Dans la phrase \zh{因为天气不好，所以我们不去市场。}, indiquez P1, P2, les corrélats et le rapport logique. 2) Complétez avec le corrélat manquant : \zh{虽然他很累，} \_\_\_\_ \zh{他还帮助我学习。}
\tcblower
\textbf{Solution détaillée.}
\begin{enumerate}
\item Corrélats : \zh{因为} (P1) et \zh{所以} (P2). P1 subordonnée : \zh{因为天气不好} « parce que le temps est mauvais » (cause) ; P2 principale : \zh{所以我们不去市场} « donc nous n'allons pas au marché » (conséquence). Rapport logique : la \textbf{cause}. Traduction : « Parce que le temps est mauvais, nous n'allons pas au marché. »
\item La phrase commence par \zh{虽然} « bien que » : le corrélat apparié attendu est \zh{但是} « mais ». Phrase complète : \zh{虽然他很累，但是他还帮助我学习。} \emph{Suīrán tā hěn lèi, dànshì tā hái bāngzhù wǒ xuéxí.} « Bien qu'il soit fatigué, il m'aide quand même à étudier. » Rapport logique : la \textbf{concession}.
\end{enumerate}
\end{exoresolu}

\begin{exercice}[À vous de jouer]
\begin{enumerate}
\item Soulignez les corrélats et traduisez : \zh{如果你有时间，我们就去看足球比赛。}
\item Complétez : \zh{因为} \_\_\_\_\zh{，} \zh{所以我不去村里。} (cause : il pleut)
\item Classez ces phrases selon le rapport logique (cause / concession / condition / but) : a) \zh{为了学好中文，我每天练习。} b) \zh{虽然很难，但是我很喜欢。} c) \zh{因为下雨，所以我不去。}
\end{enumerate}
\end{exercice}

\begin{corrige}[Exercice « À vous de jouer »]
\begin{enumerate}
\item Corrélats : \zh{如果} et \zh{就}. Traduction : « Si tu as le temps, nous irons voir le match de football. » (\zh{足球比赛} \emph{zúqiú bǐsài} : match de football.)
\item \zh{因为下雨，所以我不去村里。} \emph{Yīnwèi xiàyǔ, suǒyǐ wǒ bú qù cūnlǐ.} « Parce qu'il pleut, je ne vais pas au village. »
\item a) but (\zh{为了}) ; b) concession (\zh{虽然…但是…}) ; c) cause (\zh{因为…所以…}).
\end{enumerate}
\end{corrige}

\begin{aretenir}[L'essentiel de la section]
\begin{itemize}
\item Une \cle{phrase complexe} = P1 ，P2, deux propositions reliées par des \cle{corrélats appariés} exprimant un rapport logique.
\item Les quatre rapports de la leçon : cause \zh{因为…所以…}, concession \zh{虽然…但是…}, condition \zh{如果…就…}, but \zh{为了…}.
\item Le chinois garde souvent \textbf{deux} corrélats là où le français n'en garde qu'\textbf{un} (sauf pour le but) : ne calquez jamais le français !
\item La \textbf{virgule chinoise «，»} sépare toujours les deux propositions ; le point «。» termine la phrase.
\item Ordre logique habituel : subordonnée (cause, condition, concession, but) d'abord, principale (résultat) ensuite.
\end{itemize}
\end{aretenir}

Amina a maintenant tous les outils pour écrire sa lettre. Dans la prochaine section, nous étudierons en détail la première paire de corrélats : \zh{因为…所以…}, la phrase de \textbf{cause et conséquence}.

\coursec{Cause et conséquence : 因为…所以…}

\courssub{Structure détaillée et variantes}

La structure canonique est \zh{因为} + Cause ， \zh{所以} + Conséquence. Le mot \zh{因为} (yīnwèi) est une conjonction « parce que » ; \zh{所以} (suǒyǐ) est un adverbe « donc, par conséquent ».

\begin{propriete}[Variantes de registre]
\begin{itemize}
\item \textbf{Standard (écrit / formel) :} \zh{因为…所以…} / \zh{由于…因此…} (\emph{yóuyú… yīncǐ…}) — \zh{由于} est plus formel, souvent utilisé dans les textes administratifs, journalistiques ou académiques.
\item \textbf{Oral / familier :} \zh{因为…} seul (P2 sans \zh{所以}) ou \zh{…所以…} seul (P1 sans \zh{因为}, contexte implicite).
\item \textbf{Accentuation (focus) :} \zh{之所以…是因为…} (\emph{zhīsuǒyǐ… shì yīnwèi…}) « La raison pour laquelle..., c'est parce que... ». Structure d'insistance sur la cause.
\end{itemize}
\end{propriete}

\begin{exemplebox}[Cause : exemples variés]
\begin{enumerate}
\item \zh{由于交通堵塞，因此我迟到了。}\\
\emph{Yóuyú jiāotōng dǔsè, yīncǐ wǒ chídào le.}\\
« En raison des embouteillages, j'ai été en retard. » (Formel)
\item \zh{他之所以生气，是因为你迟到了。}\\
\emph{Tā zhīsuǒyǐ shēngqì, shì yīnwèi nǐ chídào le.}\\
« La raison pour laquelle il est en colère, c'est que tu es en retard. » (Insistance)
\item \zh{因为想吃烤鱼，所以我们去了夜市。}\\
\emph{Yīnwèi xiǎng chī kǎoyú, suǒyǐ wǒmen qù le yèshì.}\\
« Parce qu'on voulait manger du poisson grillé, on est allés au marché de nuit. »
\end{enumerate}
\end{exemplebox}

\courssub{Négation et cause}

La négation porte généralement sur la conséquence (P2) avec \zh{不} ou \zh{没}. Pour nier la cause, on utilise \zh{不是因为…而是因为…} « ce n'est pas parce que... mais parce que... ».

\begin{exemplebox}[Négation]
\begin{enumerate}
\item \zh{因为不舒服，所以他没来。} (Cause positive, conséquence négative) — « Parce qu'il ne se sent pas bien, il n'est pas venu. »
\item \zh{他不是因为懒，而是因为病了，所以没来。}\\
\emph{Tā bù shì yīnwèi lǎn, ér shì yīnwèi bìng le, suǒyǐ méi lái.}\\
« Ce n'est pas par paresse, mais parce qu'il est malade, qu'il n'est pas venu. »
\end{enumerate}
\end{exemplebox}

\courssub{Exercices d'application : Cause}

\begin{exercice}[Cause : transformation et production]
\begin{enumerate}
\item Reliez les deux phrases par \zh{因为…所以…} : \zh{天气很冷。} / \zh{我们不去爬山。} (\zh{爬山} \emph{páshān} : faire de la randonnée / grimper).
\item Traduisez en chinois : « Puisqu'il a raté le bus (错过公交车 \emph{cuòguò gōngjiāochē}), il a pris un taxi (打车 \emph{dǎchē}). »
\item Réécrivez avec \zh{由于…因此…} : \zh{因为下大雨，所以比赛取消了。} (\zh{取消} \emph{qǔxiāo} : annuler).
\item Complétez : \zh{他\_\_\_\_生气，\_\_\_\_是因为你弄坏了他的手机。} (\zh{弄坏} \emph{nònghuài} : casser/abîmer).
\end{enumerate}
\end{exercice}

\begin{corrige}[Exercice Cause]
\begin{enumerate}
\item \zh{因为天气很冷，所以我们不去爬山。}
\item \zh{因为他错过公交车，所以他打车了。} / \zh{由于他错过公交车，因此他打车了。}
\item \zh{由于下大雨，因此比赛取消了。}
\item \zh{他之所以生气，是因为你弄坏了他的手机。}
\end{enumerate}
\end{corrige}

\coursec{Concession : 虽然…但是…}

\courssub{Structure et nuances}

La structure canonique : \zh{虽然} (suīrán) + Fait admis ， \zh{但是} (dànshì) + Fait principal (contraire à l'attente).
\zh{虽然} = « bien que, même si » (conjonction). \zh{但是} = « mais, cependant » (adverbe de transition).

\begin{propriete}[Variantes des corrélats en P2]
En P2, \zh{但是} peut être remplacé selon la nuance :
\begin{itemize}
\item \zh{可是} (kěshì) : plus oral, plus léger, « mais ».
\item \zh{不过} (búguò) : « toutefois, simplement », introduit souvent une restriction ou une précision.
\item \zh{却} (què) : adverbe marquant le contraste fort, placé \textbf{avant le verbe} (pas en tête de proposition). Ex : \zh{虽然下雨，他却去了。}
\item \zh{仍然} (réngrán) / \zh{依然} (yīrán) : « encore, toujours », insiste sur la persistance de l'état/action malgré la concession. Souvent combiné : \zh{但是仍然…}
\end{itemize}
\end{propriete}

\begin{propriete}[Concession forte : 即使…也… / 就算…也…]
Pour exprimer « même si / quand bien même » (hypothèse concessionnelle), on utilise :
\zh{即使} (jíshǐ) / \zh{就算} (jiùsuàn) + Condition hypothétique ， \zh{也} (yě) / \zh{还} (hái) + Verbe.
Ex : \zh{即使下雨，我也要去。} « Même s'il pleut, j'irai. »
\end{propriete}

\begin{exemplebox}[Concession : exemples contextuels]
\begin{enumerate}
\item \zh{虽然这道菜很辣，但是很好吃。}\\
\emph{Suīrán zhège cài hěn là, dànshì hěn hǎochī.}\\
« Bien que ce plat soit épicé, il est délicieux. »
\item \zh{他虽然老了，可是身体很硬朗。}\\
\emph{Tā suīrán lǎo le, kěshì shēntǐ hěn yìnglǎng.}\\
« Bien qu'il soit vieux, il est en pleine forme. » (Sujet avant \zh{虽然}).
\item \zh{即使考试很难，我们也要努力。}\\
\emph{Jíshǐ kǎoshì hěn nán, wǒmen yě yào nǔlì.}\\
« Même si l'examen est difficile, nous devons faire des efforts. »
\item \zh{虽然路远，他却坚持走路回家。}\\
\emph{Suīrán lù yuǎn, tā què jiānchí zǒulù huí jiā.}\\
« Bien que la route soit longue, il a insisté pour rentrer à pied. » (\zh{却} avant \zh{坚持}).
\end{enumerate}
\end{exemplebox}

\courssub{Erreurs fréquentes : place de \zh{但是} et \zh{却}}

\begin{attention}[Pièges concession]
\begin{enumerate}
\item \textbf{Oublier \zh{但是} (ou variante) en P2.} \textit{Faux :} \zh{虽然很累，我很快乐。} \textit{Correct :} \zh{虽然很累，但是我很快乐。}
\item \textbf{Placer \zh{却} en tête de proposition.} \textit{Faux :} \zh{虽然下雨，却他去了。} \textit{Correct :} \zh{虽然下雨，他却去了。} (\zh{却} est un adverbe, pas une conjonction).
\item \textbf{Confondre \zh{虽然} et \zh{即使}.} \zh{虽然} = fait réel (il pleut vraiment). \zh{即使} = hypothétique (même s'il pleuvait / pleut).
\end{enumerate}
\end{attention}

\courssub{Exercices d'application : Concession}

\begin{exercice}[Concession : choix du corrélat et traduction]
\begin{enumerate}
\item Choisissez le bon mot pour P2 (\zh{但是} / \zh{可是} / \zh{不过} / \zh{却}) : \zh{虽然今天是星期天，\_\_\_\_我还要上课。} (Insistance sur le fait inattendu).
\item Traduisez : « Bien que le ndolé soit bon, je n'en mange pas tous les jours. » (\zh{恩多莱} \emph{ēnduōlái} ndolé, \zh{天天} \emph{tiāntiān} tous les jours).
\item Transformez en concession forte (\zh{即使…也…}) : \zh{虽然下雨，我们也去踢球。} (\zh{踢球} \emph{tīqiú} jouer au foot).
\item Complétez : \zh{这个手机\_\_\_\_贵，\_\_\_\_性能很好。} (Bien que cher, performance top).
\end{enumerate}
\end{exercice}

\begin{corrige}[Exercice Concession]
\begin{enumerate}
\item \zh{却} (ou \zh{但是} / \zh{还是} selon nuance, mais \zh{却} marque le contraste fort attendu). \zh{虽然今天是星期天，可/却/但是我还要上课。}
\item \zh{虽然恩多莱很好吃，但是我不天天吃。} / \zh{尽管恩多莱好吃，我不过天天吃。}
\item \zh{即使下雨，我们也要去踢球。}
\item \zh{虽然这个手机很贵，但是/不过性能很好。}
\end{enumerate}
\end{corrige}

\coursec{Condition et hypothèse : 如果…就… / 要是…就…}

\courssub{Structure canonique : 如果…就…}

\zh{如果} (rúguǒ) « si » introduit la condition (P1). \zh{就} (jiù) « alors, dans ce cas » marque la conséquence (P2), placé \textbf{immédiatement avant le verbe} (ou l'auxiliaire modal) de P2.

\begin{propriete}[La place cruciale de \zh{就}]
\zh{就} n'est pas une conjonction de liaison comme \zh{所以} ou \zh{但是}. C'est un \textbf{adverbe de portée} qui doit précéder le prédicat de la principale.
\begin{itemize}
\item \textit{Correct :} \zh{如果下雨，我就不去。} (\zh{就} + \zh{不去})
\item \textit{Faux :} \zh{如果下雨，就我不去。} (Sujet intercalé entre \zh{就} et le verbe).
\end{itemize}
Si le sujet de P2 est différent de P1, \zh{just} se place \textbf{après le sujet de P2} : \zh{如果下雨，比赛就取消。}
\end{propriete}

\courssub{Variantes selon le registre et le type de condition}

\begin{center}
\begin{tabularx}{\linewidth}{|X|X|X|X|}\hline
\rowcolor{popblueL} Type & P1 (Condition) & P2 (Conséquence) & Exemple \\ \hline
\textbf{Réelle / Future} & \zh{如果} / \zh{要是} (oral) & \zh{就} / \zh{会} / \zh{要} / \zh{得} & \zh{如果明天晴天，我们就去野餐。} \\ \hline
\textbf{Hypothétique / Contrefactuelle} & \zh{假如} / \zh{假使} (formel) & \zh{就} / \zh{会} / \zh{早就} & \zh{假如我是你，我就早走了。} \\ \hline
\textbf{Condition nécessaire} & \zh{只有} (zhǐyǒu) & \zh{才} (cái) & \zh{只有努力，才能成功。} \\ \hline
\textbf{Condition négative (sauf si)} & \zh{除非} (chúfēi) & \zh{否则} (fǒuzé) / \zh{就} & \zh{除非你邀请我，否则我不去。} \\ \hline
\end{tabularx}
\end{center}

\begin{exemplebox}[Condition : exemples progressifs]
\begin{enumerate}
\item \textbf{Real future :} \zh{要是你有空，给我打电话。} \emph{Yàoshi nǐ yǒukòng, gěi wǒ dǎ diànhuà.} « Si tu as du temps, appelle-moi. » (Oral, \zh{要是}, P2 impératif sans \zh{就}).
\item \textbf{Standard :} \zh{如果你不努力，你就考不上大学。} \emph{Rúguǒ nǐ bù nǔlì, nǐ jiù kǎo bù shàng dàxué.} « Si tu ne travailles pas, tu n'entreras pas à l'université. »
\item \textbf{Nécessaire (只有…才…) :} \zh{只有天天练习，你才能说好中文。} \emph{Zhǐyǒu tiāntiān liànxí, nǐ cái néng shuōhǎo Zhōngwén.} « Ce n'est qu'en t'entraînant chaque jour que tu pourras bien parler chinois. »
\item \textbf{Sauf si (除非…否则…) :} \zh{除非下大雪，否则我们照常上课。} \emph{Chúfēi xià dàxuě, fǒuzé wǒmen zhàocháng shàngkè.} « À moins qu'il ne neige abondamment, nous avons cours comme d'habitude. »
\item \textbf{Contrefactuel :} \zh{假如我早知道，我就不来了。} \emph{Jiǎrú wǒ zǎo zhīdào, wǒ jiù bù lái le.} « Si j'avais su plus tôt, je ne serais pas venu. »
\end{enumerate}
\end{exemplebox}

\courssub{Erreurs fréquentes : \zh{就} vs \zh{才}, omission de \zh{就}}

\begin{attention}[Pièges condition]
\begin{enumerate}
\item \textbf{Omettre \zh{就}.} \textit{Faux :} \zh{如果下雨，我回家。} \textit{Correct :} \zh{如果下雨，我就回家。} (Sans \zh{就}, le lien logique est faible/ambigu).
\item \textbf{Confondre \zh{就} et \zh{才}.} \zh{就} = conséquence directe, suffisante (si A, alors B). \zh{才} = condition nécessaire (seulement si A, alors B \rightarrow \zh{只有…才…}).
\item \textbf{Mauvaise place de \zh{就}.} Rappel : Sujet + \zh{就} + Verbe. Pas \zh{就} + Sujet + Verbe.
\item \textbf{\zh{要是} à l'écrit formel.} \zh{要is} est oral. À l'écrit (examen, rédaction), préférez \zh{如果} ou \zh{假如}.
\end{enumerate}
\end{attention}

\courssub{Exercices d'application : Condition}

\begin{exercice}[Condition : completion et transformation]
\begin{enumerate}
\item Complétez avec \zh{就} ou \zh{才} : a) \zh{只有认真复习，\_\_\_\_能考好。} b) \zh{如果天气好，我们\_\_\_\_去爬山。}
\item Traduisez : « À moins de finir les devoirs (做完作业 \emph{zuòwán zuòyè}), tu ne peux pas sortir (出去 \emph{chūqù}). » (Utilisez \zh{除非…否则…}).
\item Transformez en hypothétique (\zh{假如…就…}) : \zh{如果我有钱，我就买车。} \rightarrow \zh{假如我有钱，我就买车。} (Même structure, registre plus formel).
\item Corrigez la phrase : \zh{如果你去，就我去。}
\end{enumerate}
\end{exercice}

\begin{corrige}[Exercice Condition]
\begin{enumerate}
\item a) \zh{才} (structure \zh{只有…才…}). b) \zh{就} (structure \zh{如果…就…}).
\item \zh{除非你做完作业，否则你不能出去。} / \zh{除非做完作业，就不准出去。}
\item \zh{假如我有钱，我就买车。} (Registre plus soutenu).
\item \zh{如果你去，我就去。} (\zh{就} placé après le sujet \zh{我}).
\end{enumerate}
\end{corrige}

\coursec{But : 为了… et coordination simple}

\courssub{La structure de but : 为了 + Verbe/Objectif ， Sujet + Action}

\zh{为了} (wèile) est une préposition « pour, dans le but de ». Elle introduit une proposition de but (P1). La principale (P2) exprime l'action mise en œuvre pour atteindre ce but. \textbf{Il n'y a pas de second corrélat en P2} (pas de \zh{所以}, \zh{就}, etc.).

\begin{propriete}[Particularités de \zh{为了}]
\begin{itemize}
\item P1 = \zh{为了} + Verbe (souvent à sens volitif/actif) ou Nom d'action. Sujet de P1 souvent omis s'il est co-référent à P2.
\item P2 = Action concrète, souvent avec modalité de nécessité : \zh{要} (vouloir/falloir), \zh{得} (devoir), \zh{应该} (devoir moral), \zh{决定} (décider), \zh{开始} (commencer).
\item \textbf{Négation :} \zh{为了不…} (pour ne pas...). Ex : \zh{为了不迟到，我早早就出门了。}
\end{itemize}
\end{propriete}

\begin{exemplebox}[But : exemples]
\begin{enumerate}
\item \zh{为了买便宜的衣服，我们去了马路湾市场。}\\
\emph{Wèile mǎi piányi de yīfu, wǒmen qù le Mǎlùwān Shìchǎng.}\\
« Pour acheter des vêtements pas chers, nous sommes allés au marché Mallowan (Douala). »
\item \zh{为了不让父母担心，她每天发微信报平安。}\\
\emph{Wèile bù ràng fùmǔ dānxīn, tā měitiān fā Wēixìn bào píng'ān.}\\
« Pour ne pas inquiéter ses parents, elle envoie un WeChat chaque jour pour donner des nouvelles. » (\zh{报平安} : donner de ses nouvelles / signaler qu'on va bien).
\item \zh{为了通过HSK4，他报了一个网课。}\\
\emph{Wèile tōngguò HSK4, tā bào le yí ge wǎngkè.}\\
« Pour passer le HSK4, il s'est inscrit à un cours en ligne. »
\end{enumerate}
\end{exemplebox}

\courssub{Variantes formelles : 为 / 以便 / 以免}

\begin{center}
\begin{tabularx}{\linewidth}{|X|X|X|}\hline
\rowcolor{popblueL} Marqueur & Registre / Nuance & Exemple \\ \hline
\zh{为} (wèi) & Très formel, écrit, souvent \zh{为} + Nom (pas verbe seul) & \zh{为人民服务} (Servir le peuple) \\ \hline
\zh{以便} (yǐbiàn) & « Afin que, de manière à » — introduit une facilitation & \zh{请留下电话，以便联系。} (Laissez votre tel pour qu'on puisse vous contacter) \\ \hline
\zh{以免} (yǐmiǎn) & « De peur que, pour éviter que » — but négatif & \zh{带伞，以免淋雨。} (Prends un parapluie pour éviter d'être mouillé) \\ \hline
\end{tabularx}
\end{center}

\courssub{Coordination simple (rappel et contraste)}

Contrairement à la phrase complexe (subordination), la \textbf{coordination} relie deux éléments de même nature (noms, verbes, propositions) sans rapport de dépendance logique fort (cause/condition...).

\begin{propriete}[Coordinations nominales et verbales]
\begin{itemize}
\item \textbf{Noms :} \zh{和} (hé), \zh{跟} (gēn - oral), \zh{与} (yǔ - formel), \zh{及} (jí - formel, listes). \zh{我和他去雅温得。}
\item \textbf{Verbes / Adjectifs (simultané / addition) :} \zh{又…又…} (yòu… yòu… « à la fois... et... »), \zh{既…又…} (jì… yòu… - formel), \zh{一边…一边…} (yībiān… yībiān… « tout en... »).
\item \textbf{Propositions (parataxe) :} Juxtaposition ou \zh{而且} (érqiě « de plus »), \zh{并且} (bìngqiě « et en plus »).
\end{itemize}
\end{propriete}

\begin{attention}[Ne pas confondre coordination et subordination de but]
\begin{itemize}
\item \zh{我去市场买菜。} (Série verbale / but implicite : « Je vais au marché acheter des légumes »). Pas de \zh{为了}.
\item \zh{为了买菜，我去市场。} (But explicite, focus sur la finalité).
\item \zh{我去市场，而且买了菜。} (Coordination : deux actions successives/indépendantes).
\end{itemize}
\end{attention}

\courssub{Exercices d'application : But et Coordination}

\begin{exercice}[But et coordination]
\begin{enumerate}
\item Reliez par \zh{为了…} : \zh{我想考上大学。} / \zh{我每天复习到很晚。}
\item Choisissez : \zh{以便} ou \zh{以免} ? a) \zh{请慢点说，\_\_\_\_我听得懂。} b) \zh{系好安全带，\_\_\_\_发生意外。} (\zh{系好} \emph{jīhǎo} attacher, \zh{安全带} \emph{ānquándài} ceinture, \zh{意外} \emph{yìwài} accident).
\item Coordonnez les noms : \zh{雅温得} + \zh{杜阿拉} + \zh{巴马} (\zh{巴马} \emph{Bāmǎ} Bafoussam) avec le bon connecteur (écrit formel).
\item Traduisez : « Il étudie le chinois et l'anglais. » (Deux objets : coordination nominale).
\end{enumerate}
\end{exercice}

\begin{corrige}[Exercice But et Coordination]
\begin{enumerate}
\item \zh{为了考上大学，我每天复习到很晚。}
\item a) \zh{以便} (but positif : pour que je comprenne). b) \zh{以免} (but négatif : pour éviter l'accident).
\item \zh{雅温得、杜阿拉与巴马} (ou \zh{及}).
\item \zh{他学中文和英文。} / \zh{他学习中文与英文。}
\end{enumerate}
\end{corrige}

\coursec{Synthèse, pièges des francophones et entraînement global}

\courssub{Tableau récapitulatif des quatre structures majeures}

\begin{center}
\begin{tabularx}{\linewidth}{|X|X|X|X|X|}\hline
\rowcolor{popblueL} Rapport & Structure type & P1 (Subordonnée) & P2 (Principale) & Marqueur P2 obligatoire ? \\ \hline
Cause & \zh{因为…所以…} & \zh{因为} + Cause & \zh{所以} + Conséquence & \textbf{Oui} (\zh{所以}) \\ \hline
Concession & \zh{虽然…但是…} & \zh{虽然} + Fait admis & \zh{但是/可是/不过/却} + Fait principal & \textbf{Oui} (adv. contraste) \\ \hline
Condition & \zh{如果…就…} & \zh{如果/要是} + Condition & Sujet + \zh{就} + Conséquence & \textbf{Oui} (\zh{就} avant V) \\ \hline
But & \zh{为了…} & \zh{为了} + But & Action / Modalité & \textbf{Non} \\ \hline
\end{tabularx}
\end{center}

\courssub{Les 5 pièges majeurs des francophones (Check-list de révision)}

\begin{attention}[Top 5 des erreurs à l'examen]
\begin{enumerate}
\item \textbf{Le « donc » fantôme (Cause) :} Écrire \zh{因为下雨，我没去。} au lieu de \zh{因为下雨，所以我没去。} $\rightarrow$ \textbf{Toujours écrire le corrélat P2 à l'écrit.}
\item \textbf{Le « mais » fantôme (Concession) :} Écrire \zh{虽然很累，我很开心。} au lieu de \zh{虽然很累，但是我很开心。} $\rightarrow$ \textbf{P2 nécessite un marqueur de contraste.}
\item \textbf{Le \zh{just} égaré (Condition) :} Écrire \zh{如果下雨，就我不去。} $\rightarrow$ \textbf{\zh{就} se place APRÈS le sujet de P2, AVANT le verbe.}
\item \textbf{Le \zh{才} au lieu de \zh{就} (Condition) :} \zh{如果努力，才能成功。} $\rightarrow$ \textbf{\zh{才} n'est utilisé qu'avec \zh{只有}. Avec \zh{如果}, c'est \zh{就}.}
\item \textbf{Le double corrélat pour le But :} Écrire \zh{为了学中文，所以我努力。} $\rightarrow$ \textbf{\zh{为了} n'appelle PAS de corrélat en P2. Supprimez \zh{所以}.}
\end{enumerate}
\end{attention}

\courssub{Exercices de synthèse : Transformation et Traduction}

\begin{exercice}[Synthèse : du français au chinois (Thème)]
Traduisez les phrases suivantes en chinois (caractères + pinyin).
\begin{enumerate}
\item Parce qu'il y a la grève (罢工 \emph{bàgōng}) aujourd'hui, donc il n'y a pas de cours.
\item Bien que le trajet (路程 \emph{lùchéng}) soit long, cependant il insiste pour y aller.
\item Si tu ne viens pas demain, alors je ne viendrai pas non plus (\zh{也} \emph{yě} / \zh{都} \emph{dōu}).
\item Pour ne pas rater (错过 \emph{cuòguò}) le train (火车 \emph{huǒchē}), nous sommes partis très tôt (\zh{很早就} \emph{hěn zǎo jiù}).
\item Ce n'est pas parce qu'il est riche (有钱 \emph{yǒuqián}) qu'il est heureux, c'est parce qu'il a une famille unie (家庭和睦 \emph{jiātíng hémù}).
\end{enumerate}
\end{exercice}

\begin{exercice}[Synthèse : du chinois au français (Version)]
Traduisez en français naturel.
\begin{enumerate}
\item \zh{由于台风，因此航班取消了。} (\zh{台风} \emph{táifēng} typhon, \zh{航班} \emph{hángbān} vol)
\item \zh{即使你不帮我，我自己也能做完。} (\zh{自己} \emph{zìjǐ} moi-même, \zh{做完} \emph{zuòwán} finir de faire)
\item \zh{只有通过面试，才能得到这份工作。} (\zh{面试} \emph{miànshì} entretien, \zh{得到} \emph{dédao} obtenir)
\item \zh{为了庆祝国庆节，学校放假三天。} (\zh{庆祝} \emph{qìngzhù} célébrer, \zh{国庆节} \emph{Guóqìngjié} Fête nationale, \zh{放假} \emph{fàngjià} être en congé)
\item \zh{虽然他只学了半年，但是中文说得很流利。} (\zh{只} \emph{zhǐ} seulement, \zh{半年} \emph{bànnián} six mois, \zh{流利} \emph{liúlì} couramment)
\end{enumerate}
\end{exercice}

\begin{exercice}[Production écrite guidée : La réponse de Li Ming]
Amina a envoyé sa lettre. Li Ming lui répond depuis Shanghai. Rédigez sa réponse (60-80 caractères) en utilisant \textbf{au moins 3 types de phrases complexes} (cause, concession, condition, but). Contexte : Il pleut à Shanghai aussi ; il est content qu'elle étudie ; s'elle vient en Chine, il lui fera visiter ; pour progresser, il lui conseille de regarder des séries chinoises.
\begin{itemize}
\item Vocabulaire utile : \zh{上海} \emph{Shànghǎi}, \zh{电视剧} \emph{diànshìjù} (série TV), \zh{带你参观} \emph{dài nǐ cānguān} (te faire visiter), \zh{建议} \emph{jiànyì} (conseiller).
\end{itemize}
\end{exercice}

\begin{corrige}[Exercices de synthèse]
\textbf{Thème (Français $\rightarrow$ Chinois)}
\begin{enumerate}
\item \zh{因为今天罢工，所以没有课。} \emph{Yīnwèi jīntiān bàgōng, suǒyǐ méiyǒu kè.}
\item \zh{虽然路程很远，但是他坚持要去。} \emph{Suīrán lùchéng hěn yuǎn, dànshì tā jiānchí yào qù.}
\item \zh{如果你明天不来，我就也不来了。} \emph{Rúguǒ nǐ míngtiān bù lái, wǒ jiù yě bù lái le.}
\item \zh{为了不误火车，我们很早就出发了。} \emph{Wèile bù wù huǒchē, wǒmen hěn zǎo jiù chūfā le.}
\item \zh{他之所以快乐，不是因为有钱，而是因为家庭和睦。} \emph{Tā zhīsuǒyǐ kuàilè, bù shì yīnwèi yǒuqián, ér shì yīnwèi jiātíng hémù.}
\end{enumerate}

\textbf{Version (Chinois $\rightarrow$ Français)}
\begin{enumerate}
\item En raison du typhon, le vol a été annulé. (Registre formel \zh{由于…因此…})
\item Même si tu ne m'aides pas, je pourrai le finir tout seul. (Concession forte \zh{即使…也…})
\item Ce n'est qu'en passant l'entretien qu'on peut obtenir ce travail. (Condition nécessaire \zh{只有…才…})
\item Pour célébrer la Fête nationale, l'école est en congé pendant trois jours. (But \zh{为了…})
\item Bien qu'il n'ait étudié que six mois, il parle chinois très couramment. (Concession \zh{虽然…但是…} + \zh{说得很流利} complément de degré).
\end{enumerate}

\textbf{Production écrite (Modèle de réponse)}
\zh{亲爱的阿米娜：}\\
\zh{因为上海也在下雨，所以我今天在家看书。}\\
\zh{虽然我很忙，但是我很高兴你努力学中文。}\\
\zh{如果你将来来中国，我就带你参观外滩。}\\
\zh{为了进步快一点，我建议你看中文电视剧。}\\
\zh{祝好，李明}
\end{corrige}

\begin{aretenir}[Bilan final de la leçon 22 — Phrase complexe]
\begin{itemize}
\item \textbf{Quatre rapports logiques :} Cause (\zh{因为…所以…}), Concession (\zh{虽然…但是…}), Condition (\zh{如果…就…}), But (\zh{为了…}).
\item \textbf{Règle d'or :} Chinois = double marquage (sauf But) ; Français = simple marquage. \textbf{Ne jamais calquer le français.}
\item \textbf{Syntaxe rigide :} Virgule chinoise «，» obligatoire entre P1 et P2. \zh{就} placé après le sujet de P2. \zh{才} uniquement avec \zh{只有}.
\item \textbf{Registres :} \zh{由于…因此…} (formel), \zh{要is} (oral), \zh{假如/即使} (hypothétique), \zh{为/以便/以免} (écrit formel).
\item \textbf{Coordination vs Subordination :} \zh{和/与/跟} (noms), \zh{又…又…} (qualités), \zh{一边…一边…} (simultanéité) $\neq$ phrases complexes à corrélats.
\end{itemize}
\end{aretenir}

\coursec{Notion de phrase complexe}

En chinois, une \cle{phrase complexe} (复句, \emph{fùjù}) est une phrase qui contient au moins deux propositions (chacune avec son propre prédicat), reliées entre elles par des \cle{mots-outils corrélatifs} (关联词, \emph{guānliáncí}) — comme \zh{因为...所以...}, \zh{虽然...但是...}, \zh{如果...就...} — ou simplement par l'ordre logique et la ponctuation (parataxe). Contrairement au français qui utilise souvent des subordonnées introduites par des conjonctions uniques (\emph{parce que, si, bien que}), le chinois standard écrit privilégie la \textbf{corrélation binaire} : un mot ouvre la première proposition, un second mot ouvre la seconde, créant un cadre logique symétrique.

\begin{definition}[Types de relations logiques au programme]
Les phrases complexes du programme de Terminale couvrent principalement :
\begin{itemize}
\item \textbf{Cause -- Conséquence} : \zh{因为...所以...} (\emph{yīnwèi... suǒyǐ...})
\item \textbf{Concession -- Opposition} : \zh{虽然...但是/可是...} (\emph{suīrán... dànshì/kěshì...})
\item \textbf{Condition -- Hypothèse} : \zh{如果...就...} (\emph{rúguǒ... jiù...}) ; \zh{要是...就...} (\emph{yàoshi... jiù...})
\item \textbf{But} : \zh{为了...} (\emph{wèile...}) + verbe d'action
\item \textbf{Coordination simple} : \zh{不但...而且...} (\emph{búdàn... érqiě...}), \zh{一边...一边...} (\emph{yībiān... yībiān...})
\end{itemize}
Chaque structure impose un \textbf{ordre fixe} des propositions (cause avant conséquence, condition avant conséquence, concession avant fait principal) et une \textbf{place stricte du sujet} (unique ou double).
\end{definition}

\begin{aretenir}[Principes universels de la phrase complexe chinoise]
\begin{enumerate}
\item \textbf{Ordre logique immuable} : la cause/condition/concession précède toujours la conséquence/résultat/fait principal. Pas d'inversion possible comme en français (\emph{« Je reste car il pleut »} $\neq$ \zh{我待在家里因为下雨}).
\item \textbf{Corrélats binaires} : à l'écrit formel, les deux marqueurs sont présents. À l'oral, le second (souvent \zh{所以}, \zh{但是}, \zh{就}) peut être omis.
\item \textbf{Sujet unique} : s'il est identique dans les deux propositions, il n'apparaît qu'une seule fois (avant le premier corrélat ou entre les deux). On ne le répète jamais devant le second corrélat.
\item \textbf{Ponctuation} : une virgule chinoise ， sépare systématiquement les deux propositions.
\end{enumerate}
\end{aretenir}

\coursec{Cause et conséquence : 因为...所以...}

Nous étudions maintenant le couple de corrélats le plus fréquent : \zh{因为...所以...}, qui exprime la \cle{cause} et la \cle{conséquence}.

\courssub{Observation : lire d'abord, comprendre ensuite}

Lisons ce court texte avant toute règle. Repérez les mots \zh{因为} et \zh{所以}.

\begin{textebox}[Un lundi matin à Yaoundé]
今天早上下了大雨，路上都是水。\\
\emph{Jīntiān zǎoshang xià le dà yǔ, lù shang dōu shì shuǐ.}\\
« Ce matin, il a beaucoup plu et les routes étaient pleines d'eau. »\\[4pt]
因为路不好，所以公共汽车来得很晚。\\
\emph{Yīnwèi lù bù hǎo, suǒyǐ gōnggòng qìchē lái de hěn wǎn.}\\
« Comme la route était mauvaise, le bus est arrivé très tard. »\\[4pt]
我因为起晚了，所以没吃早饭。\\
\emph{Wǒ yīnwèi qǐ wǎn le, suǒyǐ méi chī zǎofàn.}\\
« Comme je me suis levé tard, je n'ai pas pris de petit-déjeuner. »\\[4pt]
因为考试快了，所以大家都在复习。\\
\emph{Yīnwèi kǎoshì kuài le, suǒyǐ dàjiā dōu zài fùxí.}\\
« Comme l'examen approche, tout le monde est en train de réviser. »
\end{textebox}

\textbf{Questions d'observation :}
\begin{enumerate}
\item Dans chaque phrase, quelle partie indique la raison ? Quelle partie indique le résultat ?
\item Quel mot se trouve toujours \emph{avant} la virgule ? Quel mot se trouve \emph{après} ?
\item Dans la troisième phrase, où est placé le sujet \zh{我} par rapport à \zh{因为} ?
\end{enumerate}

Vous avez remarqué : \zh{因为} ouvre la proposition de cause, \zh{所以} ouvre la proposition de conséquence, et les deux propositions gardent chacune leur ordre interne Sujet–Verbe–Objet. C'est exactement la règle que nous allons énoncer.

\courssub{La règle : schéma, forme, emploi}

\begin{propriete}[Structure de 因为...所以...]
\textbf{Schéma :} 因为 + proposition de cause ，(sujet) + 所以 + proposition de conséquence 。

\begin{itemize}
\item \zh{因为} (\emph{yīnwèi}, « parce que, comme ») introduit la \cle{cause} ;
\item \zh{所以} (\emph{suǒyǐ}, « donc, c'est pourquoi ») introduit la \cle{conséquence} ;
\item les deux propositions sont séparées par une \textbf{virgule chinoise} ， ;
\item l'ordre est fixe : la cause vient \textbf{d'abord}, la conséquence \textbf{ensuite}. On ne peut pas inverser comme en français (\emph{« Il est absent parce qu'il est malade »}).
\end{itemize}
\end{propriete}

\begin{popfigure}
\begin{tikzpicture}[node distance=0.9cm, font=\small]
\node[draw=popblue, fill=popblueL, rounded corners, align=center, minimum width=4.2cm, minimum height=1.1cm] (cause) {\zh{因为} + P1\\ \textbf{CAUSE}\\ \emph{la raison}};
\node[draw=poporange, fill=poporangeL, rounded corners, align=center, minimum width=4.2cm, minimum height=1.1cm, right=3.2cm of cause] (effet) {\zh{所以} + P2\\ \textbf{EFFET}\\ \emph{le résultat}};
\draw[-{Stealth[length=4mm]}, very thick, popdark] (cause) -- node[above, align=center]{entraîne\\ ，virgule} (effet);
\node[below=0.5cm of cause, align=center, text=popmuted] (ex1) {\zh{因为路不好，}\\ \emph{yīnwèi lù bù hǎo}};
\node[below=0.5cm of effet, align=center, text=popmuted] (ex2) {\zh{所以我们迟到了。}\\ \emph{suǒyǐ wǒmen chídào le}};
\draw[-{Stealth[length=3mm]}, dashed, popmuted] (ex1.east) -- (ex2.west);
\end{tikzpicture}
\legende{La chaîne causale : la proposition 因为 entraîne la proposition 所以.}
\end{popfigure}

\courssub{La place du sujet : trois configurations possibles}

C'est le point le plus délicat pour les francophones. Observez :

\begin{exemplebox}[Trois placements du sujet]
1. \textbf{Sujets différents} dans les deux propositions : chaque proposition garde son sujet.\\
\zh{因为天气很热，所以我们喝很多水。}\\
\emph{Yīnwèi tiānqì hěn rè, suǒyǐ wǒmen hē hěn duō shuǐ.}\\
« Comme il fait très chaud, nous buvons beaucoup d'eau. »\\[4pt]
2. \textbf{Sujet identique, placé avant 因为} (le plus naturel, thème de la phrase) : il ne se répète pas.\\
\zh{她因为生病，所以没来。}\\
\emph{Tā yīnwèi shēngbìng, suǒyǐ méi lái.}\\
« Elle n'est pas venue parce qu'elle était malade. »\\[4pt]
3. \textbf{Sujet identique, placé après 因为} (possible, met l'accent sur la cause) :\\
\zh{因为他学习努力，所以通过了考试。}\\
\emph{Yīnwèi tā xuéxí nǔlì, suǒyǐ tōngguò le kǎoshì.}\\
« Parce qu'il a travaillé dur, il a réussi l'examen. »
\end{exemplebox}

\begin{aretenir}[Règle du sujet unique]
Si le sujet est le même dans les deux propositions, il apparaît \textbf{une seule fois} : soit avant \zh{因为}, soit juste après \zh{因为}. On ne le répète jamais devant \zh{所以} dans ce cas. Si les sujets sont différents, chacun reste à sa place, à l'intérieur de sa proposition.
\end{aretenir}

\courssub{Variantes et exceptions}

\begin{propriete}[Formes allégées et usage oral]
\begin{itemize}
\item \textbf{À l'oral}, on omet souvent \zh{所以} quand la conséquence est évidente : \zh{因为下雨，我没去。} (\emph{Yīnwèi xià yǔ, wǒ méi qù.}) « Comme il pleuvait, je n'y suis pas allé. »
\item \textbf{所以 seul} peut conclure une réponse après une question \zh{为什么} :\\
— \zh{你为什么学汉语？} (\emph{Nǐ wèishénme xué Hànyǔ?})\\
— \zh{因为汉语很有意思。} (\emph{Yīnwèi Hànyǔ hěn yǒu yìsi.}) « Parce que le chinois est très intéressant. »\\
\zh{所以我要好好学习。} (\emph{Suǒyǐ wǒ yào hǎohǎo xuéxí.}) « Donc je dois bien travailler. »
\item \textbf{Exception importante} : dans une phrase complète \emph{écrite}, on peut omettre \zh{所以}, mais on n'omet \textbf{jamais} \zh{因为} en gardant \zh{所以} seul au début. \zh{所以...} tout seul n'ouvre pas une phrase de cause complète ; il suppose une cause énoncée juste avant.
\end{itemize}
\end{propriete}

\begin{savaistu}[Pourquoi deux mots là où le français n'en met qu'un ?]
En chinois classique, \zh{所以} signifiait « ce par quoi, la raison pour laquelle » (nom verbal). La langue moderne a gardé l'habitude d'encadrer la relation logique par \textbf{deux} marqueurs, comme un couple inséparable (\zh{因}... \zh{果} : cause -- effet). À l'oral rapide, les Chinois laissent souvent tomber \zh{所以} — mais à l'examen et dans vos rédactions, exigez-vous le couple complet \zh{因为...所以...} : c'est la forme correcte et valorisée.
\end{savaistu}

\courssub{Comparaison avec le français}

\begin{center}
\begin{tabularx}{\linewidth}{|X|X|X|}
\hline
\rowcolor{popblueL} \textbf{Français} & \textbf{Chinois} & \textbf{Remarque} \\
\hline
« parce que » seul & \zh{因为}...，\zh{所以}... & Le chinois utilise le couple complet à l'écrit. \\
\hline
« Comme il pleut, je reste. » & \zh{因为下雨，所以我待在家里。} & Même ordre : cause d'abord. \\
\hline
« Je reste parce qu'il pleut. » (conséquence d'abord) & \zh{我因为下雨，所以待在家里。} & Impossible d'inverser en chinois : la cause précède toujours. \\
\hline
« Parce qu'il pleut, \emph{donc} je reste. » (incorrect en français) & \zh{因为下雨，所以我待在家里。} & Ce que le français rejette, le chinois l'exige ! \\
\hline
\end{tabularx}
\end{center}

\begin{attention}[Erreur n$^\circ$ 1 des francophones]
En français, « parce que... donc... » est une faute de syntaxe (pléonasme). Beaucoup d'élèves traduisent donc « parce qu'il est malade, il est absent » par \zh{因为他病了，他没来} en supprimant \zh{所以} par réflexe français. \textbf{En chinois écrit, gardez les deux corrélats} : \zh{因为他病了，所以没来。} Inversement, ne mettez jamais \zh{所以} au début de la phrase comme le « donc » français : « Donc, il est absent » se dit \zh{所以他没来} seulement si la cause a déjà été dite dans la phrase précédente.
\end{attention}

\begin{methode}[Construire une phrase avec 因为...所以... en 4 étapes]
\begin{enumerate}
\item \textbf{Identifier} la cause et la conséquence dans la phrase française.
\item \textbf{Écrire la cause} en chinois, en respectant l'ordre SVO interne : \zh{路不好}.
\item \textbf{Écrire la conséquence} : \zh{我们迟到了}.
\item \textbf{Insérer les corrélats} sans rien changer à l'intérieur des propositions : \zh{因为} + cause + ，+ \zh{所以} + conséquence + 。 $\rightarrow$ \zh{因为路不好，所以我们迟到了。}
\end{enumerate}
Vérifiez enfin : sujet unique ou double ? Virgule présente ? Ordre cause $\rightarrow$ conséquence respecté ?
\end{methode}

\begin{exoresolu}[Traduction guidée]
Énoncé : traduisez en chinois « Comme il a beaucoup travaillé, Ali a réussi son examen de chinois. »
\tcblower
\textbf{Solution étape par étape :}
\begin{enumerate}
\item Cause : « il a beaucoup travaillé » $\rightarrow$ \zh{他学习很努力} ; conséquence : « Ali a réussi son examen » $\rightarrow$ \zh{Ali 通过了汉语考试}.
\item Sujets différents ? Non, c'est la même personne : on place le sujet une fois, avant \zh{因为} : \zh{Ali 因为...}.
\item Assemblage : \zh{Ali 因为学习很努力，所以通过了汉语考试。}\\
\emph{Ali yīnwèi xuéxí hěn nǔlì, suǒyǐ tōngguò le Hànyǔ kǎoshì.}
\end{enumerate}
\textbf{Vérification} : virgule entre les deux propositions, pas de répétition du sujet, ordre cause $\rightarrow$ conséquence. La phrase est correcte.
\end{exoresolu}

\begin{exercice}[À vous de jouer]
Traduisez en chinois avec \zh{因为...所以...} :
\begin{enumerate}
\item Comme la route est mauvaise, nous sommes arrivés en retard au lycée.
\item Elle n'a pas mangé parce qu'elle était malade.
\item Comme le match de football était très intéressant, tout le quartier l'a regardé.
\end{enumerate}
\end{exercice}

\begin{corrige}[Exercice « À vous de jouer »]
1. \zh{因为路不好，所以我们很晚才到高中。} \emph{Yīnwèi lù bù hǎo, suǒyǐ wǒmen hěn wǎn cái dào gāozhōng.}\\
2. \zh{她因为生病了，所以没有吃饭。} \emph{Tā yīnwèi shēngbìng le, suǒyǐ méiyǒu chīfàn.} (sujet unique, placé avant 因为)\\
3. \zh{因为足球比赛很有意思，所以全街区的人都看了。} \emph{Yīnwèi zúqiú bǐsài hěn yǒu yìsi, suǒyǐ quán jiēqū de rén dōu kàn le.}
\end{corrige}

\begin{definition}[Vocabulaire clé de la section]
\begin{center}
\begin{tabularx}{\linewidth}{|X|X|X|X|}
\hline
\rowcolor{popblueL} \textbf{Caractères} & \textbf{Pinyin} & \textbf{Nature} & \textbf{Traduction} \\
\hline
因为 & yīnwèi & conj. & parce que, comme \\
\hline
所以 & suǒyǐ & adv. & donc, c'est pourquoi \\
\hline
路 & lù & n. & route, chemin \\
\hline
不好 & bù hǎo & adj. & mauvais, pas bon \\
\hline
公共汽车 & gōnggòng qìchē & n. & bus (public) \\
\hline
来得晚 & lái de wǎn & v.+comp. & arriver tard \\
\hline
起晚 & qǐ wǎn & v. & se lever tard \\
\hline
早饭 & zǎofàn & n. & petit-déjeuner \\
\hline
生病 & shēngbìng & v. & tomber malade, être malade \\
\hline
没来 & méi lái & v. & ne pas venir (passé) \\
\hline
学习努力 & xuéxí nǔlì & v.+adv. & étudier avec assiduité \\
\hline
通过 & tōngguò & v. & passer (examen), traverser \\
\hline
考试 & kǎoshì & n./v. & examen / passer un examen \\
\hline
足球比赛 & zúqiú bǐsài & n. & match de football \\
\hline
有意思 & yǒu yìsi & adj. & intéressant \\
\hline
街区 & jiēqū & n. & quartier (rue, bloc) \\
\hline
\end{tabularx}
\end{center}
\end{definition}

\begin{aretenir}[L'essentiel de la section]
\begin{itemize}
\item Schéma : \zh{因为} + cause ，\zh{所以} + conséquence 。
\item À l'écrit, les deux corrélats sont exigés ; à l'oral, \zh{所以} peut tomber.
\item Sujet identique : une seule mention, avant ou après \zh{因为}.
\item Jamais d'inversion : la cause précède toujours la conséquence.
\item Piège français : ne supprimez pas \zh{所以} par habitude de « parce que » seul.
\end{itemize}
\end{aretenir}

\coursec{Concession : \zh{虽然}... \zh{但是}...}

Dans la section précédente, nous avons appris à relier une cause à sa conséquence avec \zh{因为...所以...} : les deux faits vont \emph{dans le même sens} (il pleut, donc je reste). Mais la vie n'est pas toujours aussi logique : parfois, deux faits \emph{s'opposent}. Il pleut, et pourtant je sors quand même. Le chinois exprime cette opposition avec une paire de mots-outils qui fonctionne exactement comme \zh{因为...所以...} : la paire \cle{虽然}... \cle{但是}..., équivalente du français « bien que... (mais) ». C'est la structure de la \textbf{concession}, l'une des plus fréquentes et des plus utiles du chinois.

\courssub{Observation : deux faits qui s'opposent}

\begin{textebox}[Dialogue : au lycée de Yaoundé]
甲：外面下雨了，你还去踢足球吗？\\
\emph{Jiǎ : Wàimiàn xià yǔ le, nǐ hái qù tī zúqiú ma ?}\\
« A : Il pleut dehors, tu vas quand même jouer au football ? »\\[4pt]
乙：虽然下雨了，但是我还要去。足球是我的生命！\\
\emph{Yǐ : Suīrán xià yǔ le, dànshì wǒ hái yào qù. Zúqiú shì wǒ de shēngmìng !}\\
« B : Bien qu'il pleuve, je vais quand même y aller. Le football, c'est ma vie ! »\\[4pt]
甲：你真爱足球。虽然天气很冷，可是你不怕。\\
\emph{Jiǎ : Nǐ zhēn ài zúqiú. Suīrán tiānqì hěn lěng, kěshì nǐ bú pà.}\\
« A : Tu aimes vraiment le football. Bien qu'il fasse froid, tu n'as pas peur. »\\[4pt]
乙：我虽然很累，但是每天都很快乐。\\
\emph{Yǐ : Wǒ suīrán hěn lèi, dànshì měitiān dōu hěn kuàilè.}\\
« B : Bien que je sois fatigué, je suis heureux tous les jours. »
\end{textebox}

Observons les phrases soulignées par la structure. Dans chaque cas, il y a \textbf{deux faits contradictoires} : la pluie devrait empêcher de jouer, la fatigue devrait rendre triste... et pourtant le second fait a lieu quand même. Remarquez aussi que \zh{虽然} n'apparaît \emph{jamais seul} : il est toujours suivi, dans la seconde proposition, de \zh{但是} ou de \zh{可是}.

\courssub{La règle : schéma, emploi, place du sujet}

\begin{propriete}[La structure de concession : \zh{虽然} + P1 ，\zh{但是} + P2]
\textbf{Schéma :} \zh{虽然} + fait concédé (P1) ，\zh{但是} / \zh{可是} + fait principal (P2)。
\begin{itemize}
\item \zh{虽然} (suīrán) introduit le fait que l'on \emph{concède}, que l'on reconnaît comme vrai ; il se place au début de la première proposition.
\item \zh{但是} (dànshì) ou \zh{可是} (kěshì) introduit le fait \emph{principal}, celui qui compte vraiment pour le locuteur ; il se place au début de la seconde proposition.
\item \textbf{Règle d'or :} l'information importante se trouve dans P2. La proposition P2 ne se laisse \textbf{jamais supprimer} : c'est elle qui porte le message.
\end{itemize}
\end{propriete}

\begin{popfigure}
\begin{tikzpicture}[scale=1]
% Pivot et poutre
\fill[popdark] (0,0) -- (-0.25,-1.1) -- (0.25,-1.1) -- cycle;
\draw[line width=2pt,popdark] (-3.4,0.35) -- (3.4,-0.35);
\fill[popdark] (0,0) circle (0.09);
% Plateau gauche (虽然) - plus haut
\draw[popblue,line width=1.2pt] (-3.4,0.35) -- (-4.1,-0.55) (-3.4,0.35) -- (-2.7,-0.55);
\draw[popblue,fill=popblueL,line width=1.2pt] (-4.1,-0.55) arc (180:360:0.7);
\node[align=center,popblue] at (-3.4,-1.55) {\zh{虽然} + fait concédé\\ \small il pleut / je suis fatigué};
% Plateau droit (但是) - plus bas = plus lourd
\draw[poporange,line width=1.2pt] (3.4,-0.35) -- (2.7,-1.35) (3.4,-0.35) -- (4.1,-1.35);
\draw[poporange,fill=poporangeL,line width=1.2pt] (2.7,-1.35) arc (180:360:0.7);
\node[align=center,poporange] at (3.4,-2.35) {\zh{但是} + fait principal\\ \small \textbf{plateau plus lourd : l'idée principale}};
\end{tikzpicture}
\legende{La balance de la concession : le plateau \zh{但是} pèse plus lourd, car c'est lui qui porte l'idée principale.}
\end{popfigure}

\begin{exemplebox}[Premiers exemples]
\begin{itemize}
\item \zh{虽然中文很难，但是很有意思。}\\
\emph{Suīrán Zhōngwén hěn nán, dànshì hěn yǒu yìsi.}\\
« Bien que le chinois soit difficile, il est très intéressant. »
\item \zh{虽然他家很远，但是他从不迟到。}\\
\emph{Suīrán tā jiā hěn yuǎn, dànshì tā cóng bù chídào.}\\
« Bien que sa maison soit loin, il n'est jamais en retard. »
\item \zh{虽然我很累，但是我要完成作业。}\\
\emph{Suīrán wǒ hěn lèi, dànshì wǒ yào wánchéng zuòyè.}\\
« Bien que je sois fatigué, je dois finir mes devoirs. »
\item \zh{虽然他是外国人，但是中文说得很好。}\\
\emph{Suīrán tā shì wàiguórén, dànshì Zhōngwén shuō de hěn hǎo.}\\
« Bien qu'il soit étranger, il parle très bien chinois. »
\end{itemize}
\end{exemplebox}

\courssub{La place du sujet : identique ou différent}

La place du sujet par rapport à \zh{虽然} dépend d'une question simple : les deux propositions ont-elles le \emph{même} sujet ?

\begin{propriete}[Place du sujet]
\begin{itemize}
\item \textbf{Sujets différents :} chaque proposition a son sujet, et \zh{虽然} se place \emph{avant} le sujet de P1.\\
\zh{虽然天气很冷，但是我们还出去。}\\
\emph{Suīrán tiānqì hěn lěng, dànshì wǒmen hái chūqù.}\\
« Bien qu'il fasse froid, nous sortons quand même. » (sujets : \zh{天气} / \zh{我们})
\item \textbf{Sujet identique :} le sujet se place \emph{avant} \zh{虽然} et n'est pas répété dans P2 (ou remplacé par un pronom).\\
\zh{他虽然很忙，但是帮助我。}\\
\emph{Tā suīrán hěn máng, dànshì bāngzhù wǒ.}\\
« Bien qu'il soit très occupé, il m'aide. » (sujet unique : \zh{他})
\end{itemize}
\end{propriete}

\begin{attention}[Sujet identique : ne répétez pas mécaniquement]
Quand le sujet est le même dans les deux propositions, le placer après \zh{虽然} et le répéter ensuite (\zh{虽然我很累，但是我很快乐}) n'est pas une faute grave, mais c'est lourd. La formulation naturelle est : \zh{我虽然很累，但是很快乐。} Retenez : \textbf{sujet unique = sujet devant \zh{虽然}}.
\end{attention}

\courssub{Les variantes de registre : \zh{但是}, \zh{可是}, \zh{却}}

Le chinois dispose de plusieurs mots pour marquer l'opposition dans P2. Ils ne sont pas interchangeables à 100 \% : ils diffèrent par le \textbf{registre} et par la \textbf{place dans la phrase}.

\begin{center}
\begin{tabularx}{\linewidth}{|X|X|X|X|}
\hline
\rowcolor{popblueL} Mot & Pinyin & Registre / place & Traduction \\
\hline
\zh{但是} & dànshì & neutre, écrit et oral, début de P2 & mais, pourtant \\
\hline
\zh{可是} & kěshì & oral, courant, début de P2 & mais, pourtant \\
\hline
\zh{却} & què & soutenu, \emph{adverbe} : après le sujet, avant le verbe & et pourtant \\
\hline
\zh{虽然} & suīrán & conjonction de concession, début de P1 & bien que \\
\hline
\end{tabularx}
\end{center}

\begin{exemplebox}[Les trois variantes en action]
\begin{itemize}
\item \zh{虽然这道菜很贵，但是很好吃。} \emph{Suīrán zhè dào cài hěn guì, dànshì hěn hǎochī.} « Bien que ce plat soit cher, il est très bon. » (registre neutre)
\item \zh{虽然他很年轻，可是很有经验。} \emph{Suīrán tā hěn niánqīng, kěshì hěn yǒu jīngyàn.} « Bien qu'il soit jeune, il a beaucoup d'expérience. » (oral)
\item \zh{虽然他很小，却很有力气。} \emph{Suīrán tā hěn xiǎo, què hěn yǒu lìqi.} « Bien qu'il soit petit, il est pourtant très fort. » (soutenu, \zh{却} après le sujet, avant le verbe)
\end{itemize}
\end{exemplebox}

\begin{savaistu}[\zh{却} : l'adverbe qui se glisse avant le verbe]
Contrairement à \zh{但是} et \zh{可是}, qui sont des \emph{conjonctions} placées en tête de proposition, \zh{却} est un \emph{adverbe} : il obéit à la règle d'or des adverbes chinois, c'est-à-dire qu'il se place \textbf{après le sujet et avant le verbe}. On peut même le combiner avec \zh{但是} : \zh{他虽然病了，但是却还来上课。} \emph{Tā suīrán bìng le, dànshì què hái lái shàngkè.} « Bien que malade, il est pourtant venu en cours. » Ce double marquage renforce l'opposition et est très apprécié à l'écrit.
\end{savaistu}

\courssub{Comparaison avec le français}

\begin{center}
\begin{tabularx}{\linewidth}{|X|X|X|}
\hline
\rowcolor{popblueL} Aspect & Français & Chinois \\
\hline
Nombre de mots-outils & un seul : « bien que » (le « mais » est interdit après) & \textbf{deux} : \zh{虽然}... \zh{但是}... obligatoires ensemble \\
\hline
Mode du verbe & subjonctif possible (« bien qu'il soit ») & pas de mode : le verbe reste inchangé \\
\hline
Idée principale & dans la proposition principale & dans P2, après \zh{但是} \\
\hline
Registre oral & « il pleut mais je sors » (sans « bien que ») & \zh{虽然}... \zh{可是}... très naturel à l'oral \\
\hline
\end{tabularx}
\end{center}

En français, on dit « bien qu'il pleuve, je sors » ou « il pleut, \emph{mais} je sors » : jamais les deux ensemble (« bien qu'il pleuve mais je sors » est incorrect). En chinois, c'est l'inverse : la \textbf{double marque} \zh{虽然}... \zh{但是}... est la norme, et c'est la suppression de la seconde marque qui est fautive. Bonne nouvelle : le chinois ne connaît ni subjonctif ni conjugaison, donc aucune forme verbale à modifier.

\begin{attention}[L'erreur n° 1 des francophones : \zh{虽然} tout seul]
Ne traduisez jamais « bien que » par \zh{虽然} seul : \zh{虽然中文很难。} sonne \emph{inachevée} pour un sinophone, comme si vous vous arrêtiez au milieu de votre phrase. Sans \zh{但是} ou \zh{可是}, l'opposition n'est pas marquée et l'auditeur attend la suite. Retenez le réflexe : \textbf{\zh{虽然} appelle toujours \zh{但是}}.
\end{attention}

\begin{methode}[Transformer une phrase française en phrase de concession]
Prenons : « Il pleut mais je sors. »
\begin{enumerate}
\item \textbf{Identifier l'opposition} : deux faits contradictoires (la pluie devrait m'empêcher de sortir, pourtant je sors).
\item \textbf{Repérer le fait principal} : « je sors » — c'est lui qui ira dans P2.
\item \textbf{Traduire P1} avec \zh{虽然} : \zh{虽然下雨} (suīrán xià yǔ).
\item \textbf{Traduire P2} avec \zh{但是} : \zh{但是我出去} (dànshì wǒ chūqù).
\item \textbf{Assembler} : \zh{虽然下雨，但是我出去。} « Bien qu'il pleuve, je sors. »
\end{enumerate}
Vérifiez toujours : les deux mots-outils sont-ils présents ? Le fait principal est-il bien après \zh{但是} ?
\end{methode}

\begin{exoresolu}[Traduction guidée]
Énoncé : traduisez en chinois « Bien que le ndolé soit un peu cher, je l'aime beaucoup. »
\tcblower
Solution détaillée :
\begin{enumerate}
\item Opposition : le prix (fait concédé) contre l'amour du plat (fait principal).
\item P1 : \zh{虽然 ndolé 有点儿贵} — \emph{Suīrán ndolé yǒudiǎnr guì}.
\item P2 (idée principale) : \zh{但是我很喜欢} — \emph{dànshì wǒ hěn xǐhuan}.
\item Phrase complète : \zh{虽然 ndolé 有点儿贵，但是我很喜欢。}\\
\emph{Suīrán ndolé yǒudiǎnr guì, dànshì wǒ hěn xǐhuan.}\\
« Bien que le ndolé soit un peu cher, je l'aime beaucoup. »
\end{enumerate}
Vérification : \zh{虽然} et \zh{但是} sont tous les deux présents ; le fait principal (« je l'aime ») est bien dans P2. La phrase est correcte et naturelle.
\end{exoresolu}

\begin{exercice}[À vous de jouer]
\begin{enumerate}
\item Traduisez : « Bien que Douala soit loin de Yaoundé, il y va chaque mois. »
\item Traduisez : « Bien que ce caractère soit difficile, je peux l'écrire. »
\item Complétez avec \zh{但是}, \zh{可是} ou \zh{却} : \zh{虽然今天很热，} \ldots\ldots \zh{我们还要上课。}
\item Corrigez la phrase fautive : \zh{虽然他病了，还来学校。}
\end{enumerate}
\end{exercice}

\begin{corrige}[Exercice « À vous de jouer »]
\begin{enumerate}
\item \zh{虽然 Douala 离 Yaoundé 很远，但是他每个月都去。} \emph{Suīrán Douala lí Yaoundé hěn yuǎn, dànshì tā měi ge yuè dōu qù.}
\item \zh{虽然这个汉字很难，但是我会写。} \emph{Suīrán zhè ge Hànzì hěn nán, dànshì wǒ huì xiě.}
\item \zh{虽然今天很热，但是（可是）我们还要上课。} Les deux conjonctions sont possibles ; \zh{却} est impossible ici car il se placerait après le sujet : \zh{我们却还要上课。}
\item Il manque la marque d'opposition dans P2 : \zh{虽然他病了，但是还来学校。} \emph{Suīrán tā bìng le, dànshì hái lái xuéxiào.}
\end{enumerate}
\end{corrige}

\begin{aretenir}[L'essentiel de la concession]
\begin{itemize}
\item Schéma : \zh{虽然} + fait concédé ，\zh{但是} / \zh{可是} + fait principal。
\item L'idée principale est toujours dans la \textbf{seconde} proposition.
\item \zh{虽然} ne s'emploie \textbf{jamais seul} : il appelle obligatoirement \zh{但是} ou \zh{可是}.
\item Sujet identique : le sujet se place \emph{avant} \zh{虽然} ; sujets différents : \zh{虽然} en tête de P1.
\item \zh{却} est un adverbe soutenu : après le sujet, avant le verbe.
\item Contrairement au français, le chinois garde \textbf{deux} mots-outils, et le verbe ne change jamais de forme.
\end{itemize}
\end{aretenir}

\coursec{Condition et hypothèse : 如果...就... / 要是...就...}

Dans la section précédente, nous avons appris à exprimer la \textbf{concession} avec \zh{虽然...但是...} : deux faits réels qui semblent s'opposer (\zh{虽然很累，但是我继续学习。}). Nous passons maintenant à une autre grande relation logique : la \textbf{condition}. Ici, le résultat n'est pas encore un fait : il \emph{dépend} d'une hypothèse. En français : « \emph{Si} tu travailles, tu réussiras. » En chinois : \zh{如果你努力学习，你就会成功。} C'est l'une des structures les plus fréquentes de l'examen, à l'écrit comme à l'oral.

\courssub{Observation : découvrir la structure}

Lisons d'abord ce petit dialogue entre deux élèves d'un lycée de Yaoundé, la veille d'un match de football.

\begin{textebox}[Dialogue : le match de demain]
甲：明天下午有足球比赛，你想去看吗？\\
\emph{Jiǎ : Míngtiān xiàwǔ yǒu zúqiú bǐsài, nǐ xiǎng qù kàn ma ?}\\
« A : Demain après-midi, il y a un match de football, tu veux aller le voir ? »\\[4pt]
乙：如果明天下雨，我就不去；如果天气好，我就去。\\
\emph{Yǐ : Rúguǒ míngtiān xiàyǔ, wǒ jiù bú qù ; rúguǒ tiānqì hǎo, wǒ jiù qù.}\\
« B : S'il pleut demain, je n'irai pas ; s'il fait beau, j'irai. »\\[4pt]
甲：要是你去的话，我就跟你一起去！\\
\emph{Jiǎ : Yàoshi nǐ qù dehua, wǒ jiù gēn nǐ yìqǐ qù !}\\
« A : Si tu y vas, j'irai avec toi ! »
\end{textebox}

\textbf{Observons.} Dans les réponses de B et de A, chaque phrase se coupe en \emph{deux propositions} séparées par une virgule :
\begin{itemize}
\item la première proposition commence par \zh{如果} ou \zh{要是} : c'est la \textbf{condition} (l'hypothèse) ;
\item la seconde contient \zh{就} placé juste avant le verbe : c'est le \textbf{résultat}.
\end{itemize}
On remarque aussi que le chinois n'emploie \emph{aucun temps particulier} : \zh{下雨} et \zh{去} restent à la forme simple, alors que le français dit « s'il \emph{pleut}, je n'\emph{irai} pas » (présent + futur). C'est le mot \zh{就} qui, à lui seul, assure le lien logique de conséquence.

\courssub{La règle : schéma, emploi, place de 就}

\begin{propriete}[Structure de la phrase conditionnelle]
\textbf{Schéma :} \zh{如果 / 要是} + condition \zh{，} (sujet) + \zh{就} + verbe (résultat)\zh{。}
\begin{itemize}
\item \zh{如果} (rúguǒ) introduit la condition ; \zh{要是} (yàoshi) est sa variante plus orale ; \zh{假如} (jiǎrú) est la variante soutenue, écrite.
\item \zh{就} (jiù) est \textbf{quasi obligatoire} dans la langue standard écrite : il marque le lien logique avec le résultat. Sans \zh{就}, la phrase sonne inachevée (à l'oral très familier, on l'entend parfois omis, mais à l'examen on l'écrit toujours).
\item \textbf{Place de} \zh{就} : toujours \emph{après le sujet} de la seconde proposition et \emph{avant le verbe} : \zh{如果明天下雨，我就不去。}
\item La proposition de condition se place en général \textbf{en premier}, suivie d'une virgule chinoise \zh{，}.
\item L'hypothèse exprimée est une hypothèse \textbf{réalisable} (il peut vraiment pleuvoir, tu peux vraiment y aller).
\end{itemize}
\end{propriete}

\begin{popfigure}
\begin{tikzpicture}[
node distance=0.9cm,
box/.style={draw=popblue, fill=popblueL, rounded corners, align=center, minimum width=3.4cm, minimum height=1cm, font=\small},
decision/.style={draw=poporange, fill=poporangeL, diamond, aspect=2.2, align=center, font=\small, inner sep=1pt},
res/.style={draw=popgreen, fill=popgreenL, rounded corners, align=center, minimum width=3.6cm, minimum height=1cm, font=\small},
arr/.style={-{Stealth[length=3mm]}, thick, popdark}]
\node[decision, align=center] (cond){Condition ?\\ \zh{如果} P1};
\node[res, below left=1.1cm and -0.4cm of cond, align=center] (oui){Résultat\\ sujet + \zh{就} + verbe};
\node[box, below right=1.1cm and -0.4cm of cond, align=center] (non){Pas de résultat\\ (la condition n'est pas remplie)};
\draw[arr] (cond) -- node[left, font=\small, text=popgreen]{oui} (oui);
\draw[arr] (cond) -- node[right, font=\small, text=poporange]{non} (non);
\end{tikzpicture}
\legende{Organigramme de la phrase conditionnelle : si la condition \zh{如果} P1 est remplie, le résultat (sujet + \zh{就} + verbe) suit.}
\end{popfigure}

\begin{exemplebox}[Exemples modèles]
\zh{如果你努力学习，你就会成功。}\\
\emph{Rúguǒ nǐ nǔlì xuéxí, nǐ jiù huì chénggōng.}\\
« Si tu travailles dur, tu réussiras. »\\[6pt]
\zh{要是明天晴天，我们就去爬山。}\\
\emph{Yàoshi míngtiān qíngtiān, wǒmen jiù qù páshān.}\\
« S'il fait beau demain, nous irons grimper la montagne. »\\[6pt]
\zh{如果你去杜阿拉，我就跟你一起去。}\\
\emph{Rúguǒ nǐ qù Dùālā, wǒ jiù gēn nǐ yìqǐ qù.}\\
« Si tu vas à Douala, j'irai avec toi. »\\[6pt]
\zh{要是有时间，我们就去看比赛。}\\
\emph{Yàoshi yǒu shíjiān, wǒmen jiù qù kàn bǐsài.}\\
« Si nous avons le temps, nous irons voir le match. »
\end{exemplebox}

\courssub{Les trois variantes de « si » : 如果, 要是, 假如}

Les trois mots ont exactement la même fonction et le même schéma ; seul le \textbf{registre de langue} change.

\begin{center}
\begin{tabularx}{\linewidth}{|X|X|X|X|}
\hline
\rowcolor{popblueL} Caractères & Pinyin & Registre & Traduction \\
\hline
如果 & rúguǒ & standard, oral et écrit & si \\
\hline
要是 & yàoshi & oral, familier & si \\
\hline
假如 & jiǎrú & écrit, soutenu & si, supposons que \\
\hline
就 & jiù & marque de conséquence & alors, donc \\
\hline
的话 & dehua & renforce la condition (fin de proposition) & (si ...) \\
\hline
\end{tabularx}
\end{center}

\begin{exemplebox}[La même phrase aux trois registres]
\zh{如果明天下雨，我就不去。} \emph{Rúguǒ míngtiān xiàyǔ, wǒ jiù bú qù.} « S'il pleut demain, je n'irai pas. » (standard)\\[4pt]
\zh{要是明天下雨，我就不去。} \emph{Yàoshi míngtiān xiàyǔ, wǒ jiù bú qù.} (oral, à la récréation)\\[4pt]
\zh{假如明天下雨，我就不去。} \emph{Jiǎrú míngtiān xiàyǔ, wǒ jiù bú qù.} (rédaction, lettre formelle)
\end{exemplebox}

\begin{savaistu}[Le renfort 的话]
On peut renforcer la condition en plaçant \zh{的话} (dehua) \emph{à la fin} de la première proposition : \zh{如果你去的话，我就去。} \emph{Rúguǒ nǐ qù dehua, wǒ jiù qù.} « Si tu y vas, j'irai. » \zh{的话} peut même remplacer \zh{如果} : \zh{你去的话，我就去。} Cette tournure, très fréquente à l'oral, est un excellent point bonus à l'examen : elle montre que tu maîtrises la langue authentique.
\end{savaistu}

\courssub{La négation dans la phrase conditionnelle}

La négation peut porter sur la condition, sur le résultat, ou sur les deux. On utilise \zh{不} pour les faits généraux et les volontés, \zh{没(有)} pour la possession ou l'accompli.

\begin{exemplebox}[Négations]
\zh{如果不努力，就不会进步。}\\
\emph{Rúguǒ bù nǔlì, jiù bú huì jìnbù.}\\
« Si l'on ne fait pas d'efforts, on ne progressera pas. » (double négation : \zh{不...就不...})\\[6pt]
\zh{要是没有钱，我就不买这本书。}\\
\emph{Yàoshi méiyǒu qián, wǒ jiù bù mǎi zhè běn shū.}\\
« Si je n'ai pas d'argent, je n'achèterai pas ce livre. »\\[6pt]
\zh{如果你不舒服，就不要去上课。}\\
\emph{Rúguǒ nǐ bù shūfu, jiù bú yào qù shàngkè.}\\
« Si tu ne te sens pas bien, ne va pas en cours. »
\end{exemplebox}

Remarque : quand le sujet des deux propositions est le même, on peut ne pas le répéter dans la seconde : \zh{如果不努力，就不会进步。} Mais dès que les sujets diffèrent, chacun garde le sien : \zh{如果你去，\textbf{我}就去。}

\courssub{Comparaison avec le français}

\begin{center}
\begin{tabularx}{\linewidth}{|X|X|X|}
\hline
\rowcolor{popblueL} Point comparé & Français & Chinois \\
\hline
Mot d'introduction & « si » + présent & \zh{如果 / 要是} + forme simple du verbe \\
\hline
Temps du résultat & futur (« tu réussir\emph{as} ») & aucun temps ; \zh{就} + verbe simple \\
\hline
Marque de conséquence & aucune (l'ordre suffit) & \zh{就} quasi obligatoire avant le verbe du résultat \\
\hline
« Si... si... » (deux hypothèses) & lourd à enchaîner & naturel : \zh{如果...；如果...} \\
\hline
\end{tabularx}
\end{center}

Le point essentiel : le chinois \textbf{n'a pas de conjugaison}. Le français encode l'hypothèse par ses temps (présent + futur) ; le chinois l'encode par ses \emph{mots-outils} : \zh{如果} ouvre l'hypothèse, \zh{就} ferme le raisonnement. C'est pourquoi \zh{就} ne doit jamais être oublié ni déplacé.

\begin{attention}[Erreurs fréquentes des francophones]
\begin{enumerate}
\item \textbf{Placer \zh{就} avant le sujet.} Faux : *\zh{如果下雨，就不我去。} Juste : \zh{如果下雨，我就不去。} \zh{就} \emph{suit} le sujet et \emph{précède} le verbe, comme un adverbe.
\item \textbf{Oublier \zh{就}.} *\zh{如果你努力，你会成功。} est compréhensible mais incomplet en langue standard écrite ; à l'examen, écris toujours \zh{你就会成功}.
\item \textbf{Mettre \zh{了} dans la condition.} *\zh{如果明天下雨了，...} : l'hypothèse porte sur un fait non réalisé, donc pas de \zh{了} d'accompli. Juste : \zh{如果明天下雨，我就不去。}
\item \textbf{Mal combiner \zh{那么} et \zh{就}.} Le français « si..., \emph{alors}... » se rend parfois par \zh{那么} : on peut dire \zh{如果你去，那么我就去。} (\zh{那么} se place \emph{avant le sujet} du résultat, et \zh{就} reste avant le verbe). Mais on ne met jamais \zh{那么} \emph{à la place} de \zh{就} : *\zh{如果你去，那么我去。} est incomplet.
\end{enumerate}
\end{attention}

\begin{methode}[Construire une phrase conditionnelle en trois étapes]
\begin{enumerate}
\item \textbf{Condition :} écris la première proposition en la faisant précéder de \zh{如果} (ou \zh{要是} à l'oral) : \zh{如果你去杜阿拉}...
\item \textbf{Virgule :} place la virgule chinoise \zh{，} entre les deux propositions.
\item \textbf{Résultat :} écris sujet + \zh{就} + verbe : \zh{我就跟你一起去。}
\end{enumerate}
Vérification finale : le sujet est-il bien \emph{avant} \zh{就} ? Y a-t-il un \zh{了} parasite dans la condition ? Si oui, corrige.
\end{methode}

\begin{exoresolu}[Transformation et traduction]
Énoncé. (a) Mets à la forme conditionnelle : \zh{明天下雨。我不去。} (b) Traduis en chinois : « Si tu as le temps, nous irons voir le match ensemble. »
\tcblower
Solution détaillée. (a) La première phrase devient la condition : on ajoute \zh{如果} devant ; la seconde devient le résultat : on place \zh{就} après le sujet \zh{我} et avant le verbe \zh{不去}. On obtient : \zh{如果明天下雨，我就不去。} \emph{Rúguǒ míngtiān xiàyǔ, wǒ jiù bú qù.} « S'il pleut demain, je n'irai pas. »\\[4pt]
(b) Étape 1, condition : « si tu as le temps » → \zh{要是你有时间} (registre oral, \zh{要是}). Étape 2 : virgule \zh{，}. Étape 3, résultat : sujet \zh{我们} + \zh{就} + verbe \zh{一起去看比赛}. Phrase complète : \zh{要是你有时间，我们就一起去看比赛。} \emph{Yàoshi nǐ yǒu shíjiān, wǒmen jiù yìqǐ qù kàn bǐsài.} Vérification : \zh{就} est bien après \zh{我们}, pas de \zh{了} dans la condition.
\end{exoresolu}

\begin{exercice}[À toi de jouer]
\begin{enumerate}
\item Traduis en français : \zh{如果你每天练习汉字，你就会写得很好。}
\item Complète avec \zh{就} à la bonne place : \zh{要是天气好，我们 \_\_\_\_ 去爬山。}
\item Corrige la faute : *\zh{如果他不来了，我就不等他。}
\item Traduis en chinois (registre oral) : « S'il pleut demain, nous resterons à la maison. »
\end{enumerate}
\end{exercice}

\begin{corrige}[Exercice « À toi de jouer »]
\begin{enumerate}
\item « Si tu t'entraînes aux caractères tous les jours, tu écriras très bien. »
\item \zh{要是天气好，我们就去爬山。} (\zh{就} après le sujet \zh{我们}, avant le verbe \zh{去}.)
\item La condition est une hypothèse non réalisée : pas de \zh{了}. Correction : \zh{如果他不来，我就不等他。} « S'il ne vient pas, je ne l'attendrai pas. »
\item \zh{要是明天下雨，我们就待在家里。} \emph{Yàoshi míngtiān xiàyǔ, wǒmen jiù dāi zài jiā lǐ.}
\end{enumerate}
\end{corrige}

\begin{aretenir}[L'essentiel de la section]
\begin{itemize}
\item Schéma : \zh{如果 / 要是 / 假如} + condition \zh{，} sujet + \zh{就} + verbe.
\item \zh{就} est quasi obligatoire à l'écrit et se place \emph{après le sujet, avant le verbe}.
\item Pas de temps en chinois : pas de \zh{了} dans la condition.
\item Renfort possible : \zh{的话} en fin de condition (\zh{如果你去的话，我就去。}) ; \zh{那么} possible avant le sujet du résultat, mais jamais à la place de \zh{就}.
\item Double négation possible : \zh{如果不努力，就不会进步。}
\end{itemize}
\end{aretenir}

\coursec{But : 为了... et coordination simple}

Dans la section précédente, nous avons appris à exprimer la condition et l'hypothèse avec \zh{如果...就...} et \zh{要是...就...} : « si tu travailles, tu réussiras ». Mais pourquoi travaille-t-on ? \emph{Dans quel but} ? C'est précisément ce que nous allons dire maintenant : exprimer le \cle{but} avec \zh{为了} (wèile, « pour, afin de »), puis coordonner deux qualités ou deux actions avec \zh{又...又...} et \zh{一边...一边...}. Ces trois structures sont très fréquentes à l'écrit comme à l'oral, et elles figurent régulièrement aux épreuves de traduction.

\courssub{Observation : lire d'abord, comprendre ensuite}

\begin{textebox}[Dialogue : les projets de Ngono]
A : 恩戈诺，你为什么每天学习汉语？\\
\emph{Ēngēnuò, nǐ wèishénme měitiān xuéxí Hànyǔ ?}\\
« Ngono, pourquoi étudies-tu le chinois tous les jours ? »\\
B : 为了通过考试，也为了以后去中国留学，我每天复习。\\
\emph{Wèile tōngguò kǎoshì, yě wèile yǐhòu qù Zhōngguó liúxué, wǒ měitiān fùxí.}\\
« Pour réussir l'examen, et aussi pour partir plus tard étudier en Chine, je révise chaque jour. »\\
A : 汉语难吗？\\
\emph{Hànyǔ nán ma ?}\\
« Le chinois est-il difficile ? »\\
B : 汉语又有意思又有用！我一边听中文歌，一边学新词。\\
\emph{Hànyǔ yòu yǒu yìsi yòu yǒuyòng ! Wǒ yìbiān tīng Zhōngwén gē, yìbiān xué xīn cí.}\\
« Le chinois est à la fois intéressant et utile ! J'écoute des chansons chinoises tout en apprenant des mots nouveaux. »
\end{textebox}

Observons trois phrases du dialogue :
\begin{itemize}
\item \zh{为了通过考试，我每天复习。} — le \emph{but} (réussir l'examen) est placé \textbf{avant} l'action (réviser) ;
\item \zh{汉语又有意思又有用。} — deux \emph{qualités simultanées} reliées par \zh{又...又...} ;
\item \zh{我一边听中文歌，一边学新词。} — deux \emph{actions simultanées} reliées par \zh{一边...一边...}.
\end{itemize}

\courssub{Le but avec 为了 (wèile)}

\begin{propriete}[Structure du but : 为了 + but ，sujet + verbe]
La structure est : \textbf{为了 + but ，+ sujet + verbe (action)}. Le but est \textbf{toujours exprimé AVANT l'action}, exactement comme en français dans « \emph{Pour réussir, je travaille} ».\\
Schéma : \zh{为了} + objectif visé \zh{，} + action menée pour l'atteindre.\\
Le but peut être un nom (\zh{为了身体健康}, « pour la santé ») ou une proposition verbale (\zh{为了学好汉语}, « pour bien apprendre le chinois »).\\
Remarque : on peut renforcer avec \zh{为} seul dans un style soutenu, mais à l'écrit scolaire on utilise \zh{为了}. La réponse à \zh{为了什么...？} (« dans quel but... ? ») se construit avec la même structure.
\end{propriete}

\begin{exemplebox}[为了 en action]
\zh{为了通过考试，我每天复习。} \emph{Wèile tōngguò kǎoshì, wǒ měitiān fùxí.} « Pour réussir l'examen, je révise chaque jour. »\\
\zh{为了身体健康，他每天运动。} \emph{Wèile shēntǐ jiànkāng, tā měitiān yùndòng.} « Pour sa santé, il fait du sport chaque jour. »\\
\zh{为了学好汉语，我们去中国留学。} \emph{Wèile xuéhǎo Hànyǔ, wǒmen qù Zhōngguó liúxué.} « Pour bien apprendre le chinois, nous partons étudier en Chine. »\\
\zh{为了买新电脑，他暑假在杜阿拉工作。} \emph{Wèile mǎi xīn diànnǎo, tā shǔjià zài Dù'ālā gōngzuò.} « Pour acheter un nouvel ordinateur, il travaille à Douala pendant les grandes vacances. »\\
\zh{为了家人的幸福，父母努力工作。} \emph{Wèile jiārén de xìngfú, fùmǔ nǔlì gōngzuò.} « Pour le bonheur de leur famille, les parents travaillent dur. »
\end{exemplebox}

Comparaison avec le français : le français dit « pour + infinitif » (\emph{pour réussir}) ou « pour + nom » (\emph{pour sa santé}) ; le chinois dit \zh{为了} + proposition complète ou nom. La différence essentielle : en chinois, le groupe du but se place \textbf{en tête de phrase}, suivi d'une virgule, et la phrase principale garde son sujet. On ne peut pas dire \zh{我每天复习为了通过考试} : cet ordre, possible en français (« je révise pour réussir »), est fautif ou très lourd en chinois standard.

\begin{attention}[为了 (but) ou 因为 (cause) ? Le test qui ne trompe pas]
Les deux mots commencent par \zh{为}, mais ils n'expriment pas la même relation :
\begin{itemize}
\item \zh{因为} (yīnwei) = \textbf{cause} : un fait passé ou présent qui explique. Il répond à « \emph{pourquoi ?} ». \zh{因为下雨，我没去市场。} « Parce qu'il a plu, je ne suis pas allé au marché. »
\item \zh{为了} (wèile) = \textbf{but} : un objectif \emph{futur visé}. Il répond à « \emph{dans quel but ?} ». \zh{为了健康，我每天运动。} « Pour la santé, je fais du sport chaque jour. »
\end{itemize}
\textbf{Test pratique} : posez la question mentalement. Si la réponse répond à « pourquoi ? » → \zh{因为}. Si elle répond à « dans quel but ? » → \zh{为了}. Erreur typique du francophone : traduire « pour » automatiquement par \zh{为了}, même quand « pour » signifie « à cause de » (« Pour la pluie, le match est annulé » = \zh{因为下雨，比赛取消了}).
\end{attention}

\begin{popfigure}
\begin{tikzpicture}[node distance=1.1cm, >=Stealth]
\node[draw=popgreen, fill=popgreenL, rounded corners, align=center, minimum width=3.6cm, minimum height=1.1cm] (but) {\zh{为了} + but visé\\ \emph{(futur, objectif)}};
\node[draw=popblue, fill=popblueL, rounded corners, align=center, minimum width=3.6cm, minimum height=1.1cm, right=2.6cm of but] (act) {sujet + verbe\\ \emph{(action menée)}};
\draw[<->, very thick, popdark] (but) -- node[above, align=center]{le but précède l'action} (act);
\node[draw=poporange, fill=poporangeL, rounded corners, align=center, minimum width=3.6cm, minimum height=1.1cm, below=1.4cm of but] (cause) {\zh{因为} + cause\\ \emph{(fait passé ou présent)}};
\draw[->, very thick, poporange] (cause) -- node[above, align=center]{la cause explique} (cause -| act.west) ;
\node[draw=popmuted, fill=popcream, rounded corners, align=center, below=1.4cm of act] (rep) {résultat ou conséquence\\ \zh{所以}...};
\end{tikzpicture}
\legende{为了 (but, vers le futur) et 因为 (cause, depuis un fait) : deux directions logiques à ne pas inverser.}
\end{popfigure}

\courssub{Coordination de deux qualités : 又...又...}

\begin{propriete}[又 + adjectif 1 + 又 + adjectif 2]
\zh{又...又...} (yòu... yòu...) relie deux \textbf{qualités simultanées} du même sujet : « à la fois... et... ».\\
Schéma : sujet + \zh{又} + adjectif (ou verbe d'état) + \zh{又} + adjectif.\\
Les deux éléments sont en général de même nature et de même « orientation » (tous deux positifs, ou tous deux négatifs).
\end{propriete}

\begin{exemplebox}[Deux qualités à la fois]
\zh{这个菜又好吃又便宜。} \emph{Zhège cài yòu hǎochī yòu piányi.} « Ce plat est à la fois bon et bon marché. »\\
\zh{这个学校又大又漂亮。} \emph{Zhège xuéxiào yòu dà yòu piàoliang.} « Cette école est à la fois grande et belle. »\\
\zh{雅温得的天气又热又湿。} \emph{Yǎwēndé de tiānqì yòu rè yòu shī.} « Le temps à Yaoundé est à la fois chaud et humide. »\\
\zh{她又聪明又努力。} \emph{Tā yòu cōngming yòu nǔlì.} « Elle est à la fois intelligente et travailleuse. »
\end{exemplebox}

\begin{savaistu}[又...又... : pas pour les actions ponctuelles]
\zh{又...又...} ne relie que des \textbf{adjectifs} ou des \textbf{verbes d'état} (comme \zh{喜欢}, \zh{有}), jamais des actions ponctuelles. On ne dit pas « il à la fois mange et dort » avec \zh{又...又...} ; pour deux actions, on utilise \zh{一边...一边...} (voir ci-dessous). Notez aussi que \zh{又} seul signifie « de nouveau » (\zh{他又来了}, « il est encore venu ») : le sens « à la fois » n'existe que dans la paire \zh{又...又...}.
\end{savaistu}

\courssub{Coordination de deux actions simultanées : 一边...一边...}

\begin{propriete}[一边 + verbe 1 ，一边 + verbe 2]
\zh{一边...一边...} (yìbiān... yìbiān...) exprime deux \textbf{actions réellement simultanées} accomplies par le même sujet : « tout en..., en même temps ».\\
Schéma : sujet + \zh{一边} + action 1 \zh{，} \zh{一边} + action 2.\\
Forme courte possible : \zh{边...边...} (\zh{他边走边说}, « il parle en marchant »).
\end{propriete}

\begin{exemplebox}[Deux actions en même temps]
\zh{他一边吃饭，一边看电视。} \emph{Tā yìbiān chīfàn, yìbiān kàn diànshì.} « Il mange tout en regardant la télévision. »\\
\zh{她一边唱歌，一边跳舞。} \emph{Tā yìbiān chànggē, yìbiān tiàowǔ.} « Elle chante en dansant. »\\
\zh{我们一边走路，一边聊天。} \emph{Wǒmen yìbiān zǒulù, yìbiān liáotiān.} « Nous marchons en bavardant. »\\
\zh{妈妈一边做饭，一边听音乐。} \emph{Māma yìbiān zuòfàn, yìbiān tīng yīnyuè.} « Maman écoute de la musique en cuisinant. »
\end{exemplebox}

\begin{methode}[Bien utiliser 一边...一边... en trois étapes]
\begin{enumerate}
\item Vérifiez que les deux actions sont \textbf{vraiment simultanées et duratives} : « écouter », « marcher », « chanter », « manger » conviennent ; une action ponctuelle (\zh{到}, « arriver » ; \zh{死}, « mourir ») est impossible.
\item Vérifiez que c'est \textbf{le même sujet} qui fait les deux actions ; sinon, utilisez deux phrases séparées.
\item Placez \zh{一边} devant chaque verbe, séparez les deux membres par une virgule chinoise \zh{，}, et ne mettez pas \zh{了} après le premier verbe (les actions sont en cours, pas accomplies).
\end{enumerate}
\end{methode}

\courssub{Entraînement oral : Projets et habitudes}

\begin{experience}[Jeu de rôle : « Mon objectif, mes méthodes »]
\textbf{Consigne :} En binôme. L'élève A interroge l'élève B sur un objectif personnel (examens, voyage, sport, cuisine...). L'élève B répond en utilisant \zh{为了} (but), \zh{又...又...} (qualités de l'activité) et \zh{一边...一边...} (méthodes simultanées).\\
\textbf{Exemple :}\\
A : \zh{你为什么每天跑步？} \emph{Nǐ wèishénme měitiān pǎobù ?}\\
B : \zh{为了身体健康，我每天跑步。跑步又累又快乐。我一边跑步，一边听音乐。}\\
\textbf{Variante :} Échangez les rôles. Notez 3 phrases de votre partenaire pour les lire à la classe.
\end{experience}

\courssub{Tableaux récapitulatifs}

\begin{center}
\begin{tabularx}{\linewidth}{|X|X|X|X|}
\hline
\rowcolor{popblueL} Caractères & Pinyin & Nature & Traduction \\
\hline
为了 & wèile & préposition & pour, afin de (but) \\
\hline
因为 & yīnwei & conjonction & parce que (cause) \\
\hline
又...又... & yòu... yòu... & coordination & à la fois... et... (qualités) \\
\hline
一边...一边... & yìbiān... yìbiān... & coordination & tout en..., en même temps (actions) \\
\hline
通过 & tōngguò & verbe & réussir (un examen), passer par \\
\hline
留学 & liúxué & verbe & étudier à l'étranger \\
\hline
健康 & jiànkāng & nom / adjectif & santé ; en bonne santé \\
\hline
便宜 & piányi & adjectif & bon marché \\
\hline
营养 & yíngyǎng & nom & nutrition, valeur nutritive \\
\hline
\end{tabularx}
\end{center}

\begin{center}
\begin{tabularx}{\linewidth}{|X|X|X|}
\hline
\rowcolor{popblueL} Structure & Exemple en caractères + pinyin & Traduction \\
\hline
为了 + but ，+ action & \zh{为了学好汉语，我们每天练习。} \emph{Wèile xuéhǎo Hànyǔ, wǒmen měitiān liànxí.} & Pour bien apprendre le chinois, nous nous entraînons chaque jour. \\
\hline
因为 + cause ，所以 + résultat & \zh{因为下雨，所以我没去。} \emph{Yīnwei xiàyǔ, suǒyǐ wǒ méi qù.} & Parce qu'il a plu, je n'y suis pas allé. \\
\hline
又 + adj. + 又 + adj. & \zh{这个菜又好吃又便宜。} \emph{Zhège cài yòu hǎochī yòu piányi.} & Ce plat est à la fois bon et bon marché. \\
\hline
一边 + V1 ，一边 + V2 & \zh{她一边唱歌，一边跳舞。} \emph{Tā yìbiān chànggē, yìbiān tiàowǔ.} & Elle chante en dansant. \\
\hline
\end{tabularx}
\end{center}

\courssub{Exercice résolu et pièges}

\begin{exoresolu}[Choisissez la bonne structure]
Complétez avec \zh{为了}, \zh{因为}, \zh{又...又...} ou \zh{一边...一边...} :\\
1. \zh{......通过考试，他每天复习。}\\
2. \zh{......下雨，比赛取消了。}\\
3. \zh{这个水果......甜......便宜。}\\
4. \zh{他......听音乐，......做作业。}
\tcblower
\textbf{Solution détaillée :}\\
1. \zh{为了通过考试，他每天复习。} — « Dans quel but révise-t-il ? Pour réussir l'examen » : objectif futur → \zh{为了}.\\
2. \zh{因为下雨，比赛取消了。} — « Pourquoi le match est-il annulé ? À cause de la pluie » : fait réel, cause → \zh{因为}.\\
3. \zh{这个水果又甜又便宜。} — deux qualités simultanées (adjectifs) → \zh{又...又...}.\\
4. \zh{他一边听音乐，一边做作业。} — deux actions duratives et simultanées du même sujet → \zh{一边...一边...}.\\
\textbf{Méthode :} identifiez d'abord la relation logique (but ? cause ? qualités ? actions ?), puis appliquez le test « pourquoi / dans quel but » pour distinguer \zh{因为} et \zh{为了}.
\end{exoresolu}

\begin{attention}[Les quatre pièges des francophones]
\begin{itemize}
\item \textbf{Ordre des mots} : ne dites pas \zh{我每天运动为了健康} sur le modèle français « je fais du sport pour ma santé » ; en chinois : \zh{为了健康，我每天运动。}
\item \textbf{Confusion but/cause} : « pour » français n'est pas toujours \zh{为了} ; s'il signifie « à cause de », c'est \zh{因为}.
\item \textbf{Actions incompatibles} : évitez \zh{他一边吃饭，一边睡觉} (impossible en pratique !) et les actions ponctuelles avec \zh{一边...一边...}.
\item \textbf{又...又... avec des actions} : \zh{又...又...} ne relie que des qualités ; pour deux actions, c'est \zh{一边...一边...}.
\end{itemize}
\end{attention}

\begin{exercice}[À vous de jouer]
Traduisez en chinois : 1. « Pour apprendre le chinois, elle regarde des films chinois. » 2. « Le ndolé est à la fois bon et nourrissant. » 3. « Nous chantons en dansant. » 4. « Parce qu'il est malade, il n'est pas venu à l'école. » Puis rédigez trois phrases libres sur vos projets avec \zh{为了}.
\end{exercice}

\begin{corrige}[Exercice]
1. \zh{为了学习汉语，她看中国电影。} \emph{Wèile xuéxí Hànyǔ, tā kàn Zhōngguó diànyǐng.}\\
2. \zh{恩多莱又好吃又有营养。} \emph{Ēnduōlái yòu hǎochī yòu yǒu yíngyǎng.}\\
3. \zh{我们一边唱歌，一边跳舞。} \emph{Wǒmen yìbiān chànggē, yìbiān tiàowǔ.}\\
4. \zh{因为他生病了，所以没来学校。} \emph{Yīnwei tā shēngbìng le, suǒyǐ méi lái xuéxiào.}\\
Phrases libres : vérifier que le but précède l'action et que le sujet de l'action est exprimé.
\end{corrige}

\begin{aretenir}[L'essentiel de la section]
\zh{为了} + but \zh{，} + action : le but se dit \textbf{avant} l'action et répond à « dans quel but ? ». \zh{因为} exprime la cause et répond à « pourquoi ? ». \zh{又...又...} coordonne deux qualités simultanées ; \zh{一边...一边...} coordonne deux actions duratives et simultanées du même sujet. Ces quatre outils permettent déjà de construire des phrases complexes claires et naturelles.
\end{aretenir}

\coursec{Synthèse, pièges des francophones et entraînement}

Dans la section précédente, nous avons étudié l'expression du but avec \zh{为了} (wèile) et la coordination simple. Nous disposons maintenant de toutes les pièces du puzzle : cause, concession, condition, but. Il est temps de rassembler ces cinq rapports logiques dans un tableau unique, de dresser la liste des pièges qui attendent le francophone, et de nous entraîner comme au Bac blanc. Cette section est votre boîte à outils finale : gardez-la sous les yeux pour toutes vos rédactions.

\courssub{Vocabulaire clé de la leçon}

\begin{center}
\begin{tabularx}{\linewidth}{|X|X|X|X|}
\hline
\rowcolor{popblueL} Caractères & Pinyin & Nature & Traduction \\
\hline
因为 & yīnwèi & conj. & parce que \\
\hline
所以 & suǒyǐ & adv. & donc, par conséquent \\
\hline
虽然 & suīrán & conj. & bien que, quoique \\
\hline
但是 & dànshì & conj. & mais, cependant \\
\hline
如果 & rúguǒ & conj. & si (condition) \\
\hline
就 & jiù & adv. & alors, dès lors \\
\hline
要是 & yàoshi & conj. & si (hypothèse, oral) \\
\hline
为了 & wèile & prép. & pour, afin de \\
\hline
条件 & tiáojiàn & n. & condition \\
\hline
原因 & yuányīn & n. & cause \\
\hline
但 & dàn & conj. & mais (plus formel/littéraire) \\
\hline
\end{tabularx}
\end{center}

\courssub{Tableau récapitulatif des cinq structures}

Observons d'abord une phrase type pour chaque rapport logique, tirée de la vie d'un élève de Terminale à Yaoundé :

\begin{exemplebox}[Cinq phrases, cinq rapports logiques]
\zh{因为明天下雨，所以足球比赛取消了。} \emph{Yīnwèi míngtiān xià yǔ, suǒyǐ zúqiú bǐsài qǔxiāo le.} « Parce qu'il pleut demain, le match de football est annulé. » (cause)\\[2pt]
\zh{虽然功课很多，但是我每天都复习。} \emph{Suīrán gōngkè hěn duō, dànshì wǒ měitiān dōu fùxí.} « Bien qu'il y ait beaucoup de devoirs, je révise chaque jour. » (concession)\\[2pt]
\zh{如果你努力，你就会成功。} \emph{Rúguǒ nǐ nǔlì, nǐ jiù huì chénggōng.} « Si tu travailles dur, tu réussiras. » (condition)\\[2pt]
\zh{要是明天下雨，我们就不去杜阿拉了。} \emph{Yàoshi míngtiān xià yǔ, wǒmen jiù bù qù Dù'ālā le.} « S'il pleut demain, nous n'irons pas à Douala. » (hypothèse, plus oral)\\[2pt]
\zh{为了通过中学毕业会考，他每天学习中文。} \emph{Wèile tōngguò zhōngxué bìyè huìkǎo, tā měitiān xuéxí Zhōngwén.} « Pour réussir le baccalauréat, il étudie le chinois chaque jour. » (but)
\end{exemplebox}

\begin{center}
\begin{tabularx}{\linewidth}{|X|X|X|X|}
\hline
\rowcolor{popblueL} Rapport logique & Corrélats & Sens & Exemple type \\
\hline
Cause & \zh{因为}P1，\zh{所以}P2 & parce que... donc & \zh{因为下雨，所以我不去。} \\
\hline
Concession & \zh{虽然}P1，\zh{但是}P2 & bien que... mais & \zh{虽然难，但是有意思。} \\
\hline
Condition & \zh{如果}P1，P2\zh{就} & si... alors & \zh{如果有空，我就来。} \\
\hline
Hypothèse (oral) & \zh{要是}P1，P2\zh{就} & si (supposons) & \zh{要是下雨，就不去。} \\
\hline
But & \zh{为了}P1，P2 & pour, afin de & \zh{为了健康，他运动。} \\
\hline
\end{tabularx}
\end{center}

\begin{center}
\begin{tabularx}{\linewidth}{|X|X|X|}
\hline
\rowcolor{popblueL} Structure & Exemple (caractères + pinyin) & Traduction \\
\hline
\zh{因为……所以……} & \zh{因为下雨，所以我不去。} (Yīnwèi xià yǔ, suǒyǐ wǒ bù qù.) & Parce qu'il pleut, je n'y vais pas. \\
\hline
\zh{虽然……但是……} & \zh{虽然很难，但是我喜欢。} (Suīrán hěn nán, dànshì wǒ xǐhuan.) & Bien que ce soit difficile, j'aime. \\
\hline
\zh{如果……就……} & \zh{如果有空，我就去。} (Rúguǒ yǒu kòng, wǒ jiù qù.) & Si j'ai le temps, j'irai. \\
\hline
\zh{要是……就……} & \zh{要是下雨，我们就不去。} (Yàoshi xià yǔ, wǒmen jiù bù qù.) & S'il pleut, nous n'irons pas. \\
\hline
\zh{为了……} & \zh{为了健康，他运动。} (Wèile jiànkāng, tā yùndòng.) & Pour la santé, il fait du sport. \\
\hline
\end{tabularx}
\end{center}

\begin{popfigure}
\begin{tikzpicture}[mindmap, grow cyclic, every node/.style={concept, align=center, font=\small},
  root concept/.append style={concept color=popblue, font=\small\bfseries},
  level 1 concept/.append style={level distance=3.6cm, sibling angle=72, font=\scriptsize}]
\node[root concept, align=center]{Phrase complexe\\复句}
  child[concept color=poporange] { node[align=center]{Cause\\因为...所以...} }
  child[concept color=popgreen] { node[align=center]{Concession\\虽然...但是...} }
  child[concept color=poppurple] { node[align=center]{Condition\\如果...就...} }
  child[concept color=poppink] { node[align=center]{Hypothèse\\要是...就...} }
  child[concept color=popgold] { node[align=center]{But\\为了...} };
\end{tikzpicture}
\legende{Carte mentale finale : les cinq rapports logiques et leurs corrélats.}
\end{popfigure}

\begin{propriete}[Règle d'or : les corrélats vont par paires]
En chinois écrit standard, les conjonctions corrélatives fonctionnent \textbf{par paires} : \zh{因为}/\zh{所以}, \zh{虽然}/\zh{但是}, \zh{如果}/\zh{就}, \zh{要是}/\zh{就}. Contrairement au français, où « parce que » et « donc » ne se combinent jamais (redondance), le chinois \emph{exige} la présence des deux éléments à l'écrit. Seul \zh{为了} (but) travaille seul. À l'oral familier, on peut laisser tomber le second correlat, mais à l'examen, écrivez toujours la paire complète.
\end{propriete}

\courssub{Les six pièges du francophone}

\begin{attention}[Erreurs fréquentes à bannir]
\begin{enumerate}
\item \textbf{Piège 1 — traduire mot à mot « parce que... donc » en supprimant un correlat.} Le français standard évite la redondance « parce que... donc », donc l'élève francophone écrit souvent *\zh{因为下雨，我没去。} à l'écrit. C'est acceptable à l'oral, mais à l'écrit il faut la paire complète : \zh{因为下雨，所以我没去。} \emph{Yīnwèi xià yǔ, suǒyǐ wǒ méi qù.}
\item \textbf{Piège 2 — oublier le second correlat.} *\zh{虽然很难，我学中文。} est fautif : \zh{但是} est obligatoire → \zh{虽然很难，但是我学中文。} De même, *\zh{如果你来，我很高兴。} sans \zh{就} sonne incomplet → \zh{如果你来，我就很高兴。}
\item \textbf{Piège 3 — mal placer} \zh{就}. \zh{就} se place \textbf{après le sujet et avant le verbe} de la seconde proposition : \zh{如果他有空，他就来。} et jamais *\zh{如果他有空，就他来。} ni *\zh{他就如果有空来。}
\item \textbf{Piège 4 — bouleverser l'ordre SVO.} Sous l'influence du français (« le match que j'ai vu »), on tend à déplacer les compléments. À l'intérieur de chaque proposition, gardez Sujet + Verbe + Objet : \zh{因为我很喜欢喀麦隆的足球，所以我每个周末都看比赛。} \emph{Yīnwèi wǒ hěn xǐhuan Kāmàilóng de zúqiú, suǒyǐ wǒ měi ge zhōumò dōu kàn bǐsài.}
\item \textbf{Piège 5 — confondre} \zh{因为} (cause : « parce que ») et \zh{为了} (but : « pour »). \zh{因为他病了，所以没来。} (il est malade → cause) ≠ \zh{为了通过考试，他努力学习。} (il veut réussir → but).
\item \textbf{Piège 6 — oublier la virgule} \zh{，} entre les deux propositions. En chinois écrit, elle est obligatoire : P1 \zh{，} P2.
\end{enumerate}
\end{attention}

\begin{methode}[Stratégie en quatre pas pour construire une phrase complexe]
\begin{enumerate}
\item \textbf{Identifier le rapport logique} entre les deux idées : cause ? concession ? condition ? but ? Demandez-vous : « est-ce une raison, un obstacle, une hypothèse ou un objectif ? »
\item \textbf{Choisir la paire de corrélats} correspondante : \zh{因为...所以...}, \zh{虽然...但是...}, \zh{如果...就...}, \zh{要是...就...} ou \zh{为了}.
\item \textbf{Écrire P1 +} \zh{，} \textbf{+ P2} en respectant l'ordre SVO dans chaque proposition, sans rien déplacer.
\item \textbf{Vérifier} : le second correlat est-il présent ? \zh{就} est-il bien après le sujet et avant le verbe ? La virgule est-elle là ? Le sujet est-il placé correctement (avant ou après \zh{因为}/\zh{如果} selon qu'il est identique ou non dans les deux propositions) ?
\end{enumerate}
\end{methode}

\begin{exoresolu}[Relier deux phrases simples]
Énoncé : reliez les deux phrases avec les corrélats appropriés.\\
1) \zh{他很累。他继续学习。} (concession)\\
2) \zh{明天下雨。我们不去市场。} (condition)
\tcblower
Solution détaillée :\\
1) Rapport logique : la fatigue devrait empêcher d'étudier, pourtant il continue → concession. On choisit \zh{虽然...但是...}. On place le sujet, on ajoute la virgule : \zh{虽然他很累，但是他继续学习。} \emph{Suīrán tā hěn lèi, dànshì tā jìxù xuéxí.} « Bien qu'il soit fatigué, il continue à étudier. »\\
2) Rapport logique : la pluie est une hypothèse dont dépend le départ → condition. On choisit \zh{如果...就...}. Attention à la place de \zh{就} : après le sujet \zh{我们}, avant le verbe \zh{不去} : \zh{如果明天下雨，我们就不去市场。} \emph{Rúguǒ míngtiān xià yǔ, wǒmen jiù bù qù shìchǎng.} « S'il pleut demain, nous n'irons pas au marché. »
\end{exoresolu}

\courssub{Exercices d'application et de transformation}

\begin{exercice}[Compléter avec le correlat manquant]
Complétez chaque phrase avec \zh{所以}, \zh{但是} ou \zh{就}.\\
1) \zh{因为喀麦隆的足球很有名，\_\_\_\_\_\_ 很多人都喜欢看比赛。}\\
2) \zh{虽然中文的汉字很难，\_\_\_\_\_\_ 我觉得很有意思。}\\
3) \zh{要是你有时间，我们\_\_\_\_\_\_ 一起去雅温得。}\\
4) \zh{如果明天天气好，我们\_\_\_\_\_\_ 去爬山。}
\end{exercice}

\begin{corrige}[Exercice : compléter]
1) \zh{所以} → \zh{因为喀麦隆的足球很有名，所以很多人都喜欢看比赛。}\\
2) \zh{但是} → \zh{虽然中文的汉字很难，但是我觉得很有意思。}\\
3) \zh{就} → \zh{要是你有时间，我们就一起去雅温得。}\\
4) \zh{就} → \zh{如果明天天气好，我们就去爬山。}
\end{corrige}

\begin{exercice}[Transformation : simple → complexe et inverse]
A. Transformez la phrase simple en phrase complexe avec \zh{为了} : \zh{他学习中文。他想以后在中国工作。}\\
B. Transformez la phrase complexe en deux phrases simples : \zh{因为下雨了，所以比赛取消了。}\\
C. Thème (français → chinois) : « Si tu étudies tous les jours, tu réussiras à l'examen. »
\end{exercice}

\begin{corrige}[Exercice : transformations]
A. \zh{为了以后在中国工作，他学习中文。} \emph{Wèile yǐhòu zài Zhōngguó gōngzuò, tā xuéxí Zhōngwén.} — Le but se place en tête avec \zh{为了}, virgule, puis la proposition principale en SVO.\\
B. \zh{下雨了。比赛取消了。} \emph{Xià yǔ le. Bǐsài qǔxiāo le.} — Deux phrases indépendantes ; le lien logique devient implicite, comme en français « Il a plu. Le match a été annulé. »\\
C. \zh{如果你每天学习，你就会通过考试。} \emph{Rúguǒ nǐ měitiān xuéxí, nǐ jiù huì tōngguò kǎoshì.} — Vérifiez : paire \zh{如果}/\zh{就}, \zh{就} après le sujet \zh{你}, virgule présente.
\end{corrige}

\begin{exercice}[Production guidée : ma préparation du Bac]
Rédigez cinq phrases sur le thème « ma préparation du baccalauréat » en utilisant au moins trois structures différentes parmi \zh{因为...所以...}, \zh{虽然...但是...}, \zh{如果...就...}, \zh{为了}. Modèle de départ : \zh{为了通过中学毕业会考，我每天都复习。}
\end{exercice}

\begin{corrige}[Exercice : production guidée (corrigé possible)]
\zh{为了通过中学毕业会考，我每天都复习功课。}\\
\zh{因为考试很难，所以我早上六点就起床学习。}\\
\zh{虽然有时候我很累，但是我从不放弃。}\\
\zh{如果我努力学习，我就一定能成功。}\\
\zh{要是我通过了考试，我的父母就会很高兴。}
\end{corrige}

\begin{savaistu}[Le saviez-tu ?]
Au Bac blanc camerounais, les questions de transformation de phrases portent régulièrement sur \zh{因为...所以...} et \zh{如果...就...} : on vous donne deux phrases simples et on vous demande de les relier, ou l'inverse. Les correcteurs sanctionnent presque toujours deux fautes précises : l'absence du second correlat (\zh{所以} ou \zh{就}) et la mauvaise place de \zh{就}. Entraînez-vous avec la stratégie en quatre pas : elle vous garantit de ne rien oublier le jour de l'examen.
\end{savaistu}

\begin{aretenir}[L'essentiel de la leçon]
Une phrase complexe chinoise relie deux propositions par une \textbf{paire de corrélats} : \zh{因为...所以...} (cause), \zh{虽然...但是...} (concession), \zh{如果...就...} et \zh{要是...就...} (condition et hypothèse), \zh{为了} (but, seul). Chaque proposition garde l'ordre SVO, une virgule \zh{，} sépare les deux propositions, et \zh{就} se place après le sujet, avant le verbe. À l'écrit, ne supprimez jamais le second correlat : c'est la règle d'or.
\end{aretenir}

\newpage

\coursec{Activité d'intégration}

\coursec{Leçon 22 — Grammaire : Les phrases complexes}

\courssub{Objectifs de l'activité}

Cette activité d'intégration clôt la leçon sur les phrases complexes. Elle vise à faire passer les élèves de la reconnaissance des structures à leur \cle{production autonome} dans un texte réel et personnel. À l'issue de l'activité, chaque élève sera capable de :

\begin{itemize}
\item rédiger une lettre courte et cohérente en chinois (8 à 10 phrases) à destination d'un correspondant ;
\item mobiliser correctement les cinq structures de phrases complexes étudiées : \zh{因为……所以……}, \zh{虽然……但是……}, \zh{如果……就……}, \zh{要是……就……} et \zh{为了……} ;
\item respecter l'ordre des mots chinois et la place des corrélats (avant le sujet ou après, selon les cas) ;
\item relire et corriger la production d'un pair en repérant les corrélats et en vérifiant leur emploi ;
\item poser des questions de compréhension simples sur un texte lu.
\end{itemize}

\courssub{Situation}

Dans le cadre du club chinois du lycée de Nkol-Eton à Yaoundé, votre professeur a mis en place un partenariat avec une classe de terminale d'un lycée de Chengdu, en Chine. Chaque élève camerounais a été associé à un correspondant chinois. Cette semaine, vous devez envoyer votre première vraie lettre : votre correspondant, Lǐ Míng, vous a déjà écrit et vous pose des questions sur vos projets après le baccalauréat. Voici un extrait de sa lettre :

\begin{textebox}[Extrait de la lettre de Lǐ Míng]
亲爱的朋友，你好！\\
Qīn'ài de péngyou, nǐ hǎo !\\
Cher ami, bonjour !\\[4pt]
你为什么学习汉语？是因为你喜欢中国文化吗？\\
Nǐ wèishénme xuéxí Hànyǔ ? Shì yīnwèi nǐ xǐhuan Zhōngguó wénhuà ma ?\\
Pourquoi apprends-tu le chinois ? Est-ce parce que tu aimes la culture chinoise ?\\[4pt]
如果你通过考试，你就打算去中国留学吗？\\
Rúguǒ nǐ tōngguò kǎoshì, nǐ jiù dǎsuàn qù Zhōngguó liúxué ma ?\\
Si tu réussis l'examen, as-tu l'intention d'aller étudier en Chine ?\\[4pt]
学习汉语难不难？告诉我你的计划吧！\\
Xuéxí Hànyǔ nán bu nán ? Gàosu wǒ nǐ de jìhuà ba !\\
Apprendre le chinois, est-ce difficile ? Raconte-moi tes projets !\\[4pt]
你的朋友，李明\\
Nǐ de péngyou, Lǐ Míng\\
Ton ami, Li Ming
\end{textebox}

Vous remarquez que Lǐ Míng utilise lui-même des phrases complexes (\zh{因为}, \zh{如果……就……}). À votre tour de lui répondre en montrant que vous maîtrisez ces structures !

\courssub{Consignes}

\textbf{Tâche 1 — Compréhension de la lettre reçue.} Lisez la lettre de Lǐ Míng. Soulignez les deux structures de phrases complexes qu'elle contient et recopiez-les avec leur traduction. Relevez les trois questions qu'il vous pose : votre réponse devra y répondre une par une.

\textbf{Tâche 2 — Préparation du plan.} Sur votre brouillon, complétez le plan suivant en français, puis en quelques mots chinois :

\begin{enumerate}
\item Salutations et formule d'ouverture (une phrase).
\item Pourquoi j'apprends le chinois : une phrase avec \zh{因为……所以……}.
\item Mes difficultés et mes efforts : une phrase avec \zh{虽然……但是……}.
\item Mes projets si je réussis le Bac : une phrase avec \zh{如果……就……} et une phrase avec \zh{要是……就……}.
\item Mon but à long terme : une phrase avec \zh{为了……}.
\item Conclusion et formule de politesse (une ou deux phrases).
\end{enumerate}

\textbf{Tâche 3 — Rédaction.} Rédigez votre lettre de 8 à 10 phrases en caractères chinois. Chacune des cinq structures doit apparaître \cle{au moins une fois}. Vous pouvez réutiliser le vocabulaire de la leçon (études, examen, université, travail, voyage, Chine, Cameroun).

\textbf{Tâche 4 — Vérification personnelle.} Relisez votre lettre et cochez sur la grille suivante : (a) chaque corrélat est bien placé ; (b) \zh{因为} et \zh{所以} ne traduisent pas deux fois « parce que » ; (c) l'ordre Sujet + Verbe + Complément est respecté ; (d) la ponctuation chinoise est utilisée (，。).

\textbf{Tâche 5 — Échange en binôme.} Échangez votre lettre avec votre voisin. Le partenaire doit :

\begin{enumerate}
\item souligner tous les corrélats utilisés et vérifier qu'ils sont correctement employés ;
\item signaler à l'oral une éventuelle erreur d'ordre des mots ;
\item poser deux questions de compréhension en chinois, par exemple : \zh{你为什么想去中国？} \emph{Nǐ wèishénme xiǎng qù Zhōngguó ?} « Pourquoi veux-tu aller en Chine ? »
\end{enumerate}

\textbf{Tâche 6 — Mise en commun.} Les deux ou trois meilleures lettres sont lues à voix haute devant la classe, corrigées collectivement au tableau, puis affichées au mur du club chinois.

\courssub{Ressources à mobiliser}

\begin{aretenir}[Les cinq structures à utiliser]
\begin{itemize}
\item \textbf{Cause} : \zh{因为 + cause，所以 + résultat。} — \zh{因为汉语很有用，所以我学习汉语。} \emph{Yīnwèi Hànyǔ hěn yǒuyòng, suǒyǐ wǒ xuéxí Hànyǔ.} « Parce que le chinois est très utile, je l'apprends. »
\item \textbf{Concession} : \zh{虽然 + concession，但是 + fait principal。} — \zh{虽然汉字很难，但是我每天练习。} \emph{Suīrán Hànzì hěn nán, dànshì wǒ měitiān liànxí.} « Bien que les caractères soient difficiles, je m'entraîne chaque jour. »
\item \textbf{Condition} : \zh{如果 + condition，就 + conséquence。} — \zh{如果我通过考试，我就去中国。} \emph{Rúguǒ wǒ tōngguò kǎoshì, wǒ jiù qù Zhōngguó.} « Si je réussis l'examen, j'irai en Chine. »
\item \textbf{Hypothèse familière} : \zh{要是 + hypothèse，就 + conséquence。} — \zh{要是我有时间，我就去成都看你。} \emph{Yàoshi wǒ yǒu shíjiān, wǒ jiù qù Chéngdū kàn nǐ.} « Si j'ai le temps, j'irai te voir à Chengdu. »
\item \textbf{But} : \zh{为了 + but，+ action。} — \zh{为了找到好工作，我努力学习。} \emph{Wèile zhǎodào hǎo gōngzuò, wǒ nǔlì xuéxí.} « Pour trouver un bon travail, j'étudie avec effort. »
\end{itemize}
\end{aretenir}

\begin{attention}[Pièges à éviter dans la lettre]
\begin{itemize}
\item Ne dites pas \zh{因为……所以……} en pensant à « parce que... donc... » : en français on n'utilise qu'un seul mot, en chinois les deux corrélats se répondent naturellement. À l'inverse, n'oubliez pas le second corrélat : \zh{因为……} seul laisse la phrase inachevée.
\item \zh{就} dans \zh{如果……就……} et \zh{要是……就……} se place \textbf{avant le verbe} de la seconde proposition, jamais en fin de phrase : \zh{我就去中国} et non \zh{我去中国就}.
\item \zh{为了} introduit le but \textbf{avant} l'action principale : \zh{为了学习汉语，我每天练习。} L'ordre inverse du français (« je m'entraîne pour apprendre ») est une erreur fréquente des francophones.
\item Si les deux propositions ont le même sujet, celui-ci peut être omis dans la seconde : \zh{如果我有时间，就去中国旅游。}
\end{itemize}
\end{attention}

\begin{center}
\begin{tabularx}{\linewidth}{|X|X|X|X|}
\hline
\rowcolor{popblueL} Caractères & Pinyin & Nature & Traduction \\
\hline
通过 & tōngguò & verbe & réussir (un examen) \\
\hline
考试 & kǎoshì & nom & examen \\
\hline
留学 & liúxué & verbe & étudier à l'étranger \\
\hline
大学 & dàxué & nom & université \\
\hline
计划 & jìhuà & nom & projet, plan \\
\hline
努力 & nǔlì & adverbe / adjectif & avec effort, assidu \\
\hline
有用 & yǒuyòng & adjectif & utile \\
\hline
机会 & jīhuì & nom & occasion, chance \\
\hline
翻译 & fānyì & nom / verbe & traduction, traduire ; traducteur, interprète \\
\hline
希望 & xīwàng & verbe / nom & espérer ; espoir \\
\hline
打算 & dǎsuàn & verbe & avoir l'intention de, prévoir \\
\hline
毕业 & bìyè & verbe & obtenir son diplôme, finir ses études \\
\hline
经济 & jīngjì & nom & économie \\
\hline
雅温得 & Yǎwēndé & nom propre & Yaoundé \\
\hline
喀麦隆 & Kāmàilóng & nom propre & Cameroun \\
\hline
成都 & Chéngdū & nom propre & Chengdu \\
\hline
\end{tabularx}
\end{center}

\begin{savaistu}[Le \cle{笔友} (bǐyǒu) : correspondance Chine-Cameroun]
L'échange de lettres entre élèves camerounais et chinois (souvent via WeChat ou email aujourd'hui) est une pratique pédagogique courante. En chinois, on distingue \zh{信} (xìn, lettre formelle/papier) et \zh{消息} (xiāoxi, message instantané). La structure formelle d'une lettre chinoise place l'expéditeur en bas à gauche (\zh{你的朋友}) et le destinataire en haut (\zh{亲爱的李明，你好！}). Les formules de politesse sont codifiées : \zh{请给我回信} (qǐng gěi wǒ huíxìn, « je t'en prie, réponds-moi ») marque l'attente de réponse sans familiarité excessive.

\end{savaistu}

\courssub{Proposition de production et corrigé commenté}

\begin{exoresolu}[Modèle de lettre : « Mes projets après le Bac »]
Rédigez la lettre de réponse à Lǐ Míng en 8 à 10 phrases, en utilisant les cinq structures de phrases complexes.
\tcblower
\textbf{Production modèle :}\\[4pt]
亲爱的李明，你好！\\
Qīn'ài de Lǐ Míng, nǐ hǎo !\\
Cher Li Ming, bonjour !\\[3pt]
谢谢你的信。因为汉语很有用，所以我在雅温得的中学学习汉语。\\
Xièxie nǐ de xìn. Yīnwèi Hànyǔ hěn yǒuyòng, suǒyǐ wǒ zài Yǎwēndé de zhōngxué xuéxí Hànyǔ.\\
Merci pour ta lettre. Parce que le chinois est très utile, je l'apprends dans mon lycée de Yaoundé.\\[3pt]
虽然汉字很难，但是我每天练习写汉字。\\
Suīrán Hànzì hěn nán, dànshì wǒ měitiān liànxí xiě Hànzì.\\
Bien que les caractères soient difficiles, je m'entraîne à les écrire chaque jour.\\[3pt]
如果我通过毕业考试，我就去大学学习汉语和经济。\\
Rúguǒ wǒ tōngguò bìyè kǎoshì, wǒ jiù qù dàxué xuéxí Hànyǔ hé jīngjì.\\
Si je réussis le baccalauréat, j'irai à l'université étudier le chinois et l'économie.\\[3pt]
要是我有机会，我就去中国留学，也想去成都看你。\\
Yàoshi wǒ yǒu jīhuì, wǒ jiù qù Zhōngguó liúxué, yě xiǎng qù Chéngdū kàn nǐ.\\
Si j'en ai l'occasion, j'irai étudier en Chine et je voudrais aussi aller te voir à Chengdu.\\[3pt]
为了找到好工作，我现在努力学习。\\
Wèile zhǎodào hǎo gōngzuò, wǒ xiànzài nǔlì xuéxí.\\
Pour trouver un bon travail, j'étudie avec effort maintenant.\\[3pt]
我希望以后当翻译，帮助喀麦隆和中国的朋友。\\
Wǒ xīwàng yǐhòu dāng fānyì, bāngzhù Kāmàilóng hé Zhōngguó de péngyou.\\
J'espère devenir plus tard traducteur et aider les amis du Cameroun et de la Chine.\\[3pt]
请给我回信！\\
Qǐng gěi wǒ huíxìn !\\
Réponds-moi vite, je t'en prie !\\[3pt]
你的朋友\\
Nǐ de péngyou\\
Ton ami\\[6pt]
\textbf{Commentaire des choix de langue :}
\begin{itemize}
\item La lettre compte dix phrases et emploie les cinq structures : \zh{因为……所以……} (cause de l'apprentissage), \zh{虽然……但是……} (difficulté surmontée), \zh{如果……就……} (projet conditionné au Bac), \zh{要是……就……} (hypothèse plus rêvée, ton personnel), \zh{为了……} (but professionnel).
\item \zh{就} est bien placé avant le verbe dans les deux phrases conditionnelles (\zh{我就去}, \zh{我就去}).
\item \zh{为了找到好工作} ouvre la phrase : le but précède l'action, conformément à l'ordre chinois.
\item Le complément de lieu \zh{在雅温得的中学} est placé avant le verbe \zh{学习}, comme l'exige la grammaire chinoise (Sujet + Lieu + Verbe).
\item Les formules d'ouverture (\zh{亲爱的……，你好！}) et de clôture (\zh{你的朋友}) reproduisent le modèle de la lettre reçue : c'est la structure attendue d'une correspondance chinoise simple.
\item Le vocabulaire \zh{毕业考试} (examen de fin d'études / Bac), \zh{经济} (économie), \zh{打算} (avoir l'intention) enrichit le lexique de la leçon.
\end{itemize}
\end{exoresolu}

\courssub{Bilan de l'activité}

À l'issue de cette activité, l'enseignant anime un court échange collectif : Quelles structures ont été les plus faciles à employer ? Lesquelles ont posé problème ? Les erreurs relevées pendant l'échange en binôme sont notées au tableau et corrigées ensemble — le plus souvent : l'oubli du second corrélat (\zh{所以}, \zh{但是}, \zh{就}), la mauvaise place de \zh{就}, et le calque de l'ordre français pour le but.

\begin{methode}[Grille d'auto-évaluation de la lettre]
\begin{enumerate}
\item Ma lettre compte entre 8 et 10 phrases complètes.
\item Les cinq structures complexes apparaissent, chacune avec ses deux corrélats.
\item \zh{就} est avant le verbe ; \zh{为了} est en tête de phrase.
\item L'ordre Sujet + (complément de lieu/temps) + Verbe + Complément est respecté.
\item J'ai répondu aux trois questions de Lǐ Míng (pourquoi j'apprends, difficultés, projets).
\item J'ai utilisé la ponctuation chinoise (，。？！) et les formules de lettre.
\end{enumerate}
\end{methode}

Cette production écrite, affichée ensuite au mur de la classe, constitue une trace concrète de la maîtrise des phrases complexes : elle pourra être reprise et enrichie lors des révisions générales de fin d'année, en préparation de l'épreuve écrite du baccalauréat.

\newpage

\coursec{Méthodes et exercices résolus}

\coursec{Leçon 22 — Grammaire : phrases complexes}

\begin{textebox}[Observation : Dialogue au lycée de Yaoundé]
\textbf{Contexte :} Paul et Marie, élèves de Terminale A au Lycée de Nkolbisson, discutent de leur avenir et de l'apprentissage du chinois.\\
\zh{保罗：玛丽，你为什么要学汉语？} \emph{Bǎoluó: Mǎlì, nǐ wèishénme yào xué Hànyǔ?}\\
\zh{玛丽：因为中国经济发展很快，所以学汉语很有用。虽然汉字很难写，但是语法不复杂。} \emph{Mǎlì: Yīnwèi Zhōngguó jīngjì fāzhǎn hěn kuài, suǒyǐ xué Hànyǔ hěn yǒuyòng. Suīrán Hànzì hěn nán xiě, dànshì yǔfǎ bù fùzá.}\\
\zh{保罗：如果我考上大学，我就想去北京留学。为了这个目标，我每天学习两个小时。} \emph{Bǎoluó: Rúguǒ wǒ kǎo shàng dàxué, wǒ jiù xiǎng qù Běijīng liúxué. Wèile zhège mùbiāo, wǒ měitiān xuéxí liǎng gè xiǎoshí.}\\
\zh{玛丽：好主意！要是你努力，你就一定能成功。我们一起加油吧！} \emph{Mǎlì: Hǎo zhǔyì! Yàoshi nǐ nǔlì, nǐ jiù yídìng néng chénggōng. Wǒmen yìqǐ jiāyóu ba!}\\
\textbf{Traduction :} Paul : « Marie, pourquoi apprends-tu le chinois ? » — Marie : « Parce que l'économie chinoise se développe vite, donc apprendre le chinois est très utile. Bien que les caractères soient difficiles à écrire, la grammaire n'est pas compliquée. » — Paul : « Si je suis admis à l'université, je veux aller étudier à Pékin. Pour cet objectif, j'étudie deux heures par jour. » — Marie : « Bonne idée ! Si tu fais des efforts, tu réussiras certainement. Encourageons-nous mutuellement ! »
\end{textebox}

\begin{center}
\begin{tabularx}{\linewidth}{|X|X|X|X|}
\hline
\rowcolor{popblueL} \textbf{Caractères} & \textbf{Pinyin} & \textbf{Nature} & \textbf{Traduction} \\
\hline
因为 & yīnwèi & Conj. & parce que, car \\
\hline
所以 & suǒyǐ & Adv. & donc, par conséquent \\
\hline
虽然 & suīrán & Conj. & bien que, quoique \\
\hline
但是 & dànshì & Conj. & mais, cependant \\
\hline
但 & dàn & Conj. & mais (court) \\
\hline
可是 & kěshì & Conj. & mais, pourtant \\
\hline
如果 & rúguǒ & Conj. & si (condition) \\
\hline
要是 & yàoshi & Conj. & si (oral, familier) \\
\hline
就 & jiù & Adv. & alors, dès lors (marque la conséquence) \\
\hline
为了 & wèile & Prép. & pour, afin de (but) \\
\hline
和 & hé & Conj. & et (noms) \\
\hline
又……又…… & yòu… yòu… & Structure & à la fois… et… (qualités) \\
\hline
一边……一边…… & yìbiān… yìbiān… & Structure & tout en… / en même temps que… (actions) \\
\hline
穷 & qióng & Adj. & pauvre \\
\hline
大方 & dàfang & Adj. & généreux \\
\hline
词典 & cídiǎn & N. & dictionnaire \\
\hline
本 & běn & Classif. & classificateur pour livres \\
\hline
热闹 & rènao & Adj. & animé, vivant \\
\hline
找到 & zhǎodào & V. & trouver (résultat) \\
\hline
将来 & jiānglái & N. & avenir, futur \\
\hline
打开 & dǎkāi & V. & ouvrir \\
\hline
门 & mén & N. & porte ; fig. opportunité \\
\hline
\end{tabularx}
\end{center}

\courssub{Notion : La phrase complexe en chinois}

\begin{definition}[Phrase simple vs phrase complexe]
Une \cle{phrase simple} (单句, dānjù) ne contient qu'un seul prédicat (verbe ou adjectif prédicatif). Exemple : \zh{我学习汉语。} (Un seul verbe \zh{学习}).\\
Une \cle{phrase complexe} (复句, fùjù) est formée de deux ou plusieurs propositions (分句, fènjù) liées par des connecteurs logiques (关联词, guānxì cí). Chaque proposition a son propre prédicat.\\
Exemple : \zh{因为我想去中国，所以我学习汉语。} (Deux prédicats : \zh{想去} et \zh{学习}, liés par \zh{因为……所以……}).
\end{definition}

\begin{propriete}[Caractéristiques de la phrase complexe chinoise]
\begin{enumerate}
\item \textbf{Parataxe vs Hypotaxe} : Le chinois privilégie la \emph{parataxe} (liaison par le sens, l'ordre des mots et des particules) alors que le français privilégie l'\emph{hypotaxe} (liaison par des conjonctions de subordination explicites).
\item \textbf{Ordre logique} : La cause précède la conséquence (\zh{因为……所以……}), la condition précède la conséquence (\zh{如果……就……}), le but précède l'action (\zh{为了……}). Le français permet l'inversion (« Je reste car il pleut » / « Car il pleut, je reste »), le chinois \textbf{interdit} l'inversion de la proposition principale et de la subordonnée pour ces structures.
\item \textbf{Paire corrélative} : La plupart des connecteurs fonctionnent \textbf{par paires} (corrélatifs) : \zh{因为} appelle \zh{所以}, \zh{虽然} appelle \zh{但是}, \zh{如果} appelle \zh{就}. À l'écrit formel (Bac), les deux éléments sont obligatoires.
\item \textbf{Gestion du sujet} : Si le sujet est commun aux deux propositions, il se place \textbf{avant le premier connecteur} (ex: \zh{他因为忙，所以不来}). Si les sujets diffèrent, chaque proposition garde son sujet.
\end{enumerate}
\end{propriete}

\courssub{Fiche méthode 1 : Construire une phrase de cause avec \zh{因为……所以……}}

\begin{methode}[Exprimer la cause et la conséquence]
Pour relier une cause à sa conséquence en chinois, suis ces étapes :
\begin{enumerate}
\item \textbf{Identifier} dans la phrase française quelle est la cause (« parce que », « car », « comme ») et quelle est la conséquence (« donc », « par conséquent », « c'est pourquoi »).
\item \textbf{Placer la cause en premier} : le chinois suit l'ordre logique immuable \textbf{Cause $\rightarrow$ Conséquence}. Ne calque jamais l'ordre français « Conséquence + parce que Cause ».
\item \textbf{Construire le schéma} : \cle{因为} + cause ， \cle{所以} + conséquence 。
\item \textbf{Gérer le sujet} :
      \begin{itemize}
      \item Sujet commun $\rightarrow$ placé \textbf{avant} \zh{因为} : \zh{他因为生病，所以没来。}
      \item Sujets différents $\rightarrow$ un sujet dans chaque proposition : \zh{因为下雨，所以地湿了。}
      \end{itemize}
\item \textbf{Alléger (oral seulement)} : à l'oral, on peut garder \zh{因为} seul (\zh{因为下雨，我不去}) ou \zh{所以} seul (\zh{下雨了，所以我不去}). \textbf{À l'écrit d'examen : utilisez la paire complète.}
\end{enumerate}
\textbf{Structure canonique} : [Sujet commun +] \zh{因为} + Cause ， \zh{所以} + Conséquence 。
\end{methode}

\begin{exemplebox}[Exemples variés Cause / Conséquence]
\begin{itemize}
\item \zh{因为今天是星期天，所以我们不上课。} \emph{Yīnwèi jīntiān shì xīngqítiān, suǒyǐ wǒmen bù shàngkè.} « Comme aujourd'hui est dimanche, nous n'avons pas cours. » (Sujets différents)
\item \zh{我因为没钱，所以买不起那本书。} \emph{Wǒ yīnwèi méi qián, suǒyǐ mǎi bù qǐ nà běn shū.} « Comme je n'ai pas d'argent, je ne peux pas m'acheter ce livre. » (Sujet commun \zh{我} avant \zh{因为})
\item \zh{由于交通堵塞，所以我迟到了。} \emph{Yóuyú jiāotōng dǔsè, suǒyǐ wǒ chídào le.} « En raison des embouteillages, j'ai eu du retard. » (\zh{由于} = formel pour \zh{因为}, niveau HSK 4+)
\end{itemize}
\end{exemplebox}

\begin{exoresolu}[Transformation : relier par la cause]
Énoncé — Relie les deux phrases suivantes avec \zh{因为……所以……}, puis traduis ta phrase en français.\\
1. \zh{他很忙。他不能来看你。} \emph{Tā hěn máng. Tā bù néng lái kàn nǐ.}\\
2. \zh{路太远。我们坐车去。} \emph{Lù tài yuǎn. Wǒmen zuò chē qù.}\\
3. \zh{她喜欢中国文化。她学习汉语。} \emph{Tā xǐhuan Zhōngguó wénhuà. Tā xuéxí Hànyǔ.}
\tcblower
Solution détaillée :
\begin{enumerate}
\item Cause : « il est occupé » ; Conséquence : « il ne peut pas venir ». Sujet commun \zh{他} $\rightarrow$ avant \zh{因为}.\\
\zh{他因为很忙，所以不能来看你。} \emph{Tā yīnwèi hěn máng, suǒyǐ bù néng lái kàn nǐ.} « Comme il est très occupé, il ne peut pas venir te voir. »
\item Cause : « la route est trop longue » ; Conséquence : « nous prenons le véhicule ». Sujet commun \zh{我们} $\rightarrow$ avant \zh{因为}.\\
\zh{我们因为路太远，所以坐车去。} \emph{Wǒmen yīnwèi lù tài yuǎn, suǒyǐ zuò chē qù.} « Comme la route est trop longue, nous y allons en voiture. »
\item Cause : « elle aime la culture chinoise » ; Conséquence : « elle apprend le chinois ». Sujet commun \zh{她} $\rightarrow$ avant \zh{因为}.\\
\zh{她因为喜欢中国文化，所以学习汉语。} \emph{Tā yīnwèi xǐhuan Zhōngguó wénhuà, suǒyǐ xuéxí Hànyǔ.} « Parce qu'elle aime la culture chinoise, elle apprend le chinois. »
\end{enumerate}
\textbf{Conclusion} : quand le sujet est identique, on ne le répète pas ; on le place avant \zh{因为}. Ne jamais traduire « parce que... donc... » mot à mot en gardant deux sujets (calque français).
\end{exoresolu}

\courssub{Fiche méthode 2 : Construire une phrase de concession avec \zh{虽然……但是……}}

\begin{methode}[Exprimer la concession]
La concession oppose un fait concédé (reconnu vrai) à un fait principal qui va à l'encontre de l'attente logique.
\begin{enumerate}
\item \textbf{Repérer} l'opposition en français (mots-repères : bien que, malgré, même si, pourtant, mais, cependant).
\item \textbf{Placer le fait concédé} après \cle{虽然} (suīrán) et le fait principal après \cle{但是} (dànshì) ou ses variantes \zh{但} (dàn) / \zh{可是} (kěshì).
\item \textbf{Appliquer le schéma} : \zh{虽然} + fait concédé ， \zh{但是} + fait principal 。
\item \textbf{Gérer le sujet} : Sujet commun $\rightarrow$ avant \zh{虽然} ; Sujets différents $\rightarrow$ un sujet par proposition.
\item \textbf{Règle d'or} : \textbf{En chinois, la paire \zh{虽然……但是……} est standard et obligatoire à l'écrit.} Le français interdit « Bien que... mais... » (pléonasme), mais le chinois \textbf{exige} les deux marqueurs pour bien délimiter les deux propositions.
\end{enumerate}
\textbf{Structure canonique} : [Sujet commun +] \zh{虽然} + Fait concédé ， \zh{但是} + Fait principal 。
\end{methode}

\begin{exemplebox}[Variantes du connecteur de concession]
\begin{itemize}
\item \zh{虽然他累了，但他坚持跑完了。} \emph{Suīrán tā lèi le, dàn tā jiānchí pǎo wán le.} « Bien qu'il soit fatigué, il a insisté pour finir la course. » (\zh{但} = forme courte de \zh{但是})
\item \zh{虽然下雨，可是他们还在踢足球。} \emph{Suīrán xià yǔ, kěshì tāmen hái zài tī zúqiú.} « Bien qu'il pleuve, ils jouent encore au foot. » (\zh{可是} = plus insistatif sur le contraste)
\item \zh{尽管天气冷，他也要去。} \emph{Jǐnguǎn tiānqì lěng, tā yě yào qù.} « Malgré le froid / Qu'il fasse froid, il y va quand même. » (\zh{尽管} = même si / malgré, plus fort, HSK 4).
\end{itemize}
\end{exemplebox}

\begin{exoresolu}[Grammaire : compléter et traduire]
Énoncé — Complète avec \zh{虽然} ou \zh{但是}, puis traduis en français.\\
1. \zh{……天气很冷，……他还要去上学。}\\
2. \zh{……汉语很难，……我觉得很有意思。}\\
3. Traduis en chinois : « Bien qu'il soit pauvre, il est très généreux. »
\tcblower
Solution détaillée :
\begin{enumerate}
\item Fait concédé : « il fait froid » ; Fait principal : « il va quand même à l'école » (\zh{还要} = encore vouloir/doit).\\
\zh{虽然天气很冷，但是他还要去上学。} \emph{Suīrán tiānqì hěn lěng, dànshì tā hái yào qù shàngxué.} « Bien qu'il fasse très froid, il doit quand même aller à l'école. »
\item Fait concédé : « le chinois est difficile » ; Fait principal : « je trouve ça intéressant ».\\
\zh{虽然汉语很难，但是我觉得很有意思。} \emph{Suīrán Hànyǔ hěn nán, dànshì wǒ juéde hěn yǒu yìsi.} « Bien que le chinois soit difficile, je le trouve très intéressant. »
\item « Bien que » $\rightarrow$ \zh{虽然} ; « pourtant » sous-entendu $\rightarrow$ \zh{但是} ; « pauvre » $\rightarrow$ \zh{穷} (qióng) ; « généreux » $\rightarrow$ \zh{大方} (dàfang). Sujet commun \zh{他}.\\
\zh{虽然他很穷，但是他很大方。} \emph{Suīrán tā hěn qióng, dànshì tā hěn dàfang.}
\end{enumerate}
\textbf{Conclusion} : la paire \zh{虽然……但是……} fonctionne en bloc indissociable à l'écrit. Supprimer \zh{但是} est toléré à l'oral rapide mais affaiblit la structure à l'examen.
\end{exoresolu}

\courssub{Fiche méthode 3 : Construire une hypothèse avec \zh{如果……就……} et \zh{要是……就……}}

\begin{methode}[Exprimer la condition et l'hypothèse]
\begin{enumerate}
\item \textbf{Repérer} la condition en français (« si », « au cas où ») et sa conséquence (« alors », « dans ce cas »).
\item \textbf{Choisir le mot de condition} :
      \begin{itemize}
      \item \cle{如果} (rúguǒ) : neutre, écrit et oral, standard pour le Bac.
      \item \cle{要是} (yàoshi) : plus oral, familier, conversationnel.
      \item \zh{假如} (jiǎrú) : « supposé que », plus littéraire/hypothétique (HSK 4+).
      \end{itemize}
\item \textbf{Appliquer le schéma} : \zh{如果} + Condition ， (Sujet) + \cle{就} + Conséquence 。
      \begin{itemize}
      \item \textbf{Place cruciale de \zh{juste après le sujet et avant le verbe (ou l'auxiliaire)} de la conséquence.}
      \item Ne jamais placer \zh{就} avant le sujet (\textcolor{red}{\zh{如果下雨，就我不去}} $\rightarrow$ \textbf{FAUX}).
      \end{itemize}
\item \textbf{Pas de conditionnel} : Le chinois n'a pas de morphologie verbale pour l'hypothétique/irréel. La valeur « -rais » est portée uniquement par la structure \zh{如果……就……}. N'ajoute aucun mot pour traduire le conditionnel français.
\item \textbf{Variante \zh{的话} (dehuà)} : On peut ajouter \zh{的话} à la fin de la condition pour marquer l'hypothèse : \zh{如果明天下雨的话，我就不去。} (Optionnel, renforce l'hypothèse).
\end{enumerate}
\textbf{Structure canonique} : \zh{如果} + Condition ， [Sujet +] \zh{就} + Verbe (+ Complément) 。
\end{methode}

\begin{exemplebox}[La place de \zh{就} : cas pratiques]
\begin{itemize}
\item Sujet exprimé : \zh{如果你来，我就去。} \emph{Rúguǒ nǐ lái, wǒ jiù qù.} (Sujet \zh{我} + \zh{就} + Verbe \zh{去})
\item Sujet omis (identique) : \zh{如果下雨，就不去了。} \emph{Rúguǒ xià yǔ, jiù bù qù le.} (\zh{就} directement devant \zh{不去})
\item Avec auxiliaire : \zh{如果你努力，你就能考上。} \emph{Rúguǒ nǐ nǔlì, nǐ jiù néng kǎo shàng.} (\zh{就} + \zh{能} + Verbe)
\item Avec adverbe de degré : \zh{如果很贵，我就不买。} \emph{Rúguǒ hěn guì, wǒ jiù bù mǎi.}
\end{itemize}
\end{exemplebox}

\begin{exoresolu}[Traduction de version : l'hypothèse]
Énoncé — Traduis en chinois :\\
1. « S'il pleut demain, je resterai à la maison. »\\
2. « Si tu étudies bien, tu réussiras à l'examen. »\\
3. « Si j'avais de l'argent, j'achèterais un dictionnaire. » (hypothèse irréelle : la structure reste la même !)
\tcblower
Solution détaillée :
\begin{enumerate}
\item Condition : \zh{明天下雨} ; Conséquence : \zh{就} + \zh{待} (dāi, rester) + \zh{在家里}.\\
\zh{如果明天下雨，我就待在家里。} \emph{Rúguǒ míngtiān xià yǔ, wǒ jiù dāi zài jiā lǐ.}\\
Justification : \zh{就} se place après le sujet \zh{我} et avant le verbe \zh{待}.
\item Condition : \zh{你好好学习} (\zh{好好} = bien, sérieusement) ; Conséquence : \zh{就} + \zh{会} (probabilité future) + \zh{考好}.\\
\zh{如果你好好学习，就会考好。} \emph{Rúguǒ nǐ hǎohāo xuéxí, jiù huì kǎo hǎo.}\\
Variante avec sujet répété : \zh{如果你好好学习，你就会考好。}
\item \textbf{Piège classique} : Le français utilise le conditionnel passé (« aurais », « achèterais »), le chinois \textbf{ne change rien}. L'irréel est exprimé par la seule structure. Classificateur \zh{本} (běn) obligatoire pour \zh{词典}.\\
\zh{如果我有钱，我就买一本词典。} \emph{Rúguǒ wǒ yǒu qián, wǒ jiù mǎi yì běn cídiǎn.}
\end{enumerate}
\textbf{Conclusion} : Pas de temps hypothétique en chinois ; toute la valeur repose sur le couple \zh{如果……就……}. La place de \zh{就} (après sujet, avant verbe/auxiliaire/négation) est le point n°1 contrôlé au Bac.
\end{exoresolu}

\courssub{Fiche méthode 4 : Exprimer le but avec \zh{为了} et coordonner simplement}

\begin{methode}[Exprimer le but et coordonner]
\begin{enumerate}
\item \textbf{But avec \cle{为了} (wèile, « pour, afin de »)} : Schéma \zh{为了} + But visé (verbe ou nom d'action) ， + Action principale 。
      \begin{itemize}
      \item Le but se place \textbf{en tête de phrase} (antéposition), contrairement au français « pour + infinitif » souvent en fin.
      \item \zh{为了学好汉语，我每天练习。} \emph{Wèile xué hǎo Hànyǔ, wǒ měitiān liànxí.} « Pour bien apprendre le chinois, je m'entraîne chaque jour. »
      \end{itemize}
\item \textbf{Distinction cruciale} :
      \begin{itemize}
      \item \zh{为了} + [Action/But] = \textbf{Finalité} (Pourquoi ?). Ex: \zh{为了健康，他跑步。}
      \item \zh{给} (gěi) / \zh{为} (wèi) + [Personne] = \textbf{Bénéficiaire} (Pour qui ?). Ex: \zh{我给妈妈买花。} « J'achète des fleurs \textbf{pour} ma mère. »
      \end{itemize}
\item \textbf{Coordination de deux actions simultanées} : \cle{一边……一边……} (yìbiān… yìbiān…). Même sujet, deux verbes d'action.
      \zh{他一边吃饭一边看电视。} \emph{Tā yìbiān chīfàn yìbiān kàn diànshì.} « Il mange en regardant la télé. »
\item \textbf{Coordination de deux qualités/états} : \cle{又……又……} (yòu… yòu…). « À la fois... et... ». Sujet + \zh{又} + Adj 1 + \zh{又} + Adj 2.
      \zh{这个菜又好吃又便宜。} \emph{Zhège cài yòu hǎochī yòu piányi.} « Ce plat est à la fois bon et pas cher. »
\item \textbf{Coordination de noms (seulement)} : \cle{和} (hé) ou \zh{跟} (gēn) / \zh{与} (yǔ). Ne relie \textbf{jamais} deux verbes ou deux phrases.
      \zh{我喜欢足球和篮球。} \emph{Wǒ xǐhuan zúqiú hé lánqiú.}
\end{enumerate}
\end{methode}

\begin{exoresolu}[Composition guidée : but et coordination]
Énoncé — Traduis en chinois :\\
1. « Pour trouver un bon travail, il apprend le chinois et l'anglais. »\\
2. « Elle écoute de la musique en faisant ses devoirs. »\\
3. « Ce marché est à la fois grand et animé. »
\tcblower
Solution détaillée :
\begin{enumerate}
\item But en tête : \zh{为了找到好工作} ; Coordination noms : \zh{汉语和英语}.\\
\zh{为了找到好工作，他学习汉语和英语。} \emph{Wèile zhǎodào hǎo gōngzuò, tā xuéxí Hànyǔ hé Yīngyǔ.}
\item Actions simultanées (même sujet \zh{她}) $\rightarrow$ \zh{一边……一边……}.\\
\zh{她一边做作业一边听音乐。} \emph{Tā yìbiān zuò zuòyè yìbiān tīng yīnyuè.}
\item Qualités simultanées $\rightarrow$ \zh{又……又……} ; « animé » = \zh{热闹} (rènao).\\
\zh{这个市场又大又热闹。} \emph{Zhège shìchǎng yòu dà yòu rènao.}
\end{enumerate}
\textbf{Conclusion} : Bien choisir l'outil selon la nature syntaxique : \zh{和} (noms), \zh{又……又……} (adjectifs/états), \zh{一边……一边……} (actions simultanées), \zh{为了} (finalité).
\end{exoresolu}

\courssub{Fiche méthode 5 : Transformer et combiner les connecteurs dans un court texte}

\begin{methode}[Rédiger un court texte avec phrases complexes (Niveau Bac)]
Pour une production écrite de type Bac (80-100 caractères, 8-10 lignes manuscrites) :
\begin{enumerate}
\item \textbf{Planifier} : 1 phrase d'accroche simple, puis \textbf{au moins 3 phrases complexes variées} (cause, concession, hypothèse, but), 1 phrase de conclusion.
\item \textbf{Varier les connecteurs} : Interdiction de répéter 3 fois \zh{因为……所以……}. Alterne : \zh{虽然……但是……}, \zh{如果……就……}, \zh{为了}, \zh{既然……就……} (HSK 4), \zh{只有……才……} (HSK 4).
\item \textbf{Vérifier chaque paire} : \zh{因为} $\leftrightarrow$ \zh{所以} ; \zh{虽然} $\leftrightarrow$ \zh{但是} ; \zh{如果} $\leftrightarrow$ \zh{就}. Un ouvrant sans fermant = faute lourde.
\item \textbf{Contrôler la place du sujet} : Commun = avant le 1er connecteur ; Différents = dans chaque proposition.
\item \textbf{Relire (Check-list Bac)} :
      \begin{itemize}
      \item Caractères exacts (\zh{因为} $\neq$ \zh{应为} ; \zh{所以} $\neq$ \zh{所已} ; \zh{虽然} $\neq$ \zh{虽燃}).
      \item Ponctuation chinoise pleine chasse ： ， 。 ？ ！
      \item Classificateurs présents (\zh{一本}, \zh{两个}, \zh{三个小时}).
      \item Tons du pinyin si demandés.
      \end{itemize}
\end{enumerate}
\end{methode}

\begin{exoresolu}[Production écrite : Mon projet d'avenir]
Énoncé — Rédige un court texte (6-8 phrases) sur « Pourquoi j'apprends le chinois », en utilisant \textbf{au moins une fois chacune} des structures : \zh{因为……所以……}, \zh{虽然……但是……}, \zh{如果……就……}, \zh{为了}.
\tcblower
Solution détaillée — Démarche et modèle rédigé :\\
\textbf{Étape 1 (Accroche)} : \zh{我是雅温得一所中学的学生。} \emph{Wǒ shì Yǎwēndé yì suǒ zhōngxué de xuésheng.} « Je suis élève dans un lycée de Yaoundé. »\\
\textbf{Étape 2 (Cause)} : \zh{因为中国是一个很重要的国家，所以我想学习汉语。} \emph{Yīnwèi Zhōngguó shì yí gè hěn zhòngyào de guójiā, suǒyǐ wǒ xiǎng xuéxí Hànyǔ.} « Parce que la Chine est un pays très important, je veux apprendre le chinois. »\\
\textbf{Étape 3 (Concession)} : \zh{虽然汉字很难写，但是我觉得很有意思。} \emph{Suīrán Hànzì hěn nán xiě, dànshì wǒ juéde hěn yǒu yìsi.} « Bien que les caractères soient difficiles à écrire, je trouve cela très intéressant. »\\
\textbf{Étape 4 (But)} : \zh{为了将来找到好工作，我每天学习两个小时。} \emph{Wèile jiānglái zhǎodào hǎo gōngzuò, wǒ měitiān xuéxí liǎng gè xiǎoshí.} « Pour trouver un bon travail plus tard, j'étudie deux heures par jour. »\\
\textbf{Étape 5 (Hypothèse)} : \zh{如果我的汉语很好，我就可以去中国学习。} \emph{Rúguǒ wǒ de Hànyǔ hěn hǎo, wǒ jiù kěyǐ qù Zhōngguó xuéxí.} « Si mon chinois est très bon, je pourrai aller étudier en Chine. »\\
\textbf{Étape 6 (Conclusion)} : \zh{学习汉语给我打开了很多门。} \emph{Xuéxí Hànyǔ gěi wǒ dǎkāi le hěn duō mén.} « Apprendre le chinois m'a ouvert beaucoup de portes. »\\
\textbf{Vérification} : Les 4 structures obligatoires sont présentes. Paires complètes. Sujets bien placés. Ponctuation chinoise. Vocabulaire HSK 3 (重要, 将来, 打开). Texte cohérent, niveau Terminale A validé.
\end{exoresolu}

\begin{popfigure}
\begin{tikzpicture}[
    mindmap,
    grow cyclic,
    every node/.style={concept, execute at begin node=\hskip0pt},
    root concept/.append style={concept color=popblue, minimum size=2.5cm, text=white, font=\bfseries\large},
    level 1 concept/.append style={minimum size=2cm, sibling angle=90, font=\bfseries\normalsize},
    level 2 concept/.append style={minimum size=1.6cm, sibling angle=45, font=\small},
    concept color=popblue!60,
    text=white
]
\node[align=center]{Phrase complexe\\ (复句)}
    child[concept color=poporange] { node {Cause $\rightarrow$ Conséquence}
        child { node {\zh{因为……所以……}} }
        child { node {\zh{由于……因此……} (formel)} }
    }
    child[concept color=popgreen] { node {Concession}
        child { node {\zh{虽然……但是……}} }
        child { node {\zh{尽管……也……}} }
    }
    child[concept color=poppurple] { node {Condition $\rightarrow$ Conséquence}
        child { node {\zh{如果……就……}} }
        child { node {\zh{要是……就……} (oral)} }
        child { node {\zh{只有……才……} (nécessaire)} }
    }
    child[concept color=poppink] { node {But \& Coordination}
        child { node {\zh{为了……} (But)} }
        child { node {\zh{和} (Noms)} }
        child { node {\zh{又……又……} (Qualités)} }
        child { node {\zh{一边……一边……} (Actions)} }
    };
\end{tikzpicture}
\legende{Carte mentale des connecteurs logiques majeurs (Niveau Terminale / HSK 3-4)}
\end{popfigure}

\begin{attention}[Les 5 erreurs qui coûtent des points au Bac]
\begin{enumerate}
\item \textbf{Calque français « Bien que... mais... » tronqué} : Écrire \zh{虽然他很累，他去工作。} \textbf{SANS} \zh{但是}. \textbf{Correction} : La paire \zh{虽然……但是……} est indissociable à l'écrit. N'applique pas la règle française qui interdit la redondance.
\item \textbf{Mauvaise place de \zh{就}} : Écrire \zh{如果下雨，就我不去。} \textbf{FAUX}. \zh{就} se place \textbf{après le sujet, avant le verbe/auxiliaire/négation} : \zh{如果下雨，我就不去。} C'est le point de grammaire n°1 sanctionné.
\item \textbf{Caractères homophones confondus (Zéro tolérance)} :
      \begin{itemize}
      \item \zh{因为} (yīnwèi) $\neq$ \zh{应为} (yīngwèi, « devrait être »).
      \item \zh{所以} (suǒyǐ) $\neq$ \zh{所已} (sans sens).
      \item \zh{虽然} (suīrán) $\neq$ \zh{虽燃} (suī rán, « bien que brûler »).
      \item \zh{为了} (wèile) $\neq$ \zh{为勒} (wèi lēi).
      \end{itemize}
      Vérifie systématiquement l'écriture des connecteurs, ils sont notés « Caractères » à part entière.
\item \textbf{Ordre des propositions inversé (Calque syntaxique)} : Le français dit « Je reste \textbf{parce qu'il pleut} ». Le chinois \textbf{exige} : \zh{因为下雨，所以我在家。} (Cause d'abord). Ne traduis jamais dans l'ordre linéaire français.
\item \textbf{Oubli du classificateur dans la conséquence} : Toute phrase complexe contenant un nom dénombrable exige son classificateur.
      \begin{itemize}
      \item \textcolor{red}{\zh{我买一词典}} $\rightarrow$ \textbf{Correct} : \zh{我买\cle{一本}词典。}
      \item \textcolor{red}{\zh{我们学两个小时}} $\rightarrow$ \textbf{Correct} : \zh{我们学\cle{两个}小时。} (ou \zh{两个小时}).
      \item \textcolor{red}{\zh{他有三个朋友}} $\rightarrow$ \textbf{Correct} : \zh{他有\cle{三个}朋友。}
      \end{itemize}
      L'oubli du classificateur est une faute de grammaire notée à part, même si la structure complexe est correcte.
\end{enumerate}
\end{attention}

\courssub{Entraînement : Exercices d'application (Devoir / Classe)}

\begin{exercice}[Exercices de la Leçon 22]
\textbf{Exercice 1 : Transformation (Cause / Conséquence).} Relie les phrases par \zh{因为……所以……}. Traduis.
\begin{enumerate}
\item \zh{昨天下大雨。足球比赛取消了。} \emph{Zuótiān xià dà yǔ. Zúqiú bǐsài qǔxiāo le.}
\item \zh{他不喜欢吃辣。他不去四川餐馆。} \emph{Tā bù xǐhuan chī là. Tā bù qù Sìchuān cānguǎn.}
\end{enumerate}

\textbf{Exercice 2 : Complétion (Concession).} Complète avec \zh{虽然} / \zh{但是} / \zh{但} / \zh{可是}. Traduis.
\begin{enumerate}
\item \zh{……他很年轻，……他经验很丰富。} (Il est jeune mais expérimenté)
\item \zh{……今天很累，……我还要复习功课。} (Bien que fatigué, je dois réviser)
\end{enumerate}

\textbf{Exercice 3 : Condition / Hypothèse (Le piège du \zh{就}).} Traduis en chinois. Place \zh{就} correctement.
\begin{enumerate}
\item « Si tu as le temps, viens chez moi. » (\zh{来} lái)
\item « S'il ne vient pas, j'irai seul. » (\zh{自己} zìjǐ, \zh{去} qù)
\item « Si j'étais toi, j'accepterais. » (Irréel : \zh{如果我是你，我就……} ; \zh{答应} dāying)
\end{enumerate}

\textbf{Exercice 4 : But et Coordination.} Traduis en chinois.
\begin{enumerate}
\item « Pour passer l'examen HSK 4, elle révise la grammaire et le vocabulaire. » (\zh{复习} fùxí, \zh{语法} yǔfǎ, \zh{词汇} cíhuì)
\item « Il parle et rit en même temps. » (\zh{说话} shuōhuà, \zh{笑} xiào)
\item « Ce livre est à la fois intéressant et utile. » (\zh{有用} yǒuyòng)
\end{enumerate}

\textbf{Exercice 5 : Correction d'erreurs (Niveau Bac).} Corrige les phrases suivantes (caractères, ordre, mots-outils, classificateurs).
\begin{enumerate}
\item \zh{应为他生病，所已他没来上课。}
\item \zh{虽然下雨，他没带伞。}
\item \zh{如果明天有空，就我们去爬山。}
\item \zh{为了买一新衣服，她去了商场。}
\item \zh{他一边开车一边打电话，很危险。} (Vérifie la structure \zh{一边}...)
\end{enumerate}

\textbf{Exercice 6 : Production écrite (Type Bac).} Rédige un paragraphe (8-10 phrases) sur le thème : \zh{我的大学计划} (Mon projet universitaire). Contraintes : Utilise \textbf{au moins} 1x \zh{因为……所以……}, 1x \zh{虽然……但是……}, 1x \zh{如果……就……}, 1x \zh{为了}, 1x \zh{又……又……} ou \zh{一边……一边……}. Soigne les classificateurs et la ponctuation.
\end{exercice}

\begin{corrige}[Exercices de la Leçon 22]
\textbf{Exercice 1 :}
\begin{enumerate}
\item \zh{因为昨天下大雨，所以足球比赛取消了。} \emph{Yīnwèi zuótiān xià dà yǔ, suǒyǐ zúqiú bǐsài qǔxiāo le.} « Parce qu'il a beaucoup plu hier, le match de foot a été annulé. »
\item \zh{他因为不喜欢吃辣，所以不去四川餐馆。} \emph{Tā yīnwèi bù xǐhuan chī là, suǒyǐ bù qù Sìchuān cānguǎn.} « Comme il n'aime pas manger épicé, il ne va pas au resto Sichuan. »
\end{enumerate}

\textbf{Exercice 2 :}
\begin{enumerate}
\item \zh{虽然他很年轻，但是经验很丰富。} \emph{Suīrán tā hěn niánqīng, dànshì jīngyàn hěn fēngfù.} « Bien qu'il soit jeune, son expérience est très riche. » (Sujet commun \zh{他} avant \zh{虽然}, sujet omis 2e prop).
\item \zh{虽然今天很累，但我还要复习功课。} \emph{Suīrán jīntiān hěn lèi, dàn wǒ hái yào fùxí gōngkè.} « Bien qu'aujourd'hui je sois très fatigué, je dois encore réviser mes leçons. »
\end{enumerate}

\textbf{Exercice 3 :}
\begin{enumerate}
\item \zh{如果你有时间，就来我家吧。} \emph{Rúguǒ nǐ yǒu shíjiān, jiù lái wǒ jiā ba.} (\zh{就} avant \zh{来}).
\item \zh{如果他不来，我就自己去。} \emph{Rúguǒ tā bù lái, wǒ jiù zìjǐ qù.} (\zh{就} avant \zh{自己去}).
\item \zh{如果我是你，我就答应。} \emph{Rúguǒ wǒ shì nǐ, wǒ jiù dāying.} (Irréel : structure inchangée).
\end{enumerate}

\textbf{Exercice 4 :}
\begin{enumerate}
\item \zh{为了通过HSK四级考试，她复习语法和词汇。} \emph{Wèile tōngguò HSK sìjí kǎoshì, tā fùxí yǔfǎ hé cíhuì.}
\item \zh{他一边说话一边笑。} \emph{Tā yìbiān shuōhuà yìbiān xiào.}
\item \zh{这本书又有趣又有用。} \emph{Zhè běn shū yòu yǒuqù yòu yǒuyòng.} (Classificateur \zh{本} pour livre).
\end{enumerate}

\textbf{Exercice 5 : Correction d'erreurs}
\begin{enumerate}
\item \zh{因为他生病，所以他没来上课。} (Correction : \zh{应为} $\rightarrow$ \zh{因为} ; \zh{所已} $\rightarrow$ \zh{所以}).
\item \zh{虽然下雨，但是他没带伞。} (Ajout obligatoire de \zh{但是} / \zh{但} / \zh{可是}).
\item \zh{如果明天有空，我们就去爬山。} (\zh{就} placé après sujet \zh{我们}, avant \zh{去}).
\item \zh{为了买\cle{一件}新衣服，她去了商场。} (Ajout classificateur \zh{件} jiàn pour vêtement \zh{衣服}).
\item \zh{他一边开车一边打电话，很危险。} \textbf{Phrase correcte.} (Structure \zh{一边……一边……} bien utilisée, sujet commun \zh{他}).
\end{enumerate}

\textbf{Exercice 6 : Production écrite (Modèle de correction)}
\zh{我的大学计划}\\
\zh{我是喀麦隆的一名高三学生。} \emph{Wǒ shì Kāmàilóng de yì míng gāosān xuésheng.}\\
\zh{因为我想当外交官，所以我必须学好外语。} \emph{Yīnwèi wǒ xiǎng dāng wàijiāoguān, suǒyǐ wǒ bìxū xué hǎo wàiyǔ.}\\
\zh{虽然法语是我的母语，但是我觉得中文更有挑战性。} \emph{Suīrán Fǎyǔ shì wǒ de mǔyǔ, dànshì wǒ juéde Zhōngwén gèng yǒu tiǎozhànxìng.}\\
\zh{为了考上北京外国语大学，我每天复习三个小时。} \emph{Wèile kǎo shàng Běijīng Wàiguóyǔ Dàxué, wǒ měitiān fùxí sān gè xiǎoshí.}\\
\zh{如果我的HSK成绩很好，我就可以申请奖学金。} \emph{Rúguǒ wǒ de HSK chéngjī hěn hǎo, wǒ jiù kěyǐ shēnqǐng jiǎngxuéjīn.}\\
\zh{学习中文又累又快乐，因为我能认识新朋友。} \emph{Xuéxí Zhōngwén yòu lèi yòu kuàilè, yīnwèi wǒ néng rènshi xīn péngyou.}\\
\zh{我一边准备高考，一边练习汉字书写。} \emph{Wǒ yìbiān zhǔnbèi gāokǎo, yìbiān liànxí Hànzì shūxiě.}\\
\zh{我相信努力就会有结果。} \emph{Wǒ xiāngxìn nǔlì jiù huì yǒu jiéguǒ.}\\
\textbf{Analyse} : 6 phrases complexes variées, 4 structures obligatoires + coordination, classificateurs (\zh{一名}, \zh{三个}, \zh{个小时}), ponctuation correcte, vocabulaire HSK 3-4 (\zh{外交官}, \zh{挑战性}, \zh{申请}, \zh{奖学金}, \zh{相信}). Note attendue : 18-20/20.
\end{corrige}

\newpage

\coursec{Exercices}

\courssub{Je vérifie mes connaissances}

\begin{exercice}[QCM : le bon corrélatif]
Pour chaque phrase, choisis la bonne réponse (A, B ou C).

\begin{enumerate}
\item \zh{＿＿＿明天下雨，我们＿＿＿不去踢足球。}\par
A. \zh{虽然……但是……} \quad B. \zh{如果……就……} \quad C. \zh{为了……所以……}
\item \zh{＿＿＿他生病了，＿＿＿他没有来上课。}\par
A. \zh{因为……所以……} \quad B. \zh{如果……就……} \quad C. \zh{虽然……但是……}
\item \zh{＿＿＿学好中文，我每天练习汉字。}\par
A. \zh{因为} \quad B. \zh{虽然} \quad C. \zh{为了}
\item \zh{＿＿＿中文字很难，＿＿＿我觉得很有意思。}\par
A. \zh{如果……就……} \quad B. \zh{虽然……但是……} \quad C. \zh{因为……所以……}
\end{enumerate}
\end{exercice}

\begin{exercice}[Vrai ou faux ? Justifie ta réponse]
Dis si les affirmations suivantes sont vraies (V) ou fausses (F), puis justifie en une phrase en français.

\begin{enumerate}
\item Dans une phrase avec \zh{如果……就……}, le mot \zh{就} se place avant le sujet de la deuxième proposition.
\item On peut utiliser \zh{因为} et \zh{所以} ensemble dans la même phrase en chinois.
\item \zh{为了} exprime le but, tandis que \zh{因为} exprime la cause.
\item En chinois, contrairement au français, on peut dire à la fois \zh{虽然} et \zh{但是} dans la même phrase.
\end{enumerate}
\end{exercice}

\begin{exercice}[Vocabulaire : complète le tableau]
Recopie et complète le tableau avec le pinyin et la traduction française.

\begin{center}
\begin{tabularx}{\linewidth}{|X|X|X|X|}
\hline
\rowcolor{popblueL} Caractères & Pinyin & Nature & Traduction \\
\hline
\zh{因为} & & conjonction & \\
\hline
\zh{所以} & & conjonction & \\
\hline
\zh{虽然} & & conjonction & \\
\hline
\zh{但是} & & conjonction & \\
\hline
\zh{如果} & & conjonction & \\
\hline
\zh{为了} & & préposition & \\
\hline
\end{tabularx}
\end{center}
\end{exercice}

\begin{exercice}[Associe le pinyin aux caractères]
Associe chaque mot en pinyin (1--6) au bon mot en caractères (A--F).

\begin{enumerate}
\item \emph{suīrán} \qquad 2. \emph{wèile} \qquad 3. \emph{yīnwèi}
\item \emph{rúguǒ} \qquad 5. \emph{suǒyǐ} \qquad 6. \emph{dànshì}
\end{enumerate}

\begin{itemize}
\item[A.] \zh{但是} \qquad B. \zh{因为} \qquad C. \zh{如果}
\item[D.] \zh{虽然} \qquad E. \zh{所以} \qquad F. \zh{为了}
\end{itemize}
\end{exercice}

\courssub{Je m'entraîne}

\begin{exercice}[Relie les deux phrases]
Relie chaque paire de phrases simples en une phrase complexe, en utilisant le corrélatif indiqué entre parenthèses. Écris la phrase complète en caractères.

\begin{enumerate}
\item \zh{他生病了。他没来上课。} (\zh{因为……所以……})
\item \zh{他很累。他还学习。} (\zh{虽然……但是……})
\item \zh{明天下雨。我们不去杜阿拉。} (\zh{如果……就……})
\item \zh{我想去中国留学。我努力学习中文。} (\zh{为了})
\end{enumerate}
\end{exercice}

\begin{exercice}[Texte à trous : les mots-outils]
Complète le texte suivant avec les mots de la liste : \zh{因为、所以、虽然、但是、如果、就、为了}。

\begin{textebox}[À la gare de Yaoundé]
\zh{＿＿＿我想去中国学习，＿＿＿我每天学中文。＿＿＿汉字很难，＿＿＿我很喜欢写。＿＿＿明天下雨，我＿＿＿在家练习。＿＿＿通过考试，我要更加努力。}\par
\emph{...wǒ xiǎng qù Zhōngguó xuéxí, ...wǒ měitiān xué Zhōngwén. ...Hànzì hěn nán, ...wǒ hěn xǐhuan xiě. ...míngtiān xià yǔ, wǒ...zài jiā liànxí. ...tōngguò kǎoshì, wǒ yào gèngjiā nǔlì.}\par
« Parce que je veux aller étudier en Chine, j'apprends le chinois tous les jours. Bien que les caractères soient difficiles, j'aime beaucoup les écrire. S'il pleut demain, je m'entraînerai à la maison. Pour réussir l'examen, je dois travailler encore plus. »
\end{textebox}
\end{exercice}

\begin{exercice}[Remets les mots dans l'ordre]
Remets les mots dans le bon ordre pour former une phrase complexe correcte. Écris la phrase complète.

\begin{enumerate}
\item \zh{就 / 如果 / 去 / 你 / 我 / 雅温得 / 也 / 去}
\item \zh{但是 / 虽然 / 喜欢 / 很难 / 我 / 中文 / 很}
\item \zh{所以 / 因为 / 没来 / 他 / 上课 / 生病了}
\item \zh{为了 / 每天 / 中文 / 练习 / 学好 / 他}
\end{enumerate}
\end{exercice}

\begin{exercice}[Transformation et correction]
\textbf{A.} Transforme chaque phrase complexe en deux phrases simples, comme dans l'exemple.\par
\emph{Exemple :} \zh{虽然中文很难，但是我很喜欢。} → \zh{中文很难。我很喜欢中文。}

\begin{enumerate}
\item \zh{虽然他生病了，但是他还是来上课了。}
\item \zh{因为今天下雨，所以我们不去市场。}
\end{enumerate}

\textbf{B.} Chaque phrase contient une erreur. Relève-la et écris la phrase corrigée.

\begin{enumerate}
\item \zh{如果明天下雨，就不我去学校。}
\item \zh{因为他很累，但是他睡觉了。}
\item \zh{为了我生病了，所以我不去上课。}
\end{enumerate}
\end{exercice}

\courssub{Je me prépare au Bac}

\begin{exercice}[Compréhension de texte et questions de langue]
Lis le texte suivant, puis réponds aux questions.

\begin{textebox}[Mon rêve : la Chine]
\zh{我叫恩东，是雅温得的高中生。因为我喜欢中国文化，所以我学习中文已经两年了。虽然汉字很难，但是我每天练习写二十个汉字。我的老师说，如果我想去中国留学，就要更加努力。为了通过HSK考试，我每天早上听中文广播。要是我能去中国，我就去参观长城。虽然我的家离中国很远，但是我相信我的梦想会实现。}\par
\emph{Wǒ jiào Ēndōng, shì Yǎwēndé de gāozhōngshēng. Yīnwèi wǒ xǐhuan Zhōngguó wénhuà, suǒyǐ wǒ xuéxí Zhōngwén yǐjīng liǎng nián le. Suīrán Hànzì hěn nán, dànshì wǒ měitiān liànxí xiě èrshí ge Hànzì. Wǒ de lǎoshī shuō, rúguǒ wǒ xiǎng qù Zhōngguó liúxué, jiù yào gèngjiā nǔlì. Wèile tōngguò HSK kǎoshì, wǒ měitiān zǎoshang tīng Zhōngwén guǎngbō. Yàoshi wǒ néng qù Zhōngguó, wǒ jiù qù cānguān Chángchéng. Suīrán wǒ de jiā lí Zhōngguó hěn yuǎn, dànshì wǒ xiāngxìn wǒ de mèngxiǎng huì shíxiàn.}\par
« Je m'appelle Ndong, je suis lycéen à Yaoundé. Parce que j'aime la culture chinoise, j'apprends le chinois depuis deux ans. Bien que les caractères soient difficiles, je m'entraîne à en écrire vingt chaque jour. Mon professeur dit que si je veux aller étudier en Chine, je dois travailler encore plus. Pour réussir l'examen HSK, j'écoute la radio chinoise tous les matins. Si je peux aller en Chine, je visiterai la Grande Muraille. Bien que ma maison soit loin de la Chine, je crois que mon rêve se réalisera. »
\end{textebox}

\textbf{Questions de compréhension} (réponds en français) :
\begin{enumerate}
\item Pourquoi Ndong apprend-il le chinois ?
\item Que fait-il chaque jour pour progresser ? Cite deux actions.
\item Que fera-t-il s'il peut aller en Chine ?
\end{enumerate}

\textbf{Questions de langue} :
\begin{enumerate}
\item Relève dans le texte trois structures complexes différentes et recopie les phrases qui les contiennent.
\item Dans la phrase \zh{如果我想去中国留学，就要更加努力}, pourquoi \zh{就} est-il placé avant \zh{要} ?
\item Transforme en deux phrases simples : \zh{虽然汉字很难，但是我每天练习写二十个汉字。}
\end{enumerate}
\end{exercice}

\begin{exercice}[Thème et version]
\textbf{A. Thème.} Traduis en chinois (caractères) les phrases suivantes.

\begin{enumerate}
\item « Si tu viens à Yaoundé, je te montrerai ma ville. » (\zh{如果……就……})
\item « Parce qu'il pleut, nous n'allons pas au marché. » (\zh{因为……所以……})
\item « Bien que le chinois soit difficile, je l'aime beaucoup. » (\zh{虽然……但是……})
\item « Pour réussir au baccalauréat, elle travaille tous les jours. » (\zh{为了})
\end{enumerate}

\textbf{B. Version.} Traduis en français le texte suivant.

\begin{textebox}[Version]
\zh{虽然我家离学校很远，但是我从不迟到。因为我想学好中文，所以我每天六点起床。如果今年考试成功，我就去中国学习。}\par
\emph{Suīrán wǒ jiā lí xuéxiào hěn yuǎn, dànshì wǒ cóng bù chídào. Yīnwèi wǒ xiǎng xué hǎo Zhōngwén, suǒyǐ wǒ měitiān liù diǎn qǐchuáng. Rúguǒ jīnnián kǎoshì chénggōng, wǒ jiù qù Zhōngguó xuéxí.}
\end{textebox}
\end{exercice}

\begin{exercice}[Expression écrite : Pourquoi j'apprends le chinois]
Rédige en chinois un paragraphe de \textbf{5 à 7 phrases} (environ 80 à 100 caractères) sur le sujet : « Pourquoi j'apprends le chinois » (\zh{我为什么学习中文}).

\textbf{Consignes obligatoires :}
\begin{itemize}
\item Utilise \textbf{au moins trois structures complexes différentes} parmi : \zh{因为……所以……}, \zh{虽然……但是……}, \zh{如果……就……}, \zh{要是……就……}, \zh{为了……}.
\item Écris uniquement en caractères chinois (pas de pinyin).
\item Respecte l'ordre des mots et la place de \zh{就} avant le verbe.
\item Termine par une phrase qui exprime ton projet d'avenir.
\end{itemize}

\textbf{Éléments de vocabulaire utiles :} \zh{文化} (culture), \zh{梦想} (rêve), \zh{留学} (étudier à l'étranger), \zh{努力} (faire des efforts), \zh{将来} (à l'avenir), \zh{工作} (travailler).
\end{exercice}

\newpage

\coursec{Corrigés des exercices}

\begin{corrige}[Exercice 1 — QCM : le bon corrélatif]
\begin{enumerate}
\item \textbf{B.} \zh{如果明天下雨，我们就不去踢足球。}\par
\emph{Rúguǒ míngtiān xià yǔ, wǒmen jiù bù qù tī zúqiú.} « S'il pleut demain, nous n'irons pas jouer au football. »\par
Justification : il s'agit d'une \cle{hypothèse} (il n'est pas sûr qu'il pleuve) ; seul \zh{如果……就……} convient. \zh{为了……所以……} n'existe pas : \zh{为了} (but) n'a pas de corrélatif dans la seconde proposition.
\item \textbf{A.} \zh{因为他生病了，所以他没有来上课。}\par
\emph{Yīnwèi tā shēngbìng le, suǒyǐ tā méiyǒu lái shàngkè.} « Parce qu'il était malade, il n'est pas venu en cours. »\par
Justification : la maladie est la \cle{cause} réelle de l'absence → \zh{因为……所以……}.
\item \textbf{C.} \zh{为了学好中文，我每天练习汉字。}\par
\emph{Wèile xué hǎo Zhōngwén, wǒ měitiān liànxí Hànzì.} « Pour bien apprendre le chinois, je m'entraîne aux caractères tous les jours. »\par
Justification : « bien apprendre le chinois » est le \cle{but} de l'entraînement quotidien → \zh{为了}.
\item \textbf{B.} \zh{虽然中文字很难，但是我觉得很有意思。}\par
\emph{Suīrán Zhōngwén zì hěn nán, dànshì wǒ juéde hěn yǒu yìsi.} « Bien que les caractères chinois soient difficiles, je les trouve très intéressants. »\par
Justification : opposition entre « difficile » et « intéressant » → \cle{concession} \zh{虽然……但是……}.
\end{enumerate}
\end{corrige}

\begin{corrige}[Exercice 2 — Vrai ou faux ?]
\begin{enumerate}
\item \textbf{Faux.} \zh{就} se place \textbf{après} le sujet de la deuxième proposition et \textbf{avant le verbe} : \zh{如果明天下雨，我就待在家里。} (\emph{wǒ jiù...}), et non \zh{就我}. Si le sujet des deux propositions est identique, il peut être omis dans la seconde : \zh{如果明天下雨，就不去学校。}
\item \textbf{Vrai.} Contrairement au français (on ne dit pas « parce que... donc... »), le chinois utilise volontiers les deux corrélatifs ensemble : \zh{因为他生病了，所以他没来上课。} On peut aussi n'en garder qu'un seul : \zh{因为他生病了，他没来上课。}
\item \textbf{Vrai.} \zh{为了} + nom ou verbe introduit le \cle{but} (« pour, afin de »), \zh{因为} introduit la \cle{cause} (« parce que »). Attention à ne pas les confondre : on ne dit jamais \zh{为了他生病了}.
\item \textbf{Vrai.} En français on dit « bien que... » ou « ...mais... », jamais les deux ; en chinois, \zh{虽然……但是……} s'emploient ensemble dans la même phrase : \zh{虽然很难，但是很有意思。}
\end{enumerate}
\end{corrige}

\begin{corrige}[Exercice 3 — Vocabulaire : complète le tableau]
\begin{center}
\begin{tabularx}{\linewidth}{|X|X|X|X|}
\hline
\rowcolor{popblueL} Caractères & Pinyin & Nature & Traduction \\
\hline
\zh{因为} & \emph{yīnwèi} & conjonction & parce que \\
\hline
\zh{所以} & \emph{suǒyǐ} & conjonction & donc, c'est pourquoi \\
\hline
\zh{虽然} & \emph{suīrán} & conjonction & bien que, quoique \\
\hline
\zh{但是} & \emph{dànshì} & conjonction & mais, cependant \\
\hline
\zh{如果} & \emph{rúguǒ} & conjonction & si (hypothèse) \\
\hline
\zh{为了} & \emph{wèile} & préposition & pour, afin de \\
\hline
\end{tabularx}
\end{center}
\textbf{Point de vigilance :} le caractère \zh{为} a deux prononciations. Il se lit \emph{wèi} (4\up{e} ton) quand il exprime la cause ou le but : \zh{为了} (\emph{wèile}, « pour »), \zh{为什么} (\emph{wèishénme}, « pourquoi »). Il se lit \emph{wéi} (2\up{e} ton) dans d'autres mots comme \zh{认为} (\emph{renwéi}, « penser, estimer ») ou \zh{成为} (\emph{chéngwéi}, « devenir »).
\end{corrige}

\begin{corrige}[Exercice 4 — Associe le pinyin aux caractères]
\begin{center}
\begin{tabularx}{\linewidth}{|X|X|X|}
\hline
\rowcolor{popblueL} Pinyin & Caractères & Traduction \\
\hline
1. \emph{suīrán} & D. \zh{虽然} & bien que \\
\hline
2. \emph{wèile} & F. \zh{为了} & pour, afin de \\
\hline
3. \emph{yīnwèi} & B. \zh{因为} & parce que \\
\hline
4. \emph{rúguǒ} & C. \zh{如果} & si \\
\hline
5. \emph{suǒyǐ} & E. \zh{所以} & donc \\
\hline
6. \emph{dànshì} & A. \zh{但是} & mais \\
\hline
\end{tabularx}
\end{center}
Réponses : 1--D, 2--F, 3--B, 4--C, 5--E, 6--A.\par
\textbf{Astuce mémoire :} \zh{因} contient \zh{大} dans une enceinte \zh{囗} (la cause est « enfermée » dans le fait) ; \zh{果} (clé du bois \zh{木}) signifie « fruit, résultat » — \zh{如果} introduit donc bien une condition dont dépend le « fruit ».
\end{corrige}

\begin{corrige}[Exercice 5 — Relie les deux phrases]
\begin{enumerate}
\item \zh{因为他生病了，所以他没有来上课。}\par
\emph{Yīnwèi tā shēngbìng le, suǒyǐ tā méiyǒu lái shàngkè.} « Parce qu'il était malade, il n'est pas venu en cours. »
\item \zh{虽然他很累，但是他还学习。}\par
\emph{Suīrán tā hěn lèi, dànshì tā hái xuéxí.} « Bien qu'il soit très fatigué, il étudie encore. »
\item \zh{如果明天下雨，我们就不去杜阿拉。}\par
\emph{Rúguǒ míngtiān xià yǔ, wǒmen jiù bù qù Dùālā.} « S'il pleut demain, nous n'irons pas à Douala. »\par
Point de vigilance : \zh{就} se place après le sujet \zh{我们} et avant \zh{不去}.
\item \zh{为了去中国留学，我努力学习中文。}\par
\emph{Wèile qù Zhōngguó liúxué, wǒ nǔlì xuéxí Zhōngwén.} « Pour aller étudier en Chine, je travaille dur le chinois. »\par
Point de vigilance : avec \zh{为了}, on supprime \zh{我想} ; le but s'exprime directement par \zh{为了} + verbe.
\end{enumerate}
\end{corrige}

\begin{corrige}[Exercice 6 — Texte à trous : les mots-outils]
Texte complété :\par
\zh{因为我想去中国学习，所以我每天学中文。虽然汉字很难，但是我很喜欢写。如果明天下雨，我就在家练习。为了通过考试，我要更加努力。}\par
\emph{Yīnwèi wǒ xiǎng qù Zhōngguó xuéxí, suǒyǐ wǒ měitiān xué Zhōngwén. Suīrán Hànzì hěn nán, dànshì wǒ hěn xǐhuan xiě. Rúguǒ míngtiān xià yǔ, wǒ jiù zài jiā liànxí. Wèile tōngguò kǎoshì, wǒ yào gèngjiā nǔlì.}\par
Dans l'ordre : \zh{因为} → \zh{所以} → \zh{虽然} → \zh{但是} → \zh{如果} → \zh{就} → \zh{为了}.\par
\textbf{Justifications :} la 1\up{re} paire exprime la cause (« je veux aller en Chine » → « j'étudie chaque jour ») ; la 2\up{e} paire exprime la concession (« difficile » mais « j'aime ») ; la 3\up{e} exprime l'hypothèse (« s'il pleut » → \zh{就} avant le verbe \zh{在}) ; la dernière exprime le but (\zh{为了通过考试}).
\end{corrige}

\begin{corrige}[Exercice 7 — Remets les mots dans l'ordre]
\begin{enumerate}
\item \zh{如果你去雅温得，我也去。}\par
\emph{Rúguǒ nǐ qù Yǎwēndé, wǒ yě qù.} « Si tu vas à Yaoundé, j'y vais aussi. »\par
(Variante acceptée avec les deux adverbes : \zh{如果你去雅温得，我就也去。} — \zh{就} puis \zh{也}, dans cet ordre, avant le verbe.)
\item \zh{虽然中文很难，但是我很喜欢。}\par
\emph{Suīrán Zhōngwén hěn nán, dànshì wǒ hěn xǐhuan.} « Bien que le chinois soit difficile, je l'aime beaucoup. »
\item \zh{因为他生病了，所以没来上课。}\par
\emph{Yīnwèi tā shēngbìng le, suǒyǐ méi lái shàngkè.} « Parce qu'il était malade, il n'est pas venu en cours. »\par
(Sujet \zh{他} non répété dans la 2\up{e} proposition : omission correcte et naturelle.)
\item \zh{为了学好中文，他每天练习。}\par
\emph{Wèile xué hǎo Zhōngwén, tā měitiān liànxí.} « Pour bien apprendre le chinois, il s'entraîne tous les jours. »
\end{enumerate}
\end{corrige}

\begin{corrige}[Exercice 8 — Transformation et correction]
\textbf{A. Transformation en deux phrases simples}
\begin{enumerate}
\item \zh{虽然他生病了，但是他还是来上课了。} → \zh{他生病了。他还是来上课了。}\par
\emph{Tā shēngbìng le. Tā háishi lái shàngkè le.} « Il était malade. Il est quand même venu en cours. »
\item \zh{因为今天下雨，所以我们不去市场。} → \zh{今天下雨。我们不去市场。}\par
\emph{Jīntiān xià yǔ. Wǒmen bù qù shìchǎng.} « Il pleut aujourd'hui. Nous n'allons pas au marché. »
\end{enumerate}
\textbf{B. Correction des erreurs}
\begin{enumerate}
\item Erreur : place de \zh{就} et de la négation. \zh{就} doit précéder le verbe (et donc la négation \zh{不}), et le sujet \zh{我} doit précéder \zh{就}.\par
Correction : \zh{如果明天下雨，我就不去学校。}\par
\emph{Rúguǒ míngtiān xià yǔ, wǒ jiù bù qù xuéxiào.} « S'il pleut demain, je n'irai pas à l'école. »
\item Erreur : mélange de corrélatifs incompatibles — \zh{因为} (cause) ne se combine pas avec \zh{但是} (concession) ; le corrélatif de \zh{因为} est \zh{所以}.\par
Correction : \zh{因为他很累，所以他睡觉了。}\par
\emph{Yīnwèi tā hěn lèi, suǒyǐ tā shuìjiào le.} « Parce qu'il était très fatigué, il s'est couché. »
\item Erreur : confusion but/cause — « être malade » est une cause, pas un but ; on n'utilise pas \zh{为了}.\par
Correction : \zh{因为我生病了，所以我不去上课。}\par
\emph{Yīnwèi wǒ shēngbìng le, suǒyǐ wǒ bù qù shàngkè.} « Parce que je suis malade, je ne vais pas en cours. »
\end{enumerate}
\end{corrige}

\begin{corrige}[Exercice 9 — Compréhension de texte et questions de langue]
\textbf{Barème indicatif (sur 10) :} compréhension 3 × 1 pt ; questions de langue : 3 pts (relevé) + 2 pts (place de \zh{就}) + 2 pts (transformation).

\textbf{Questions de compréhension}
\begin{enumerate}
\item Ndong apprend le chinois \textbf{parce qu'il aime la culture chinoise} (\zh{因为我喜欢中国文化，所以我学习中文}).
\item Chaque jour, il \textbf{s'entraîne à écrire vingt caractères} (\zh{每天练习写二十个汉字}) et il \textbf{écoute la radio chinoise tous les matins} (\zh{每天早上听中文广播}).
\item S'il peut aller en Chine, il \textbf{visitera la Grande Muraille} (\zh{我就去参观长城}).
\end{enumerate}

\textbf{Questions de langue}
\begin{enumerate}
\item Trois structures complexes relevées dans le texte :
\begin{itemize}
\item Cause : \zh{因为我喜欢中国文化，所以我学习中文已经两年了。}
\item Concession : \zh{虽然汉字很难，但是我每天练习写二十个汉字。}
\item Hypothèse : \zh{如果我想去中国留学，就要更加努力。} (ou \zh{要是我能去中国，我就去参观长城。})
\item But (également accepté) : \zh{为了通过HSK考试，我每天早上听中文广播。}
\end{itemize}
\item Dans \zh{如果我想去中国留学，就要更加努力}, le sujet de la deuxième proposition (\zh{我}) est le même que celui de la première : il est donc \textbf{omis}. \zh{就} se place alors directement \textbf{avant le verbe} (ici le verbe modal \zh{要} « devoir »). La règle reste : \zh{就} = après le sujet (exprimé ou sous-entendu), avant le verbe.
\item \zh{虽然汉字很难，但是我每天练习写二十个汉字。} → \zh{汉字很难。我每天练习写二十个汉字。}\par
\emph{Hànzì hěn nán. Wǒ měitiān liànxí xiě èrshí ge Hànzì.} « Les caractères sont difficiles. Je m'entraîne à en écrire vingt chaque jour. »
\end{enumerate}
\textbf{Points de vigilance :} ne pas traduire \zh{已经两年了} par « déjà deux ans » au futur — c'est une durée accomplie (« depuis deux ans ») ; \zh{要是……就……} est un équivalent plus oral de \zh{如果……就……}.
\end{corrige}

\begin{corrige}[Exercice 10 — Thème et version]
\textbf{Barème indicatif (sur 10) :} thème 4 × 1,5 pt (structure correcte 0,5 + ordre des mots 0,5 + caractères 0,5) ; version 4 pts (fidélité 2 + qualité du français 2).

\textbf{A. Thème}
\begin{enumerate}
\item \zh{如果你来雅温得，我就给你看看我的城市。}\par
\emph{Rúguǒ nǐ lái Yǎwēndé, wǒ jiù gěi nǐ kànkan wǒ de chéngshì.}\par
(Variante : \zh{我就带你参观我的城市。} « je te ferai visiter ma ville ».)
\item \zh{因为下雨，所以我们不去市场。}\par
\emph{Yīnwèi xià yǔ, suǒyǐ wǒmen bù qù shìchǎng.}\par
(Variante : \zh{因为今天下雨，所以我们不去市场。})
\item \zh{虽然中文很难，但是我很喜欢。}\par
\emph{Suīrán Zhōngwén hěn nán, dànshì wǒ hěn xǐhuan.}
\item \zh{为了通过高考，她每天学习。}\par
\emph{Wèile tōngguò gāokǎo, tā měitiān xuéxí.}\par
(Variante : \zh{为了考上好大学，她每天努力学习。} « Pour entrer dans une bonne université, elle travaille dur chaque jour. »)
\end{enumerate}
\textbf{Points de vigilance :} en 1, \zh{就} après le sujet \zh{我} et avant le verbe ; en 2, ne pas écrire \zh{因为……但是……} ; en 4, \zh{为了} + verbe directement, sans \zh{她想}.

\textbf{B. Version}\par
« Bien que ma maison soit très loin de l'école, je ne suis jamais en retard. Parce que je veux bien apprendre le chinois, je me lève tous les jours à six heures. Si je réussis l'examen cette année, j'irai étudier en Chine. »\par
\textbf{Points de vigilance :} \zh{从不} = « ne... jamais » (\zh{从} marque l'habitude constante) ; \zh{六点起床} = « se lever à six heures » ; \zh{考试成功} se rend naturellement par « réussir l'examen ».
\end{corrige}

\begin{corrige}[Exercice 11 — Expression écrite : Pourquoi j'apprends le chinois]
\textbf{Barème indicatif (sur 10) :} respect des consignes et des 3 structures complexes 3 pts ; correction grammaticale (ordre des mots, place de \zh{就}) 3 pts ; exactitude des caractères 2 pts ; richesse du vocabulaire et cohérence 2 pts.

\textbf{Proposition de rédaction modèle :}\par
\zh{我是雅温得的高中生。因为我很喜欢中国文化，所以我学习中文已经一年了。虽然汉字很难，但是我觉得很有意思。为了通过HSK考试，我每天练习说和写。如果我的成绩好，我就去中国留学。将来我想在喀麦隆的中国公司工作。}\par
\emph{Wǒ shì Yǎwēndé de gāozhōngshēng. Yīnwèi wǒ hěn xǐhuan Zhōngguó wénhuà, suǒyǐ wǒ xuéxí Zhōngwén yǐjīng yì nián le. Suīrán Hànzì hěn nán, dànshì wǒ juéde hěn yǒu yìsi. Wèile tōngguò HSK kǎoshì, wǒ měitiān liànxí shuō hé xiě. Rúguǒ wǒ de chéngjì hǎo, wǒ jiù qù Zhōngguó liúxué. Jiānglái wǒ xiǎng zài Kāmàilóng de Zhōngguó gōngsī gōngzuò.}\par
« Je suis lycéen à Yaoundé. Parce que j'aime beaucoup la culture chinoise, j'apprends le chinois depuis un an. Bien que les caractères soient difficiles, je les trouve très intéressants. Pour réussir l'examen HSK, je m'entraîne chaque jour à parler et à écrire. Si mes résultats sont bons, j'irai étudier en Chine. À l'avenir, je voudrais travailler dans une entreprise chinoise au Cameroun. »

\textbf{Vérification des consignes dans ce modèle :}
\begin{itemize}
\item 4 structures complexes utilisées : \zh{因为……所以……}, \zh{虽然……但是……}, \zh{为了……}, \zh{如果……就……} (minimum de 3 exigé).
\item 7 phrases, environ 95 caractères : conforme.
\item Dernière phrase = projet d'avenir avec \zh{将来}.
\end{itemize}

\textbf{Points de vigilance pour les élèves :}
\begin{itemize}
\item Ne jamais écrire \zh{因为……但是……} ni \zh{为了……所以……} : vérifier les paires de corrélatifs.
\item \zh{就} toujours après le sujet, avant le verbe : \zh{我就去}, jamais \zh{就我去}.
\item Ne pas répéter inutilement le sujet dans la deuxième proposition quand il est identique.
\item Soigner les caractères des mots imposés : \zh{梦想} (rêve), \zh{留学} (étudier à l'étranger), \zh{努力} — attention à \zh{努} (clé de la force \zh{力} en bas).
\end{itemize}
\end{corrige}

\newpage

\coursec{Fiche bilan}

\begin{popfigure}
\begin{tikzpicture}[
  mindmap,
  grow cyclic,
  every node/.style={concept, align=center, font=\small},
  concept color=popblue,
  level 1 concept/.append style={sibling angle=60, level distance=3.4cm, font=\small},
  root concept/.append style={concept color=popblue, font=\bfseries},
]
\node[concept, align=center]{Phrases\\complexes\\复合句}
  child[concept color=poporange] { node[concept, align=center]{Cause\\因为\ldots 所以\ldots\\yīnwèi\ldots suǒyǐ\ldots} }
  child[concept color=popgreen] { node[concept, align=center]{Concession\\虽然\ldots 但是\ldots\\suīrán\ldots dànshì\ldots} }
  child[concept color=poppurple] { node[concept, align=center]{Condition\\如果\ldots 就\ldots\\rúguǒ\ldots jiù\ldots} }
  child[concept color=poppink] { node[concept, align=center]{Hypothèse\\要是\ldots 就\ldots\\yàoshi\ldots jiù\ldots} }
  child[concept color=popgold] { node[concept, align=center]{But\\为了\ldots\\wèile\ldots} }
  child[concept color=popblue!60] { node[concept, align=center]{Pièges\\ordre des mots\\pas de double\\traduction} };
\end{tikzpicture}
\legende{Carte mentale de la leçon : les phrases complexes en chinois}
\end{popfigure}

\begin{bilan}
\textbf{Lexique clé (connecteurs logiques)}\par
\begin{center}
\begin{tabularx}{\linewidth}{|X|X|X|X|}
\hline
\rowcolor{popblueL} Caractères & Pinyin & Nature & Traduction \\
\hline
因为 & yīnwèi & conjonction & parce que \\
\hline
所以 & suǒyǐ & conjonction & donc, c'est pourquoi \\
\hline
虽然 & suīrán & conjonction & bien que \\
\hline
但是 & dànshì & conjonction & mais \\
\hline
如果 & rúguǒ & conjonction & si (condition) \\
\hline
要是 & yàoshi & conjonction & si (hypothèse, familier) \\
\hline
就 & jiù & adverbe & alors, aussitôt \\
\hline
为了 & wèile & préposition & afin de, pour \\
\hline
\end{tabularx}
\end{center}

\textbf{Règles de grammaire à connaître par cœur}\par
\begin{center}
\begin{tabularx}{\linewidth}{|X|X|X|}
\hline
\rowcolor{popblueL} Structure & Exemple (caractères + pinyin) & Traduction \\
\hline
因为 + cause，所以 + résultat & \zh{因为下雨，所以我不去。} Yīnwèi xià yǔ, suǒyǐ wǒ bú qù. & Comme il pleut, je n'y vais pas. \\
\hline
虽然 + concession，但是 + opposition & \zh{虽然很难，但是我喜欢。} Suīrán hěn nán, dànshì wǒ xǐhuan. & Bien que ce soit difficile, j'aime ça. \\
\hline
如果 + condition，就 + résultat & \zh{如果你努力，就会成功。} Rúguǒ nǐ nǔlì, jiù huì chénggōng. & Si tu travailles dur, tu réussiras. \\
\hline
要是 + hypothèse，就 + résultat & \zh{要是明天下雨，我们就不去杜阿拉。} Yàoshi míngtiān xià yǔ, wǒmen jiù bú qù Dù'ālā. & S'il pleut demain, nous n'irons pas à Douala. \\
\hline
为了 + but，sujet + verbe & \zh{为了通过考试，他每天学习。} Wèile tōngguò kǎoshì, tā měitiān xuéxí. & Pour réussir l'examen, il étudie chaque jour. \\
\hline
\end{tabularx}
\end{center}

\textbf{Méthodes express}\par
\begin{itemize}
  \item \cle{Repérer} d'abord la relation logique : cause ? opposition ? condition ? but ?
  \item \cle{Choisir} la paire de connecteurs correspondante (ils vont souvent par deux).
  \item \cle{Placer} la subordonnée en premier, puis la principale après la virgule.
  \item \cle{Vérifier} qu'un seul connecteur par membre de phrase : jamais de double traduction du français.
  \item \cle{Traduire} en dernier, en adaptant l'ordre au français.
\end{itemize}
\end{bilan}

\begin{aretenir}[Checklist avant le Bac]
\begin{itemize}
  \item $\square$ Je connais par cœur les cinq structures : \zh{因为\ldots 所以\ldots}, \zh{虽然\ldots 但是\ldots}, \zh{如果\ldots 就\ldots}, \zh{要是\ldots 就\ldots}, \zh{为了\ldots}.
  \item $\square$ Je sais qu'en chinois on peut utiliser les deux connecteurs ensemble (\zh{因为} + \zh{所以}), contrairement au français.
  \item $\square$ Je place la subordonnée avant la principale, séparées par une virgule chinoise ，.
  \item $\square$ Je distingue \zh{如果} (condition neutre) et \zh{要是} (hypothèse, registre plus oral).
  \item $\square$ Je sais que \zh{就} se place juste avant le verbe de la principale, après le sujet.
  \item $\square$ Je n'oublie pas que \zh{为了} + nom ou verbe exprime le but et se met en tête de phrase.
  \item $\square$ J'évite de traduire mot à mot « parce que \ldots donc » par deux idées séparées : je garde la paire de connecteurs.
  \item $\square$ Je fais attention aux tons : yīnwèi, suǒyǐ, suīrán, dànshì, rúguǒ, yàoshi, wèile.
  \item $\square$ Je sais transformer une phrase simple en phrase complexe (ajout d'une cause, d'une condition ou d'un but).
  \item $\square$ Je relis ma production pour vérifier l'ordre : subordonnée + ，+ sujet + 就/但是 + verbe.
  \item $\square$ Je peux produire un court texte (5-6 phrases) utilisant au moins trois connecteurs différents.
  \item $\square$ Je vérifie la ponctuation chinoise pleine chasse ：，。？ dans mes phrases.
\end{itemize}
\end{aretenir}

\findecours

\end{document}
